\chapter{First-order deformations and the graded first-order rigidity of Virasoro minimal models}\label{ch:L07}
\begingroup\setlength{\emergencystretch}{3em}

\providecommand{\lsevenvac}{\mathbf{1}}
\providecommand{\lsevenV}{\mathcal{V}}
\providecommand{\lsevenR}{R}
\providecommand{\lsevenE}{\mathcal{E}}
\providecommand{\lsevenK}{\mathsf{K}}
\providecommand{\lsevenN}{\mathsf{N}}
\providecommand{\lsevenLt}{\widetilde{L}}
\providecommand{\lsevencpq}{c_{p,q}}

This chapter documents the release of the LeanVOA library named in the front matter. Every Lean quotation below is checked by \texttt{scripts/textbook/sync.py}
to occur verbatim (up to whitespace and comments) in the source files of that release. The chapter was
first written against an earlier state of the library; checks that were run at that earlier state and not re-run for this edition are marked as such. The chapter covers the files assigned to it in
\texttt{Notes/Textbook/assignment.txt}: the general first-order deformation theory in
\leanpath{VOA/\allowbreak Core/\allowbreak Deformation*.lean}, \leanpath{VOA/\allowbreak Core/\allowbreak DeformedForm.lean},
\leanpath{VOA/\allowbreak Core/\allowbreak RigidityBridge.lean} and a few auxiliary linear-algebra files,
the example file \leanpath{VOA/\allowbreak Examples/\allowbreak NonzeroDeformation.lean}, and the
roughly eighty files \leanpath{VOA/\allowbreak Examples/\allowbreak Virasoro*.lean} whose subject is the rigidity of
the simple Virasoro vertex operator algebras.

The headline of the chapter is a single theorem,
\lean{VOA.VirModule.minimal_graded_rigidity}: \emph{for every pair of coprime integers
$p,q\ge2$, every graded first-order deformation of the simple Virasoro vertex operator algebra
$L(\lsevencpq,0)$ of central charge $\lsevencpq = 1-6(p-q)^2/(pq)$ is trivial}, that is, it is
undone by an infinitesimal change of coordinates. The theorem is proved without any hypothesis
beyond the label conditions, and \texttt{\#print axioms} reports only Lean's three standard axioms;
in particular \lean{VOA.pbw} (the PBW theorem of Chapter~\ref{ch:L02}; formerly a project axiom, now a theorem) is not used.

Three qualifications travel with the headline everywhere in this chapter (they are the caveats of the
release audit \texttt{Notes/Textbook/referee/RELEASE\_AUDIT\_RIGIDITY.md}). First, the theorem is
\emph{graded first-order} rigidity: the parameter ring is $\C[t]/(t^2)$, the vertex operation $Y$ is
deformed, the vacuum is held fixed, the deformation must respect the weight grading, and ``trivial''
means ``the coboundary of a grade-preserving, vacuum-fixing $\varphi$''. Nothing is claimed about formal
($\C[[t]]$), ungraded or grading-moving deformations; the library's \emph{ungraded} notion
\lean{IsRigid} is false for every universal Virasoro VOA $V(c,0)$, for a definitional reason
(Warning~\ref{l07:warn:isrigid}), and is never quoted as true in this book. Second, the notion
\lean{IsRigidGraded} is not vacuous: the library constructs the $c$-derivative $\partial_cY$ on the universal
Virasoro VOA $V(c,0)$ and proves \lean{VOA.}\allowbreak\lean{VirPoly.}\allowbreak\lean{not_}\allowbreak\lean{isRigidGraded_}\allowbreak\lean{virasoroVOA}, i.e.\ $V(c,0)$ is not graded first-order rigid,
for every $c\in\C$ (Remark~\ref{l07:rem:nonvacuity}). The witness deformation moves the central charge, which the
library's notion of first-order deformation allows (it is not required to preserve $\omega$ or $c$). This caveat of
the release audit has therefore been discharged since the audit was written. Third, at $(p,q) = (2,3)$ the central charge is $c_{2,3} = 0$
and $L(0,0) = \C\lsevenvac$, so that instance is trivially true; the content lies in the other labels.

The chapter has four jobs.
\begin{enumerate}
\item Explain from first principles what a first-order deformation of a vertex operator algebra
  is, what it means for one to be \emph{trivial}, and what \emph{rigidity} means, both
  mathematically (vanishing of a second cohomology group, in the spirit of Gerstenhaber's
  deformation theory of associative algebras) and physically (absence of exactly marginal
  deformations of a chiral algebra: the theory is an isolated point of its moduli space).
\item Present the library's definitions --- \lean{FirstOrderDeformation}, \lean{IsGraded},
  \lean{IsCoboundary}, \lean{IsTrivial}, \lean{IsRigid}, \lean{IsRigidGraded} --- with their
  verbatim Lean source, and run the fidelity check on them with particular care. This is the heart
  of the chapter: the referee must be convinced that \lean{IsRigidGraded} is the rigidity notion
  the authors intended, must know exactly which deformations it quantifies over, and must know
  how it could have been vacuous and why it is not. We shall see that one of the two notions the
  library defines, \lean{IsRigid}, is \emph{false for every Virasoro vertex operator algebra for a
  definitional reason}, and that the library knows this and uses only \lean{IsRigidGraded}.
\item Explain the proof of the headline theorem on paper. The proof has three steps that can be
  stated in a paragraph each: (a) a graded deformation of $L(c,0)$ has exactly one gauge-invariant
  parameter $\delta$, the first-order change of the central charge; (b) if $\delta = 0$ the
  deformation is trivial (a reconstruction theorem using simplicity); (c) if $c = \lsevencpq$, a
  null vector of the universal vacuum module, whose norm has a \emph{simple} zero at $c$, forces
  $\delta = 0$. Step (c) needs the Kac determinant of the vacuum module, which is the subject of the
  companion Chapter~\ref{ch:L06}; we state precisely which results of that chapter are used.
\item Present the worked cases that the library proved on the way (central charge $0$,
  Lee--Yang $-22/5$, Ising $1/2$, $M(2,7)$ at $-68/7$, tricritical Ising $7/10$, the families
  $M(2,q)$ and $M(p,3)$), the dead ends that were recorded (a tautological ``structural''
  rigidity theorem, an unsatisfiable ``bridge'' hypothesis, the false belief that an isolated
  degenerate charge suffices), and close with the referee's checklist.
\end{enumerate}

The foundations --- formal distributions, modes $u_{(n)}$, the definitions of vertex algebra and
vertex operator algebra, the Borcherds identity, ideals, simple quotients and invariant bilinear
forms --- are those of Chapters~\ref{ch:L01}, \ref{ch:L03} and~\ref{ch:L04}, to which we point for every such notion. The Virasoro
vertex operator algebra itself (its construction, its Verma modules, the Kac determinant) is the
subject of Chapter~\ref{ch:L06}; since that chapter is a companion rather than a prerequisite we
recall in Section~\ref{l07:sec:virasoro} every Virasoro definition our statements use, with its
verbatim Lean source.

%======================================================================
\section{Why deformations, and why minimal models should be rigid}\label{l07:sec:motivation}
%======================================================================

\subsection{Deformations in algebra: Gerstenhaber's picture}

Let $A$ be an associative algebra over $\C$ with product $\mu(a,b) = ab$. A \emph{first-order
deformation} of $A$ is an associative product on $A[t]/(t^2) = A\oplus tA$ of the form
\begin{equation*}
 \mu_t(a,b) = \mu(a,b) + t\,\mu_1(a,b),
\end{equation*}
extended $\C[t]/(t^2)$-bilinearly, where $\mu_1 : A\times A\to A$ is bilinear. Associativity of
$\mu_t$ modulo $t^2$ is the condition
\begin{equation*}
 \mu_1(ab,c) + \mu_1(a,b)c = \mu_1(a,bc) + a\mu_1(b,c),
\end{equation*}
which says that $\mu_1$ is a Hochschild $2$-cocycle. Two deformations are \emph{equivalent} if
they differ by an infinitesimal change of coordinates $a\mapsto a + t\varphi(a)$ with
$\varphi : A\to A$ linear; such a change transforms $\mu_1$ into
$\mu_1 + \delta\varphi$ with $(\delta\varphi)(a,b) = \varphi(a)b - \varphi(ab) + a\varphi(b)$,
a Hochschild coboundary. A deformation is \emph{trivial} if it is equivalent to the zero
deformation, and $A$ is \emph{rigid} (to first order) if every first-order deformation is
trivial; equivalently the second Hochschild cohomology $HH^2(A,A)$ vanishes. This is the theory
created by Gerstenhaber in 1964. Two features of it will recur:
\begin{itemize}
\item \emph{The cocycle condition is linear in $\mu_1$.} The original axiom (associativity) is
  quadratic in the product; substituting $\mu + t\mu_1$ and discarding $t^2$ leaves the
  ``derivative'' of the axiom in the direction $\mu_1$, which is linear. Hence the
  deformations form a vector space, and the trivial ones a subspace.
\item \emph{Rigidity is a statement about a quotient.} A nonzero $\mu_1$ can be trivial; what is
  asked is that every cocycle is a coboundary. A rigidity theorem that forgot the quotient (``every
  $\mu_1$ is zero'') would be false for every algebra, since $\delta\varphi$ is almost never zero.
\end{itemize}

\subsection{Deformations of vertex algebras}

A vertex algebra (Definition~\ref{l03:def:VA} of Chapter~\ref{ch:L03}) has infinitely many
products $u_{(n)}v$, $n\in\Z$, packaged in the state--field correspondence
$Y(u,z) = \sum_n u_{(n)}z^{-n-1}$, and its main identity, the Borcherds identity
(Theorem~\ref{l03:thm:borcherds}), is quadratic in $Y$. Deforming $Y$ to $Y + tY_1$ over the dual
numbers and discarding $t^2$ therefore gives, exactly as for associative algebras, a linear
``cocycle'' condition on $Y_1$, and infinitesimal changes of coordinates $1 + t\varphi$ give
``coboundaries''. Cohomology theories for vertex algebras in which the second cohomology
classifies such first-order deformations were developed by Huang (for grading-restricted vertex
algebras) and by others; see the Notes at the end of the chapter. The library constructs no
cohomology theory at all: it works with the cocycle and coboundary conditions directly, as
elementary statements about modes. This is a deliberate design decision, recorded in the
docstring of \leanpath{VOA/\allowbreak Core/\allowbreak Deformation.lean}, and it makes the fidelity check easier,
because nothing is hidden inside a homological construction.

\subsection{The physical picture: moduli of conformal field theories}

In two-dimensional conformal field theory one asks whether a given theory sits in a continuous
family. The standard example is the compactified free boson: the radius $R$ of the target circle
can be varied continuously, and every value gives a consistent CFT; the operator that moves $R$ is
the ``exactly marginal'' operator $\partial X\bar\partial X$ of conformal weight $(1,1)$. A theory
with no exactly marginal deformations is \emph{isolated}: it is a point, not a curve, in the space
of theories.

A vertex operator algebra is the \emph{chiral algebra} of a CFT: the algebra of fields depending
holomorphically on $z$, with their operator product expansions. Deforming the vertex algebra
means deforming the OPE coefficients of the chiral fields while keeping the space of states fixed.
For the free boson the chiral algebra is the Heisenberg vertex operator algebra, generated by a
current $J(z)$ with $J(z)J(w)\sim k/(z-w)^2$; changing $k$ looks like a deformation, but it is
undone by rescaling $J\mapsto\lambda J$, i.e.\ by a change of coordinates, so it is trivial. For
the Virasoro algebra the analogous parameter is the central charge in
\begin{equation*}
 T(z)T(w)\sim\frac{c/2}{(z-w)^4} + \frac{2T(w)}{(z-w)^2} + \frac{\partial T(w)}{z-w},
\end{equation*}
and here rescaling $T$ does \emph{not} undo a change of $c$: rescaling $T\mapsto\lambda T$ changes
the coefficient $2$ of the second-order pole, which is fixed by the requirement that $T$ generate
the conformal symmetry. So the universal Virasoro vertex operator algebra $V(c,0)$ has a genuine
one-parameter family of deformations, $c\mapsto c+t$. (We shall see that the library proves the
``coordinate change cannot move $c$'' half of this statement as a theorem,
Theorem~\ref{l07:thm:gaugeline}, and also constructs the deformation $\partial_c$ itself,
\lean{VOA.}\allowbreak\lean{VirPoly.}\allowbreak\lean{cDeformation}; Section~\ref{l07:sec:nonvacuity}.)

\subsection{Why the minimal models should be rigid}

The \emph{minimal models} of Belavin, Polyakov and Zamolodchikov are the CFTs whose chiral algebra
is the simple quotient $L(c,0)$ of $V(c,0)$ at one of the special central charges
\begin{equation*}
 \lsevencpq = 1-\frac{6(p-q)^2}{pq},\qquad p,q\ge2,\ \gcd(p,q) = 1 .
\end{equation*}
(Examples: $c_{2,3} = 0$, the Lee--Yang model $c_{2,5} = -22/5$, the Ising model $c_{3,4} = 1/2$, the
tricritical Ising model $c_{4,5} = 7/10$.) At these values, and only there, the vacuum module
$V(c,0)$ contains a \emph{null vector} --- a nonzero state $s$ of positive weight orthogonal to
everything for the invariant form --- and $L(c,0) = V(c,0)/\langle s\rangle$. The lowest one sits at
weight $(p-1)(q-1)$. The physical heuristic for rigidity is then:
\begin{enumerate}
\item The chiral algebra $L(c,0)$ is generated by $T$, and the only parameter in the $TT$ OPE is
  $c$. So any deformation of $L(c,0)$, after a change of coordinates, can only move $c$, to
  $c + t\delta$.
\item In $L(c,0)$ the null vector is \emph{zero}: $s = 0$ is a relation of the algebra. A
  deformation must keep this relation (to first order, after correcting $s$ by $t s'$).
\item But the norm $\langle s,s\rangle$ of the corresponding universal state, as a function of the
  central charge, vanishes at $c = \lsevencpq$ to \emph{first} order. Moving $c$ to $c + t\delta$
  changes the norm by $t\delta\cdot\frac{d}{dc}\langle s,s\rangle\ne0$ unless $\delta = 0$, and a
  state with nonzero norm cannot be zero. Hence $\delta = 0$, and by step 1 the deformation is
  trivial.
\end{enumerate}
That is the whole idea, and the Lean proof follows it faithfully. The work is in making each step
precise: in step 1, ``after a change of coordinates'' and ``generated by $T$'' become a
reconstruction theorem (Section~\ref{l07:sec:reconstruction}); in step 2, ``keep the relation''
becomes a pairing identity over the dual numbers (Section~\ref{l07:sec:obstruction}); in step 3 one
needs to know that a suitable norm, or the vacuum Gram determinant, has a \emph{simple} zero at
$\lsevencpq$ (Section~\ref{l07:sec:all}). The word ``simple'' is essential: the minimal-model
charges accumulate at $c = 1$, so ``$\lsevencpq$ is an isolated degenerate value'' is not even true
across all levels, and at a fixed level an isolated root of multiplicity two would \emph{not}
obstruct a first-order motion (Warning~\ref{l07:warn:doubleroot}).

The physical statement is about first order only, and so is the theorem. It says nothing about
analytic families, about the full (non-chiral) CFT, or about the modules of $L(c,0)$; we list
precisely what it does and does not say in Section~\ref{l07:sec:fidelity-rigid}.

%======================================================================
\section{First-order deformations of a vertex operator algebra}\label{l07:sec:deformations}
%======================================================================

Throughout, $W$ is a vertex operator algebra in the sense of Definition~\ref{l03:def:VOA} of
Chapter~\ref{ch:L03}, with state space $\lsevenV = \bigoplus_{n\ge0}\lsevenV_n$ (CFT type:
$\lsevenV_0 = \C\lsevenvac$, $\lsevenV_n = 0$ for $n<0$), state--field map $Y$, vacuum $\lsevenvac$,
conformal vector $\omega$ and central charge $c$. We write $u_{(n)}v$ for the $n$-th product
(\lean{W.smul u n v} in Lean) and $L_n = \omega_{(n+1)}$.

\subsection{Derivation of the axioms}\label{l07:sec:derivation}

Let $\lsevenR = \C[t]/(t^2)$, the ring of dual numbers, and let $\lsevenV_\lsevenR = \lsevenR\otimes_\C
\lsevenV = \lsevenV\oplus t\lsevenV$. We look for a vertex algebra structure on $\lsevenV_\lsevenR$ over
$\lsevenR$, with the same vacuum, of the form
\begin{equation}\label{l07:eq:Yt}
 Y^t(u,z)v = Y(u,z)v + t\,Y_1(u,z)v,\qquad\text{i.e.}\qquad u^t_{(n)}v = u_{(n)}v + t\,u^1_{(n)}v,
\end{equation}
where $Y_1 : \lsevenV\to\End(\lsevenV)[[z^{\pm1}]]$ is $\C$-linear, $u^1_{(n)}v$ denotes the coefficient
of $z^{-n-1}$ in $Y_1(u,z)v$, and everything is extended $\lsevenR$-bilinearly. Each vertex-algebra
axiom for $(\lsevenV_\lsevenR, Y^t,\lsevenvac)$ splits into a $t^0$ part, which is the axiom for $W$ and
holds automatically, and a $t^1$ part, which is a condition on $Y_1$:
\begin{enumerate}[label=(D\arabic*)]
\item \emph{Truncation} (VA5). $u^t_{(n)}v = 0$ for $n\gg0$; at order $t$: for all $u,v$ there is
  $N$ with $u^1_{(n)}v = 0$ for $n\ge N$.
\item \emph{Vacuum field} (VA2). $Y^t(\lsevenvac,z) = \mathrm{id}$; the right side has no $t$, so
  $Y_1(\lsevenvac,z) = 0$.
\item \emph{Creation} (VA3). $u^t_{(n)}\lsevenvac = 0$ for $n\ge0$ and $u^t_{(-1)}\lsevenvac = u$; the
  vector $u$ on the right has no $t$-component (we deform $Y$, not the space), so
  $u^1_{(n)}\lsevenvac = 0$ for all $n\ge-1$.
\item \emph{Borcherds identity.} Every term of \eqref{l03:eq:borcherds} contains exactly two
  modes. Substituting $u_{(k)}\mapsto u_{(k)} + tu^1_{(k)}$ and keeping the coefficient of $t$
  gives, for all $u,v,x\in\lsevenV$ and $\ell,m,n\in\Z$,
  \begin{equation}\label{l07:eq:jac1}
  \begin{aligned}
   &\sum_{i\ge0}(-1)^i\binom{\ell}{i}\Big(u^1_{(\ell+m-i)}v_{(n+i)}x + u_{(\ell+m-i)}v^1_{(n+i)}x
     - (-1)^\ell\big(v^1_{(\ell+n-i)}u_{(m+i)}x + v_{(\ell+n-i)}u^1_{(m+i)}x\big)\Big)\\
   &\qquad = \sum_{i\ge0}\binom mi\Big((u^1_{(\ell+i)}v)_{(m+n-i)}x + (u_{(\ell+i)}v)^1_{(m+n-i)}x\Big).
  \end{aligned}
  \end{equation}
  On each side one sums over the two ways of placing a single $Y_1$. This is the ``$Y_1$ is a
  $2$-cocycle'' condition, and it is linear in $Y_1$.
\end{enumerate}
Locality (VA4) is not imposed separately: over a commutative base ring the Borcherds identity
together with vacuum, creation and truncation is equivalent to the full set of vertex-algebra
axioms (this is the Jacobi-identity definition of Frenkel--Lepowsky--Meurman and Lepowsky--Li, and
its equivalence with the locality definition; for the library's own axioms over $\C$ see
Chapter~\ref{ch:L03}). The grading axiom (VA1) is \emph{not} imposed at this stage; its first-order
part will be the separate predicate \lean{IsGraded} below.

\subsection{The definition}

\begin{definition}[First-order deformation]\label{l07:def:fod}
A \emph{first-order deformation} of $W$ is a $\C$-linear map
$Y_1 : \lsevenV\to\End(\lsevenV)[[z^{\pm1}]]$ satisfying (D1)--(D4). We write
$u^1_{(n)}v := (Y_1(u))_{(n)}v$ for its products.
\end{definition}

\begin{leancode}
structure FirstOrderDeformation (W : VertexOperatorAlgebra 𝒱) where
  Y1 : 𝒱 →ₗ[ℂ] (End ℂ 𝒱)⟦z^±⟧
  truncation : ∀ u v : 𝒱, ∃ N : ℤ, ∀ n : ℤ, N ≤ n → mode (Y1 u) n v = 0
  vacuum_Y1 : Y1 W.vacuum = 0
  creation : ∀ (u : 𝒱) (n : ℤ), -1 ≤ n → mode (Y1 u) n W.vacuum = 0
  jacobi : ∀ (u v x : 𝒱) (ℓ m n : ℤ),
      (∑ᶠ i : ℤ, (negOnePow i * cbz ℓ i) •
          (mode (Y1 u) (ℓ + m - i) (W.smul v (n + i) x)
            + W.smul u (ℓ + m - i) (mode (Y1 v) (n + i) x)
            - negOnePow ℓ •
              (mode (Y1 v) (ℓ + n - i) (W.smul u (m + i) x)
                + W.smul v (ℓ + n - i) (mode (Y1 u) (m + i) x))))
        = ∑ᶠ i : ℤ, cbz m i •
            (W.smul (mode (Y1 u) (ℓ + i) v) (m + n - i) x
              + mode (Y1 (W.smul u (ℓ + i) v)) (m + n - i) x)

def smul₁ (D : FirstOrderDeformation W) (u : 𝒱) (n : ℤ) (v : 𝒱) : 𝒱 := mode (D.Y1 u) n v
\end{leancode}
(\lean{VOA.Deform.FirstOrderDeformation}, \lean{VOA.Deform.FirstOrderDeformation.smul₁},
\leanpath{VOA/\allowbreak Core/\allowbreak Deformation.lean}; the ambient context is
\lean{variable {𝒱 : Type u} [AddCommGroup 𝒱] [Module ℂ 𝒱]}, i.e.\ $\lsevenV$ is a complex vector space;
field docstrings removed.)

\begin{dictionary}
\textbf{Decoding.}
\begin{itemize}
\item \lean{Y1 : 𝒱 →ₗ[ℂ] (End ℂ 𝒱)⟦z^±⟧} is a $\C$-linear map from $\lsevenV$ to formal distributions
  with coefficients in $\End_\C(\lsevenV)$ --- \emph{exactly the type} of the undeformed
  \lean{W.Y} (Chapter~\ref{ch:L03}, Definition~\ref{l03:def:VA}). Linearity in $u$ is built into
  \lean{→ₗ[ℂ]}; linearity of each $u^1_{(n)}$ in $v$ is built into \lean{End ℂ 𝒱}. So $Y_1$ is a
  bilinear family of products, with nothing assumed beyond bilinearity.
\item \lean{mode (Y1 u) n v} is $u^1_{(n)}v$, by the convention \lean{mode f n = f (-n-1)} of
  Definition~\ref{l01:def:mode}: the coefficient of $z^{-n-1}$. \lean{smul₁ D u n v} is the same
  thing with a name, mirroring \lean{W.smul u n v} $= u_{(n)}v$.
\item \lean{truncation} is (D1), with $N$ depending on $u,v$. \lean{vacuum_Y1} is (D2):
  \lean{Y1 W.vacuum = 0} means the whole distribution $Y_1(\lsevenvac,z)$ is zero.
  \lean{creation} is (D3): $u^1_{(n)}\lsevenvac = 0$ for every $n\ge-1$ --- note \emph{including}
  $n = -1$, which is the first-order part of $u_{(-1)}\lsevenvac = u$.
\item \lean{jacobi} is \eqref{l07:eq:jac1}. \lean{negOnePow i} is $(-1)^i$, \lean{cbz ℓ i} is the
  binomial coefficient $\binom{\ell}{i}$ with integer lower index, \emph{zero for $i<0$}
  (Definition~\ref{l01:def:gbinom}), so the sums over all $i\in\Z$ are the sums over $i\ge0$ of the
  displayed formula. The left side has four terms: $u^1v$, $uv^1$, $v^1u$, $vu^1$; the right side
  has two: $(u^1_{(\ell+i)}v)_{(m+n-i)}x$ and $(u_{(\ell+i)}v)^1_{(m+n-i)}x$. This is
  term-by-term the Leibniz derivative of the library's undeformed Borcherds identity
  \lean{VertexOperatorAlgebra.jacobi} (Theorem~\ref{l03:thm:borcherds}), whose summands are
  $u_{(\ell+m-i)}v_{(n+i)}x - (-1)^\ell v_{(\ell+n-i)}u_{(m+i)}x$ and $\binom mi(u_{(\ell+i)}v)_{(m+n-i)}x$.
\item \lean{∑ᶠ} is Mathlib's \lean{finsum}: the sum of a family with finite support, and \emph{by
  definition $0$ if the support is infinite}.
\end{itemize}

\textbf{Could \lean{jacobi} be vacuous?} Yes, in principle, through \lean{finsum}: if one of the
two families had infinite support, its \lean{finsum} would be $0$ by fiat and the axiom would
degenerate. The library proves that this does not happen: \lean{jacobiLHS_support_finite} and
\lean{jacobiRHS_support_finite} (same file) show both supports are finite, using
\lean{truncation} and the undeformed VA5. (Negative $i$ are killed by \lean{cbz}; large $i$ by
truncation.) So \lean{jacobi} is a genuine finite identity, the one displayed.

\textbf{What is fixed and what is deformed.}
\begin{itemize}
\item The \emph{state space} is fixed: $\lsevenV_\lsevenR = \lsevenV\oplus t\lsevenV$, a free $\lsevenR$-module.
  Deformations that change the underlying graded dimensions (as happens at a jump of the family
  $c\mapsto L(c,0)$) are \emph{not} first-order deformations in this sense; this is the standard
  convention of algebraic deformation theory.
\item The \emph{vacuum} is fixed. The source records why this loses nothing: a deformed vacuum
  $\lsevenvac + t\lsevenvac'$ can be moved back to $\lsevenvac$ by the coordinate change of
  $\varphi = -(\,v\mapsto\lsevenvac'_{(-1)}v\,)$. We have not checked this remark in Lean; it is not
  used.
\item The \emph{conformal vector} is \textbf{not} part of the data. Nothing requires $Y^t$ to have
  a conformal vector, let alone $\omega$. So these are deformations of the \emph{vertex algebra}
  underlying $W$, not of the vertex operator algebra. For the rigidity question this is the
  \emph{harder} (larger) class: every deformation of $W$ as a VOA (moving $Y$ and possibly $\omega$)
  gives a deformation of the vertex algebra by forgetting $\omega$.
\item The \emph{translation operator} is not mentioned. In a vertex algebra it is determined by
  $Tu = u_{(-2)}\lsevenvac$, so the deformed one is $T + tT_1$ with $T_1u = u^1_{(-2)}\lsevenvac$;
  (D3) does not force $T_1 = 0$ because it constrains only $n\ge-1$. Similarly skew-symmetry is not
  an axiom; it follows from (D4) as it does in the undeformed case.
\item The \emph{base ring} is the dual numbers only. This is ``first order''. There is no
  statement over $\C[t]/(t^3)$ or $\C[[t]]$.
\end{itemize}

\textbf{Same as the literature?} With these conventions, a \lean{FirstOrderDeformation} is exactly
a vertex algebra structure on $\lsevenR\otimes\lsevenV$ over $\lsevenR$ reducing to $W$ modulo $t$ and
having vacuum $\lsevenvac$: (D1)--(D4) are necessary (they are the order-$t$ parts of axioms that any
such structure satisfies) and, by the Jacobi-identity form of the axioms, sufficient. The
direction that matters for rigidity is the first: \emph{no genuine first-order deformation is
excluded by the definition}, so a rigidity theorem quantifying over \lean{FirstOrderDeformation}
is not weakened by a too-narrow notion of deformation.
\end{dictionary}

\begin{status}
A definition. The support-finiteness lemmas \lean{jacobiLHS_}\allowbreak\lean{support_}\allowbreak\lean{finite},
\lean{jacobiRHS_}\allowbreak\lean{support_}\allowbreak\lean{finite} and the linearity lemmas for \lean{smul₁}
(\lean{smul₁_add_left}, \lean{smul₁_smul_right}, \ldots) are proved in the same file, axiom-clean.
\end{status}

\begin{physics}
A first-order deformation is the chiral-algebra version of perturbing a CFT by an infinitesimal
amount: one keeps the space of states and the vacuum and changes every operator product expansion
coefficient $C_{uv}^{w}$ by $t\,\delta C_{uv}^{w}$, discarding $t^2$. Conditions (D1)--(D4) say that the perturbed
OPEs are still consistent: the OPE of two fields is still a Laurent series (truncation), the identity
field is still the identity (vacuum), the state--operator map still sends $u$ to the field creating $u$
from the vacuum (creation), and the perturbed OPE is still associative and local (Borcherds). In
conformal perturbation theory such a change is produced by integrating a \emph{marginal} operator over the
plane; for a purely chiral algebra the analogous objects are changes of the OPE structure constants, and
the question whether they are ``exactly marginal'' is the question whether they integrate to a family.
The definition does not require the perturbed theory to keep a stress tensor; that only enlarges the class
of perturbations considered.
\end{physics}

\subsection{Graded deformations}

\begin{definition}[Graded deformation]\label{l07:def:graded}
A first-order deformation is \emph{graded} if $u^1_{(n)}v\in\lsevenV_{k+\ell-n-1}$ whenever
$u\in\lsevenV_k$, $v\in\lsevenV_\ell$, $n\in\Z$; that is, each $u^1_{(n)}$ shifts the weight exactly as
$u_{(n)}$ does (VA1 at order $t$).
\end{definition}
\begin{leancode}
def IsGraded (D : FirstOrderDeformation W) : Prop :=
  ∀ (k ℓ : ℤ) (u : 𝒱), u ∈ W.grading k → ∀ (v : 𝒱), v ∈ W.grading ℓ → ∀ (n : ℤ),
    D.smul₁ u n v ∈ W.grading (k + ℓ - n - 1)
\end{leancode}
(\lean{VOA.Deform.FirstOrderDeformation.IsGraded}, \leanpath{VOA/\allowbreak Core/\allowbreak Deformation.lean}.)

\begin{dictionary}
\lean{W.grading k} is the subspace $\lsevenV_k$ (the field \lean{grading} of
Definition~\ref{l03:def:VA}). The quantifiers are over all integers $k,\ell,n$ and all
homogeneous $u,v$; for $k<0$ or $\ell<0$ the hypothesis forces $u=0$ or $v=0$ (CFT type) and the
condition is empty. The predicate is exactly VA1 for the $t$-component.

\textbf{Why a predicate and not a field.} The source explains: an ungraded first-order deformation
is a legitimate object, and making gradedness a field would silently strengthen
\lean{FirstOrderDeformation}. We shall see in Warning~\ref{l07:warn:isrigid} that this choice has
a consequence for the \emph{ungraded} rigidity notion \lean{IsRigid}, which the library
acknowledges and avoids.

\textbf{What gradedness buys.} For graded $Y_1$ the mode index is determined by three weights:
$u^1_{(n)}$ maps $\lsevenV_a\otimes\lsevenV_b$ to $\lsevenV_{a+b-n-1}$, so the component of $Y_1$ from
$\lsevenV_a\otimes\lsevenV_b$ to $\lsevenV_d$ lives in the single mode $n = a+b-d-1$, and is a bilinear
map between finite-dimensional spaces. Since $\lsevenV_d = 0$ for $d<0$, truncation (D1) is then
automatic (\lean{Deform.truncated_of_graded} in \leanpath{VOA/\allowbreak Core/\allowbreak DeformationGraded.lean}).
\end{dictionary}

\begin{physics}
Gradedness says the perturbation preserves conformal weights: the field of a state of weight $k$
acting on a state of weight $\ell$ still produces, at the pole $(z-w)^{-n-1}$, states of weight $k+\ell-n-1$.
Physically this is the requirement that the perturbation commutes with dilatations $z\mapsto\lambda z$, i.e.\ is
\emph{marginal} rather than relevant or irrelevant: a perturbation by an operator of the wrong weight would
introduce a scale and break conformal invariance. (Heuristic: a CFT perturbed by a marginal operator stays
conformal to first order; that is what gradedness models here.)
\end{physics}

\subsection{Trivial deformations}

Let $\varphi:\lsevenV\to\lsevenV$ be linear and consider the $\lsevenR$-linear automorphism $g = 1+t\varphi$
of $\lsevenV_\lsevenR$, with inverse $1-t\varphi$. Transporting $Y$ along $g$ gives the vertex algebra
structure $Y^g(u,z) = g^{-1}Y(gu,z)g$, and modulo $t^2$
\begin{equation*}
 Y^g(u,z)v = Y(u,z)v + t\big(Y(\varphi u,z)v - \varphi(Y(u,z)v) + Y(u,z)\varphi v\big).
\end{equation*}
Its vacuum is $g^{-1}\lsevenvac = \lsevenvac - t\varphi(\lsevenvac)$, which is $\lsevenvac$ precisely when
$\varphi(\lsevenvac) = 0$. So the first-order deformations ``obtained by a change of coordinates'' are
\begin{equation}\label{l07:eq:delta}
 (\delta\varphi)(u,n,v) := (\varphi u)_{(n)}v - \varphi(u_{(n)}v) + u_{(n)}(\varphi v),
 \qquad \varphi(\lsevenvac) = 0 .
\end{equation}
In words: $\delta\varphi$ measures the failure of $\varphi$ to be a derivation of every $n$-th
product.

\begin{definition}[Coboundary, trivial deformation]\label{l07:def:trivial}
A linear map $\varphi:\lsevenV\to\lsevenV$ is \emph{grade-preserving} if $\varphi(\lsevenV_k)\subseteq\lsevenV_k$
for all $k$. A cochain $f:\lsevenV\to\End(\lsevenV)[[z^{\pm1}]]$ is a \emph{coboundary} if there is a
grade-preserving $\varphi$ with $\varphi(\lsevenvac) = 0$ and $f(u)_{(n)}v = (\delta\varphi)(u,n,v)$ for
all $u,n,v$. A first-order deformation is \emph{trivial} if its $Y_1$ is a coboundary.
\end{definition}
\begin{leancode}
def GradePreserving (W : VertexOperatorAlgebra 𝒱) (φ : 𝒱 →ₗ[ℂ] 𝒱) : Prop :=
  ∀ (k : ℤ) (u : 𝒱), u ∈ W.grading k → φ u ∈ W.grading k

def IsCoboundary (W : VertexOperatorAlgebra 𝒱) (f : 𝒱 →ₗ[ℂ] (End ℂ 𝒱)⟦z^±⟧) : Prop :=
  ∃ φ : 𝒱 →ₗ[ℂ] 𝒱, φ W.vacuum = 0 ∧ GradePreserving W φ ∧
    ∀ (u : 𝒱) (n : ℤ) (v : 𝒱),
      mode (f u) n v = W.smul (φ u) n v - φ (W.smul u n v) + W.smul u n (φ v)

def IsTrivial (D : FirstOrderDeformation W) : Prop := IsCoboundary W D.Y1
\end{leancode}
(\lean{VOA.Deform.GradePreserving}, \lean{VOA.Deform.IsCoboundary},
\lean{VOA.Deform.FirstOrderDeformation.IsTrivial}, \leanpath{VOA/\allowbreak Core/\allowbreak Deformation.lean}.)

\begin{dictionary}
\lean{∃ φ : 𝒱 →ₗ[ℂ] 𝒱} is ``there is a linear map $\varphi$''; the three conjuncts are
$\varphi(\lsevenvac) = 0$, grade preservation, and the formula \eqref{l07:eq:delta} for every
$u\in\lsevenV$ (not only homogeneous), $n\in\Z$, $v\in\lsevenV$. \lean{W.smul (φ u) n v} is
$(\varphi u)_{(n)}v$, \lean{φ (W.smul u n v)} is $\varphi(u_{(n)}v)$, \lean{W.smul u n (φ v)} is
$u_{(n)}(\varphi v)$.

\textbf{Sign.} The opposite overall sign ($-\delta\varphi$) gives the same set, since $\varphi$ may be
replaced by $-\varphi$. Indeed the library's reconstruction theorem produces its trivialising map as
\lean{-ψ} (Theorem~\ref{l07:thm:closedlift}).

\textbf{The gauge group.} Allowed coordinate changes are $1+t\varphi$ with $\varphi$ grade-preserving
and vacuum-fixing. $\varphi$ is \emph{not} required to fix $\omega$, nor to commute with anything.
So triviality is generous in the expected way: it allows every infinitesimal graded,
vacuum-preserving change of basis. Restricting $\varphi$ to be grade-preserving is natural for
graded deformations; it matters only for ungraded ones (Warning~\ref{l07:warn:isrigid}).
\end{dictionary}

The first thing to check about any notion of ``trivial deformation'' is that the trivial ones are
deformations at all: that $\delta\varphi$ satisfies (D1)--(D4).

\begin{physics}
A trivial deformation is a field redefinition: one relabels the states $u\mapsto u+t\varphi(u)$, as a physicist
relabels operators by mixing operators of the same dimension. Such a ``deformation'' changes the OPE
coefficients written in the old basis but not the theory; it is the analogue of a redundant operator in
conformal perturbation theory. Requiring $\varphi$ to preserve weights says operators only mix with operators
of the same scaling dimension; requiring $\varphi(\lsevenvac)=0$ says the identity operator is not redefined.
\end{physics}

\begin{theorem}[Coboundaries are cocycles]\label{l07:thm:cobcoc}
For every linear $\varphi$ with $\varphi(\lsevenvac) = 0$, the cochain $\delta\varphi$ satisfies
(D1)--(D4). In particular every coboundary is the $Y_1$ of a first-order deformation.
\end{theorem}
\begin{leancode}
def cocycles (W : VertexOperatorAlgebra 𝒱) : Submodule ℂ (Cochain 𝒱) where
  carrier := {f | Truncated f ∧ f W.vacuum = 0 ∧ CreationAx W f ∧ JacobiAx W f}

theorem coboundaries_le_cocycles (W : VertexOperatorAlgebra 𝒱) :
    coboundaries W ≤ cocycles W := by

theorem deltaCochain_mem_cocycles (W : VertexOperatorAlgebra 𝒱) {φ : 𝒱 →ₗ[ℂ] 𝒱}
    (hφ0 : φ W.vacuum = 0) : deltaCochain W φ ∈ Deform.cocycles W := by
\end{leancode}
(\lean{VOA.Deform.cocycles}, \lean{VOA.Deform.coboundaries_le_cocycles},
\leanpath{VOA/\allowbreak Core/\allowbreak DeformationGraded.lean}; \lean{VOA.NonzeroDeform.deltaCochain_mem_cocycles},
\leanpath{VOA/\allowbreak Examples/\allowbreak NonzeroDeformation.lean}. The proof
of the submodule axioms of \lean{cocycles} is omitted.)

\begin{dictionary}
\lean{Cochain 𝒱} abbreviates \lean{𝒱 →ₗ[ℂ] (End ℂ 𝒱)⟦z^±⟧}, the type of \lean{Y1}.
\lean{Truncated}, \lean{CreationAx}, \lean{JacobiAx} are (D1), (D3), (D4) restated for an arbitrary
cochain, verbatim copies of the fields of \lean{FirstOrderDeformation}; so \lean{cocycles W} is the
set of $Y_1$'s of first-order deformations, and the library proves the bijection
\lean{cocycleEquivDeformation : ↥(cocycles W) ≃ FirstOrderDeformation W}. That it is a
\lean{Submodule} --- closed under sums and scalar multiples --- is the statement that the axioms
are linear in $Y_1$. \lean{coboundaries W} is the submodule of coboundaries
(Definition~\ref{l07:def:trivial}); \lean{≤} is inclusion. \lean{deltaCochain W φ} is the cochain
$\delta\varphi$ of \eqref{l07:eq:delta}, defined for \emph{any} linear $\varphi$, graded or not.
\end{dictionary}

\begin{proof}[Proof on paper]
(D2) and (D3): if $\varphi\lsevenvac = 0$ then $(\delta\varphi)(\lsevenvac,n,v) = 0 - \varphi(\lsevenvac_{(n)}v)
+ \lsevenvac_{(n)}\varphi v = -\delta_{n,-1}\varphi v + \delta_{n,-1}\varphi v = 0$, and for $n\ge-1$,
$(\delta\varphi)(u,n,\lsevenvac) = (\varphi u)_{(n)}\lsevenvac - \varphi(u_{(n)}\lsevenvac)$, which is
$0-0$ for $n\ge0$ and $\varphi u - \varphi u$ for $n=-1$. (D1) is immediate from VA5 for the three
terms. For (D4), write $\mathcal{D}$ for the operation that, applied to a trilinear expression
$F(u,v,x)$, returns $F(\varphi u,v,x)+F(u,\varphi v,x)+F(u,v,\varphi x)-\varphi F(u,v,x)$. A direct
expansion shows that substituting $\delta\varphi$ for $Y_1$ in each summand of the left side of
\eqref{l07:eq:jac1} gives $\mathcal{D}$ of the corresponding summand of the undeformed Borcherds
identity (the inner $\varphi$'s cancel in pairs), and likewise on the right. Applying the undeformed
identity at $(\varphi u,v,x)$, $(u,\varphi v,x)$, $(u,v,\varphi x)$ and $(u,v,x)$ concludes. This is the
mode-by-mode form of the statement that $Y^g$ is again a vertex algebra.
\end{proof}

\begin{status}
\lean{coboundaries_le_cocycles}, \lean{deltaCochain_mem_cocycles} and
\lean{cocycleEquivDeformation} are proved; \texttt{\#print axioms} gives
\lean{[propext, Classical.choice, Quot.sound]} (checked for \lean{coboundaries_le_cocycles}, see
Section~\ref{l07:sec:axioms}). Unconditional.
\end{status}

\begin{physics}
Field redefinitions never spoil consistency: if one starts from a consistent set of OPEs and renames
operators, the renamed OPEs are consistent. This is the trivial direction of every deformation theory,
and it is why ``deformations modulo field redefinitions'' is a well-defined space (the tangent space of the
moduli space at the given theory).
\end{physics}

\subsection{Rigidity}

\begin{definition}[Rigidity]\label{l07:def:rigid}
$W$ is \emph{rigid} if every first-order deformation is trivial, and \emph{graded rigid} (in full:
\emph{graded first-order rigid}) if every \emph{graded} first-order deformation is trivial. Only the second
notion is meaningful (Warning~\ref{l07:warn:isrigid}); whenever this chapter says ``graded rigid'' it means
graded first-order rigid in this sense.
\end{definition}
\begin{leancode}
def IsRigid (W : VertexOperatorAlgebra 𝒱) : Prop :=
  ∀ D : FirstOrderDeformation W, D.IsTrivial

def IsRigidGraded (W : VertexOperatorAlgebra 𝒱) : Prop :=
  ∀ D : FirstOrderDeformation W, D.IsGraded → D.IsTrivial
\end{leancode}
(\lean{VOA.Deform.IsRigid}, \leanpath{VOA/\allowbreak Core/\allowbreak Deformation.lean};
\lean{VOA.Deform.IsRigidGraded}, \leanpath{VOA/\allowbreak Core/\allowbreak DeformationGraded.lean}.)

The library also proves the cohomological reading: with $H(W) := \mathrm{cocycles}/\mathrm{coboundaries}$
(an honest quotient module),
\begin{leancode}
abbrev H (W : VertexOperatorAlgebra 𝒱) : Type u := ↥(cocycles W) ⧸ coboundariesIn W

theorem isRigid_iff_H_trivial (W : VertexOperatorAlgebra 𝒱) :
    IsRigid W ↔ ∀ z : H W, z = 0 := by

theorem isRigidGraded_of_gradedCocycles_le (W : VertexOperatorAlgebra 𝒱)
    (h : gradedCocycles W ≤ gradedCoboundaries W) : IsRigidGraded W := by
\end{leancode}
(\lean{VOA.Deform.H}, \lean{VOA.Deform.isRigid_iff_H_trivial},
\lean{VOA.Deform.isRigidGraded_of_gradedCocycles_le}, \leanpath{VOA/\allowbreak Core/\allowbreak DeformationGraded.lean}.
\lean{coboundariesIn W} is \lean{coboundaries W} viewed inside \lean{cocycles W};
\lean{gradedCocycles} and \lean{gradedCoboundaries} are the intersections of the two submodules with
the submodule \lean{gradedCochains W} of graded cochains.)

\begin{warning}\label{l07:warn:isrigid}
\textbf{\lean{IsRigid} is false for every Virasoro vertex operator algebra, for a definitional
reason; only \lean{IsRigidGraded} is meaningful.} The cocycle condition (D1)--(D4) ignores the
grading, while \lean{IsCoboundary} insists on a \emph{grade-preserving} $\varphi$. Consequently
\lean{IsTrivial} implies \lean{IsGraded} (\lean{NonzeroDeform.isGraded_of_isTrivial}: each of the
three terms of $\delta\varphi$ lands in the right weight), and \emph{any} ungraded cocycle is
non-trivial. Such cocycles are easy to produce: $\varphi = L_1$ kills the vacuum, so
$\delta L_1$ is a cocycle (Theorem~\ref{l07:thm:cobcoc}), but $L_1$ lowers weights by one and
$\delta L_1$ is not graded: on $V(c,0)$ one computes $\omega^1_{(-3)}\lsevenvac = -5L_{-3}\lsevenvac$,
of weight $3$, whereas gradedness would put it in weight $2+0-(-3)-1 = 4$. The library proves
\begin{leancode}
theorem isGraded_of_isTrivial {W : VertexOperatorAlgebra 𝒱}
    (D : Deform.FirstOrderDeformation W) (hD : D.IsTrivial) : D.IsGraded := by

def l1Deformation (W : VertexOperatorAlgebra 𝒱) : Deform.FirstOrderDeformation W :=
  deformOf W (L_one_vacuum W)

theorem not_isRigid_virasoroVOA (c : ℂ) : ¬ Deform.IsRigid (virasoroVOA c) :=
  fun h => not_isTrivial_l1Deformation c (h (l1Deformation (virasoroVOA c)))
\end{leancode}
(\lean{VOA.NonzeroDeform.*}, \leanpath{VOA/\allowbreak Examples/\allowbreak NonzeroDeformation.lean}.) The deformation
$\delta L_1$ is of course trivial in the ordinary sense --- it is a change of coordinates by the
\emph{ungraded} map $1+tL_1$ --- so \lean{not_isRigid_virasoroVOA} records an asymmetry of the
definitions, not a modulus of $V(c,0)$. The same argument applies to any VOA with $L_1\ne0$ on
some weight space. The source itself says so (``\lean{Deform.IsRigid} is false for essentially every
VOA and \lean{IsRigid}-shaped goals should not be pursued'') and recommends exactly the repair the
library then adopted: state rigidity as \lean{IsRigidGraded}. Every rigidity theorem of this chapter
concerns \lean{IsRigidGraded}. A reader who meets the word ``rigid'' attached to \lean{IsRigid}
anywhere in the library should treat it as a trap: \lean{IsRigid} must never be quoted as true, and
it is proved to fail for every $V(c,0)$
(\lean{VOA.NonzeroDeform.not_isRigid_virasoroVOA}, \texttt{\#print axioms}-clean).
\end{warning}

\begin{physics}
Rigidity means that the tangent space to the space of chiral algebras at $W$ is zero: every
infinitesimal change of OPE coefficients is a field redefinition. In CFT language, the chiral algebra has
no exactly marginal deformations and is an isolated point of its moduli space (its ``conformal manifold''
is a point). The free boson is the model counterexample for the full CFT (its radius is a modulus), and the
universal Virasoro algebra, whose central charge can be moved, is the model counterexample for chiral
algebras. The graded version restricts to perturbations that preserve scaling dimensions, i.e.\ to
marginal ones, which are the ones relevant for families of conformal theories.
\end{physics}

\subsection{Fidelity of \lean{IsRigidGraded}}\label{l07:sec:fidelity-rigid}

Since the headline theorem is a statement of the form \lean{IsRigidGraded (virasoroSimple c)},
everything rests on this definition, and we examine it from every side.

\begin{dictionary}
\textbf{What \lean{IsRigidGraded W} says, unfolded completely.} For every $\C$-linear
$Y_1:\lsevenV\to\End(\lsevenV)[[z^{\pm1}]]$ satisfying (D1)--(D4) and such that every $u^1_{(n)}$ shifts
weight as $u_{(n)}$ does, there is a linear $\varphi:\lsevenV\to\lsevenV$ with $\varphi(\lsevenvac) = 0$,
$\varphi(\lsevenV_k)\subseteq\lsevenV_k$ for all $k$, and
$u^1_{(n)}v = (\varphi u)_{(n)}v - \varphi(u_{(n)}v) + u_{(n)}\varphi v$ for all $u,v\in\lsevenV$,
$n\in\Z$. One single $\varphi$ works for all $u,n,v$ at once.

\textbf{Which deformations are allowed.} Deformations of the \emph{vertex algebra} structure $Y$
(not just of the VOA; $\omega$ may become non-conformal), over the dual numbers $\C[t]/(t^2)$, on the
fixed graded space $\lsevenV$, with fixed vacuum, compatible with the original weight grading.

\textbf{Which coordinate changes are allowed.} $1+t\varphi$ with $\varphi$ grade-preserving and
$\varphi(\lsevenvac) = 0$.

\textbf{Is the graded gauge group too small?} One might worry that a graded deformation could be
trivialised by an \emph{ungraded} $\varphi$ but not by a graded one, which would make
\lean{IsRigidGraded} stronger than intended (harder to satisfy) --- harmless for a rigidity
\emph{theorem}, but worth knowing. It cannot happen. Argued on paper (not formalised): decompose an
arbitrary linear $\varphi$ with $\varphi\lsevenvac = 0$ into homogeneous components
$\varphi = \sum_d\varphi_d$, $\varphi_d(\lsevenV_k)\subseteq\lsevenV_{k+d}$ (possible because $\lsevenV =
\bigoplus\lsevenV_k$). For homogeneous $u,v$, the three terms of $(\delta\varphi_d)(u,n,v)$ lie in
$\lsevenV_{k+\ell-n-1+d}$, so the weight-$(k+\ell-n-1)$ part of $\delta\varphi(u,n,v)$ is
$\delta\varphi_0(u,n,v)$. If $Y_1 = \delta\varphi$ is graded, it equals its weight-$(k+\ell-n-1)$ part,
so $Y_1 = \delta\varphi_0$ on homogeneous vectors and hence everywhere; and $\varphi_0\lsevenvac$ is the
weight-$0$ part of $\varphi\lsevenvac = 0$. So for graded deformations ``trivial by a graded gauge'' and
``trivial by any gauge'' coincide.

\textbf{Is the graded class of deformations too small?} Restricting to graded deformations is
standard (the cohomology theories in the literature for graded vertex algebras are
grading-restricted), and it is the only sensible choice with this gauge group by
Warning~\ref{l07:warn:isrigid}. It does exclude deformations in which the weight grading itself
moves. For a deformation of $W$ \emph{as a vertex operator algebra}, with deformed conformal vector
$\omega+t\omega_1$, the deformed $L_0$ is diagonalisable on each finite-dimensional
$\lsevenV_n\oplus t\lsevenV_n$-block with the same integer eigenvalues, and one expects to conjugate it back
to the original grading by a coordinate change; we have not formalised this, and the library does
not claim it. The theorem should therefore be read as \emph{graded} rigidity, and the book never
says more.

\textbf{Same as / weaker than / stronger than the literature.} Up to the remarks just made,
\lean{IsRigidGraded W} is the vanishing of the first-order (``$H^2$'') deformation space of $W$ as a
graded vertex algebra with fixed vacuum. It is \emph{weaker} than rigidity among all
(ungraded) deformations with ungraded gauge (a notion the library does not define), and
\emph{stronger} than rigidity of $W$ as a VOA with $\omega$ fixed, because $\omega$ is allowed to
move. It is a first-order statement: it does not by itself say that every formal deformation over
$\C[[t]]$ is trivial (that follows by the usual Gerstenhaber induction for graded deformations with
finite-dimensional weight spaces, an argument we mention but which the library does not contain),
and it says nothing about convergent families.

\textbf{Comparison with Linshaw--Qi.} In the standard terminology, \lean{IsRigidGraded W} is the vanishing
of the \emph{graded} second deformation cohomology $H^2$ of $W$ with the vacuum fixed: graded cocycles modulo
coboundaries of grade-preserving, vacuum-fixing $\varphi$. It is \emph{at least as strong as} the notion of
Linshaw and Qi (arXiv:2601.12017, \S2). The gauge group is the same as theirs (grade-preserving $\varphi$ with
$\varphi\lsevenvac = 0$), while the class of deformations is \emph{larger}: Linshaw--Qi require the
deformation to preserve the translation (derivative) operator and the grading operator, whereas
\lean{FirstOrderDeformation} does not ask $T = L_{-1}$ to be preserved (the trivial deformation $\delta L_0$
already has $\omega^1_{(-2)}\lsevenvac = -L_{-3}\lsevenvac\neq0$) and lets $\omega$ stop being conformal. Since
\lean{IsRigidGraded} asserts triviality for every member of the larger class, with the same notion of
triviality, it implies Linshaw--Qi rigidity. The affine theorem of Chapter~\ref{ch:L09} is therefore
Linshaw--Qi's Theorem~5.1 in their own sense, not merely an analogue of it.

\textbf{How could it have been vacuous?} Three ways, each checked:
\begin{enumerate}
\item \emph{An unsatisfiable or degenerate cocycle condition} would make every VOA rigid. The
  \lean{finsum} trap is excluded by the finiteness lemmas above, and the condition is satisfiable
  by nonzero cochains: the library constructs the scaling deformation $\delta L_0$, which is graded
  with $u^1_{(n)}v = (n+1)u_{(n)}v$ and has $Y_1\neq0$ on $V(c,0)$ for every $c$
  (\lean{NonzeroDeform.exists_nonzero_deformation}, Example~\ref{l07:ex:scaling}).
  So the rigidity theorem has to do real work even on trivial deformations: it must find $\varphi$.
\item \emph{A gauge group so large that everything is trivial.} The gauge group is exactly the
  graded vacuum-fixing changes of coordinates. The library proves that it is \emph{not} large
  enough to absorb a change of the central term: on any VOA with $\lsevenV_2 = \C\omega$ a coboundary
  has $\omega^1_{(1)}\omega = 2\mu\omega$ and $\omega^1_{(3)}\omega = \mu c\lsevenvac$ for one scalar $\mu$
  (Theorem~\ref{l07:thm:gaugeline}). A deformation with $\omega^1_{(1)}\omega = 0$ and
  $\omega^1_{(3)}\omega\ne0$ --- what $\partial_c$ would be --- is provably non-trivial.
\item \emph{A conclusion true for every VOA.} Is \lean{IsRigidGraded} perhaps true for all $W$?
  On paper, no: the universal Virasoro VOA $V(c,0)$ at generic $c$ has the nontrivial graded
  deformation $\partial_c$. In Lean too: the library constructs $\partial_c$ (\lean{VOA.}\allowbreak\lean{VirPoly.}\allowbreak\lean{cDeformation}),
  proves that it is a nontrivial graded deformation, and so proves \lean{¬ IsRigidGraded W} for the concrete
  $W = V(c,0)$, every $c$ (\lean{VOA.}\allowbreak\lean{VirPoly.}\allowbreak\lean{not_}\allowbreak\lean{isRigidGraded_}\allowbreak\lean{virasoroVOA}; Section~\ref{l07:sec:nonvacuity}).
\end{enumerate}
\end{dictionary}

\subsection{Is \texorpdfstring{\lean{IsRigidGraded}}{IsRigidGraded} vacuous? The non-vacuity check}\label{l07:sec:nonvacuity}

A rigidity notion that every vertex operator algebra satisfied would make the headline theorem worthless.
We therefore searched the library (all files assigned to this chapter, and every occurrence of
\lean{¬ IsRigidGraded}, \lean{¬ Deform.IsRigidGraded} and \lean{¬ D.IsTrivial} under \leanpath{VOA/}) for a
\emph{graded} deformation proved \emph{non-trivial}. When this section was first written,
the outcome was: \textbf{there was none}. That has since changed: the files \leanpath{VirasoroPoly*.lean},
\leanpath{VirasoroCDefBuild.lean}, \leanpath{VirasoroCDefJacobi.lean} and \leanpath{VirasoroCDefFinal.lean} under
\leanpath{VOA/\allowbreak Examples/} construct $\partial_cY$ and prove
\begin{leancode}
theorem isCDeformation_cDeformation (c : ℂ) : VirCDef.IsCDeformation c (cDeformation c) :=

theorem not_isRigidGraded_virasoroVOA (c : ℂ) : ¬ Deform.IsRigidGraded (virasoroVOA c) :=
\end{leancode}
(\lean{VOA.}\allowbreak\lean{VirPoly.}\allowbreak\lean{isCDeformation_}\allowbreak\lean{cDeformation}, \lean{VOA.}\allowbreak\lean{VirPoly.}\allowbreak\lean{not_}\allowbreak\lean{isRigidGraded_}\allowbreak\lean{virasoroVOA},
\leanpath{VOA/\allowbreak Examples/\allowbreak VirasoroCDefFinal.lean}; the deformation \lean{cDeformation c} itself is
defined in \leanpath{VOA/\allowbreak Examples/\allowbreak VirasoroCDefJacobi.lean}). Here \lean{virasoroVOA c} is the
\emph{universal} $V(c,0)$, not the simple quotient, and $c\in\C$ is arbitrary. The rest of this section records the
earlier search and the gauge-group half on which the new theorem rests. Apart from $\partial_cY$, the deformations the
library constructs are the zero deformation and $\delta\varphi$ for explicit $\varphi$; those with graded $\varphi$
(e.g.\ $\delta L_0$) are trivial by construction, and those with ungraded $\varphi$ (e.g.\ $\delta L_1$) are non-trivial
only because they are ungraded (Warning~\ref{l07:warn:isrigid}). Every other statement of the form
\lean{¬ IsRigidGraded W} in the library is conditional: \lean{VirCDef.not_isRigidGraded_of_isCDeformation} and
\lean{not_isRigidGraded_of_omega_witness} (\leanpath{VOA/\allowbreak Examples/\allowbreak VirasoroCDeformation.lean}) assume a
deformation (now supplied by \lean{cDeformation}), and \lean{AffineSl2RigidityCollapse.no_uniform_derivation_piece2}
and \lean{…no_uniform_derivation_sl2Basis_tripletW2} (\leanpath{VOA/\allowbreak Examples/\allowbreak AffineSl2RigidityCollapse.lean})
assume a uniform derivation. What \emph{is} proved unconditionally is the half of the non-vacuity argument
that concerns the gauge group:
\begin{leancode}
theorem not_isTrivial_of_omega_witness (D : FirstOrderDeformation (virasoroVOA c))
    (h1 : D.smul₁ (omegaVir c) 1 (omegaVir c) = 0)
    (h3 : D.smul₁ (omegaVir c) 3 (omegaVir c) ≠ 0) : ¬ D.IsTrivial :=
\end{leancode}
(\lean{VOA.VirCDef.not_isTrivial_of_omega_witness}, \leanpath{VOA/\allowbreak Examples/\allowbreak VirasoroCDeformation.lean}.)

\begin{dictionary}
\textbf{Fidelity check of the non-vacuity claim.} The claim to be checked is ``there exist VOAs that are not graded
rigid''. \emph{Machine status:} proved, for $V(c,0)$ at every $c$ (\lean{VOA.}\allowbreak\lean{VirPoly.}\allowbreak\lean{not_}\allowbreak\lean{isRigidGraded_}\allowbreak\lean{virasoroVOA}). The theorem
quoted just above says that \emph{if} a deformation of $V(c,0)$ leaves
$\omega_{(1)}\omega$ unchanged and moves $\omega_{(3)}\omega$, it is not trivial (for every $c$, including $c = 0$), and
\lean{not_isRigidGraded_of_isCDeformation} turns any \emph{graded} such deformation into \lean{¬ IsRigidGraded (virasoroVOA c)}.
The inhabitant --- a graded first-order deformation with
$\omega^1_{(1)}\omega = 0$, $\omega^1_{(3)}\omega = \frac12\lsevenvac$, namely $\partial_cY$ --- is \lean{cDeformation c}, and
\lean{isCDeformation_cDeformation} checks the three conditions. \emph{The argument, which the Lean construction follows:}
the creation monomials in sorted order form a basis of $V(c,0)$ with a $c$-independent index set (proved,
\lean{gradedBasis}); in that basis every structure constant $\langle$basis$\rangle_{(n)}\langle$basis$\rangle$ is a polynomial in $c$
(normal ordering, as in Definition~\ref{l07:def:words}). Transport all $V(c',0)$ to one space via this basis; the
Borcherds identity holds identically in $c'$, so differentiating at $c' = c$ gives a cochain $Y_1 = \partial_cY$ satisfying
(D1)--(D4); it is graded because the basis is homogeneous; its values on $(\omega,\omega)$ are $\partial_c(2\omega) = 0$ and
$\partial_c(\frac c2\lsevenvac) = \frac12\lsevenvac$. By the theorem above it is non-trivial, so $V(c,0)$ is \emph{not} graded rigid, at
every $c$. (Note this concerns the universal $V(c,0)$; at generic $c$ it coincides with $L(c,0)$, so generic simple
Virasoro VOAs are not graded rigid either, while the minimal ones are.) The file \leanpath{VirasoroCDeformation.lean} names two routes
to formalising this; the library took the second (differentiation of polynomial structure constants on the
$c$-independent basis of sorted creation words, \leanpath{VirasoroPolyStructure.lean} to \leanpath{VirasoroCDefFinal.lean}).
\emph{Verdict:} \lean{IsRigidGraded} is non-vacuous, machine-checked: $V(c,0)$ is not graded first-order rigid at any
$c\in\C$. The witness moves the central charge; the library's \lean{FirstOrderDeformation} does not require the conformal
vector or $c$ to be preserved, which is the same convention under which the rigidity theorems of this chapter are proved.
\end{dictionary}

\begin{warning}\label{l07:warn:nonvacuity}
\textbf{Non-vacuity of \lean{IsRigidGraded}: once a paper argument, now machine-checked.} When this chapter was first written, the library contained no proof that any concrete vertex operator algebra fails
\lean{IsRigidGraded}; the files cited in Section~\ref{l07:sec:nonvacuity} now supply one
(\lean{VOA.}\allowbreak\lean{VirPoly.}\allowbreak\lean{not_}\allowbreak\lean{isRigidGraded_}\allowbreak\lean{virasoroVOA}). The conditional statement on which that proof rests is
\begin{leancode}
def IsCDeformation (D : FirstOrderDeformation (virasoroVOA c)) : Prop :=
  D.IsGraded ∧
    D.smul₁ (omegaVir c) 1 (omegaVir c) = 0 ∧
    D.smul₁ (omegaVir c) 3 (omegaVir c) = (2⁻¹ : ℂ) • vacVec c

theorem not_isRigidGraded_of_isCDeformation {D : FirstOrderDeformation (virasoroVOA c)}
    (hD : IsCDeformation c D) : ¬ IsRigidGraded (virasoroVOA c) :=
\end{leancode}
(\lean{VOA.VirCDef.IsCDeformation}, \lean{VOA.VirCDef.not_isRigidGraded_of_isCDeformation},
\leanpath{VOA/\allowbreak Examples/\allowbreak VirasoroCDeformation.lean}): a graded deformation of the universal
$V(c,0)$ that leaves $\omega_{(1)}\omega = 2\omega$ unchanged and moves $\omega_{(3)}\omega = \frac c2\lsevenvac$ at
rate $\frac12\lsevenvac$ --- the derivative $\partial_c$ --- would refute graded rigidity of $V(c,0)$.
The file \leanpath{VirasoroCDeformation.lean} does not itself inhabit \lean{IsCDeformation}; it says so
(``\lean{IsCDeformation c} is a predicate and this file does \textbf{not} inhabit it''; some of its docstrings,
written before the construction, still say that the predicate is inhabited nowhere), and
explains that constructing $\partial_cY$ needs either a square-zero extension
$U(\mathfrak{Vir})/J_2(c)$ with a PBW-type freeness statement, or differentiation of the family
of structure constants on a $c$-independent basis. The second route is the one later carried out, in
\leanpath{VirasoroCDefJacobi.lean} and \leanpath{VirasoroCDefFinal.lean}, which inhabit the predicate
(\lean{VOA.}\allowbreak\lean{VirPoly.}\allowbreak\lean{isCDeformation_}\allowbreak\lean{cDeformation}). Two routes that do \emph{not} produce a
nontrivial graded deformation are also recorded there: the background-charge family of free bosons
(whose $Y$ is independent of the parameter, so the cochain is zero) and naive differentiation of
$c\mapsto V(c,0)$ (the carrier \lean{VirUniv c} depends on $c$, so $\partial_cY$ is not even
type-correct without a trivialisation). The project notes also contain the observation that a
uniform derivation ``$\forall W$, \lean{IsRigidGraded W} $\Rightarrow$ (some normalisation
statement)'' would refute graded rigidity of a free boson (\lean{AffineSl2RigidityCollapse.}\allowbreak
\lean{no_uniform_derivation_piece2}), which again is a conditional, not a witness.

Consequently \lean{IsRigidGraded} fails for $V(c,0)$ at every $c$, machine-checked (the standard argument
that $c\mapsto V(c,0)$ is a nontrivial family, formalised, together with
Theorem~\ref{l07:thm:gaugeline}); since $L(c,0) = V(c,0)$ at generic $c$, it is a
genuine property that the minimal models have and generic Virasoro VOAs lack. As further internal
evidence that the minimal-model theorem has content: its proof uses, at the decisive step, a null
relation in the \emph{simple} quotient together with the simplicity of the quotient; at generic $c$
there is no such relation, and the proof does not apply.
\end{warning}

\begin{remark}[Non-vacuity, in one paragraph]\label{l07:rem:nonvacuity}
Here is the whole non-vacuity argument. Every step of it is now machine-checked; the paragraph marked
\emph{(Construction.)} is the step that was missing when this chapter was first written.
\emph{(Construction.)} The universal $V(c,0)$ has the $c$-independent PBW basis
$L_{-n_1}\cdots L_{-n_k}\lsevenvac$, $n_1\ge\dots\ge n_k\ge2$, in which every structure constant of $Y_c$ is a
polynomial in $c$. Transporting all $V(c',0)$ to one space along this basis and differentiating at $c' = c$
gives $Y_1 := \partial_cY$. It satisfies (D1)--(D4): truncation, since the vanishing bound does not depend on
$c$; the vacuum and creation conditions, since $Y_c(\lsevenvac) = \mathrm{id}$ and $u_{(-1)}\lsevenvac = u$ hold for
every $c$; the linearised Borcherds identity, by differentiating the Borcherds identity in $c$. It is graded,
because the basis is homogeneous, and
\begin{equation*}
 \omega^1_{(1)}\omega = \partial_c(2\omega) = 0,\qquad \omega^1_{(3)}\omega = \partial_c\big(\tfrac c2\lsevenvac\big) = \tfrac12\lsevenvac .
\end{equation*}
So $Y_1$ is a graded first-order deformation. In Lean this is \lean{VOA.}\allowbreak\lean{VirPoly.}\allowbreak\lean{cDeformation}
(\leanpath{VirasoroCDefJacobi.lean}; the linearised Borcherds identity is \lean{Y1_jacobi}), with
\lean{cDeformation_isGraded}, \lean{cDeformation_omega_one} ($\omega^1_{(1)}\omega = 0$) and \lean{cDeformation_omega_three}
($\omega^1_{(3)}\omega = \frac12\lsevenvac$) in \leanpath{VirasoroCDefFinal.lean}. \emph{(Non-triviality.)} On any VOA with $\lsevenV_2 = \C\omega$, every
coboundary $f$ has $(f(\omega)_{(1)}\omega, f(\omega)_{(3)}\omega) = (2\mu\omega,\mu c\lsevenvac)$ for one scalar $\mu$
(\lean{VOA.VirCDef.cmode_conformal_of_isCoboundary}, Theorem~\ref{l07:thm:gaugeline}). For $Y_1$ this would
need $0 = 2\mu$ and $\frac12 = \mu c$, which is impossible; so $Y_1$ is non-trivial. The library proves exactly
this step: \lean{VOA.VirCDef.not_isTrivial_of_omega_witness} (quoted in Section~\ref{l07:sec:nonvacuity}) and
\lean{VOA.VirCDef.not_isRigidGraded_of_isCDeformation} (quoted in Warning~\ref{l07:warn:nonvacuity}), both
\texttt{\#print axioms}-clean. \emph{(Conclusion, in Lean.)} \lean{VOA.}\allowbreak\lean{VirPoly.}\allowbreak\lean{not_}\allowbreak\lean{isRigidGraded_}\allowbreak\lean{virasoroVOA}: $V(c,0)$ is not graded
first-order rigid at any $c$, nor is $L(c,0) = V(c,0)$ at generic $c$ (the deformation moves the central charge, as the
library's notion allows); so the notion separates the minimal-model charges from the
generic ones, as a rigidity theorem should.
\end{remark}

\begin{physics}
In conformal-field-theory terms the remark says: the chiral algebra generated by a stress tensor alone has
one genuine modulus, the central charge, and rescaling $T$ cannot imitate a change of $c$ (it moves the
double pole $2T$ and the quartic pole $c/2$ together). The minimal models kill this modulus through their
null vectors; the generic Virasoro algebra keeps it.
\end{physics}

\begin{status}
Machine-checked: \lean{VOA.VirCDef.cmode_conformal_of_isCoboundary},
\lean{VOA.VirCDef.not_isTrivial_of_omega_witness}, \lean{VOA.VirCDef.not_isRigidGraded_of_omega_witness},
\lean{VOA.VirCDef.not_isRigidGraded_of_isCDeformation}; each reports \lean{[propext, Classical.choice,
Quot.sound]} (re-run 2026-10-27, Section~\ref{l07:sec:axioms}). Also machine-checked since:
\lean{VOA.VirPoly.cDeformation}, \lean{VOA.VirPoly.isCDeformation_cDeformation} and
\lean{VOA.VirPoly.not_isRigidGraded_virasoroVOA} (every $c\in\C$), so a VOA that is not graded first-order rigid
exists in Lean. Their axiom lists are reported as the three standard axioms by the release follow-ups audit
(2026-10-28); \texttt{\#print axioms} was not re-run for this chapter.
\end{status}

\begin{example}[The zero and scaling deformations]\label{l07:ex:scaling}
The zero cochain is a deformation (\lean{Deform.zeroDeformation}), trivial with $\varphi = 0$. More
interestingly, $\varphi = L_0$ is grade-preserving and kills $\lsevenvac$, so $\delta L_0$ is a trivial
graded deformation. On homogeneous $u\in\lsevenV_a$, $v\in\lsevenV_b$,
\begin{equation*}
 (\delta L_0)(u,n,v) = a\,u_{(n)}v - (a+b-n-1)u_{(n)}v + b\,u_{(n)}v = (n+1)\,u_{(n)}v,
\end{equation*}
i.e.\ $Y_1(u,z) = -z\partial_zY(u,z)$: the infinitesimal form of the rescaling $z\mapsto\lambda z$.
In particular $\omega^1_{(1)}\omega = 2\cdot2\omega = 4\omega$ and $\omega^1_{(3)}\omega = 4\cdot\frac c2\lsevenvac = 2c\lsevenvac$,
so this trivial deformation \emph{does} change the structure constant $\omega_{(3)}\omega$ --- together
with $\omega_{(1)}\omega$, in the ratio allowed by Theorem~\ref{l07:thm:gaugeline} ($\mu = 2$). This is the
example that refutes any naive rigidity statement forgetting the quotient by coboundaries.
\begin{leancode}
def scalingDeformation (W : VertexOperatorAlgebra 𝒱) : Deform.FirstOrderDeformation W :=
  deformOf W (L_zero_vacuum W)

theorem smul₁_scalingDeformation_homog (W : VertexOperatorAlgebra 𝒱) {a b : ℤ} {u v : 𝒱}
    (hu : u ∈ W.grading a) (hv : v ∈ W.grading b) (n : ℤ) :
    (scalingDeformation W).smul₁ u n v = ((n : ℂ) + 1) • W.smul u n v := by

theorem exists_nonzero_deformation (c : ℂ) :
    ∃ D : Deform.FirstOrderDeformation (virasoroVOA c), D.Y1 ≠ 0 ∧ D.IsGraded :=
\end{leancode}
(\lean{VOA.NonzeroDeform.*}, \leanpath{VOA/\allowbreak Examples/\allowbreak NonzeroDeformation.lean}; \lean{deformOf W h}
is the deformation with $Y_1 = \delta\varphi$ built from Theorem~\ref{l07:thm:cobcoc}, and
\lean{virasoroVOA c} is the universal Virasoro VOA of Section~\ref{l07:sec:virasoro}. The source
also proves that $\delta L_0$ is the only deformation of the ``mode-diagonal'' form
$u^1_{(n)}v = g(n)u_{(n)}v$, up to scale.)
\end{example}

\subsection{The finite-dimensional picture and its limit}\label{l07:sec:tridegree}

Because a graded $Y_1$ is determined by its components $\lsevenV_a\times\lsevenV_b\to\lsevenV_d$ (in the mode
$n = a+b-d-1$), one might hope to decide rigidity one ``tridegree'' $(a,b,d)$ at a time, by finite
linear algebra. The file \leanpath{VOA/\allowbreak Core/\allowbreak DeformationGraded.lean} sets this up carefully:
\lean{GradedCochain W a b d} is the finite-dimensional space
$\mathrm{Hom}(\lsevenV_a,\mathrm{Hom}(\lsevenV_b,\lsevenV_d))$, \lean{gradedComponent} is the linear component map,
\lean{gradedCocyclesAt} and \lean{gradedCoboundariesAt} are the images of cocycles and coboundaries,
and \lean{DeformationSpaceAt W a b d} is their (finite-dimensional) quotient. It proves the easy
direction --- rigidity implies vanishing at every tridegree --- and then states the converse as a
named \emph{non-theorem}:
\begin{leancode}
theorem gradedCocyclesAt_le_gradedCoboundariesAt_of_isRigid {W : VertexOperatorAlgebra 𝒱}
    (h : IsRigid W) (a b c : ℤ) :
    gradedCocyclesAt W a b c ≤ gradedCoboundariesAt W a b c := by

def GradedToGlobalRigidity (W : VertexOperatorAlgebra 𝒱) : Prop :=
  (∀ a b c : ℤ, gradedCocyclesAt W a b c ≤ gradedCoboundariesAt W a b c) → IsRigid W
\end{leancode}
(\lean{VOA.Deform.*}, \leanpath{VOA/\allowbreak Core/\allowbreak DeformationGraded.lean}.)

\begin{warning}\label{l07:warn:tridegree}
\textbf{Tridegree-wise vanishing does not give rigidity.} The coboundary of $\varphi$ at tridegree
$(a,b,d)$ depends only on the three blocks $\varphi_a,\varphi_b,\varphi_d$ of $\varphi$. Vanishing of
every \lean{DeformationSpaceAt} produces one $\varphi$ \emph{per tridegree}, and nothing forces these
choices to patch into a single $\varphi$. The source gives a concrete instance on $V(c,0)$
(\leanpath{VOA/\allowbreak Examples/\allowbreak VirasoroDeformationLow.lean}): the candidate $\partial_cY$ has component
$\frac12\lsevenvac$ at $(2,2,0)$ and $0$ at $(2,2,2)$; on $\lsevenV_2 = \C\omega$ a gauge acts by one scalar,
$\varphi\omega = \lambda\omega$, and kills it at $(2,2,0)$ only for $\lambda = 1/(2c)$ and at $(2,2,2)$ only
for $\lambda = 0$. Both deformation spaces vanish (for $c\ne0$), yet no single $\varphi$ works. So
\lean{GradedToGlobalRigidity} is \emph{not} proved, is not used, and the headline theorem is stated
against \lean{IsRigidGraded} directly, with one global $\varphi$; it does not inherit this gap. (The
tables of \lean{VirasoroDeformationLow} --- the coboundary space fills the cochain space at the
tridegrees $(2,2,d)$, $d\in\{0,\dots,4\}$, except $(2,2,0)$ at $c = 0$ --- are correct finite
statements, but they are partial information, not rigidity.)
\end{warning}

\subsection{Regular deformations and the pole coefficient}

For completeness we record the remaining definitions of \leanpath{VOA/\allowbreak Core/\allowbreak Deformation.lean},
which serve the affine-algebra rigidity proofs of Chapter~\ref{ch:L09} (the paper of Linshaw and Qi
on $L_{-4/3}(\mathfrak{sl}_2)$ that the deformation layer was first built to formalise). A
deformation is \emph{regular} if $Y_1(a,x)b$ has no negative powers of $x$ for weight-one $a,b$,
i.e.\ $a^1_{(n)}b = 0$ for $n\ge0$; its \emph{pole vector} is the coefficient $a^1_{(1)}b$ of
$x^{-2}$, which for graded $D$ lies in $\lsevenV_0 = \C\lsevenvac$.
\begin{leancode}
def Regular (D : FirstOrderDeformation W) : Prop :=
  ∀ a ∈ W.grading 1, ∀ b ∈ W.grading 1, ∀ n : ℤ, 0 ≤ n → D.smul₁ a n b = 0

def poleVec (D : FirstOrderDeformation W) (a b : 𝒱) : 𝒱 := D.smul₁ a 1 b
\end{leancode}
(\lean{VOA.Deform.FirstOrderDeformation.Regular}, \lean{…poleVec}, same file.) The Virasoro
algebras have $\lsevenV_1 = 0$, so these notions are empty there and play no role in this chapter.
The cautionary tale attached to them --- the statement ``every graded deformation is regular'' is
\emph{false}, refuted by $\delta L_0$, which has $a^1_{(0)}b = a_{(0)}b$ --- is told in
Chapter~\ref{ch:L09}.

%======================================================================
\section{The Virasoro vertex operator algebras, recalled}\label{l07:sec:virasoro}
%======================================================================

This section recalls, briefly but completely, the Virasoro objects that occur in the statements
of this chapter. Proofs (and the construction of the vertex operators) belong to
Chapter~\ref{ch:L06}; here we only need the definitions and a handful of facts, each with its Lean
name.

\subsection{The Virasoro algebra and its vacuum module}

\begin{definition}[Virasoro algebra]\label{l07:def:vir}
The Virasoro algebra $\mathfrak{Vir}$ is the Lie algebra with basis $\{L_n\}_{n\in\Z}\cup\{C\}$ and
bracket
\begin{equation}\label{l07:eq:vir}
 [L_m,L_n] = (m-n)L_{m+n} + \frac{m^3-m}{12}\,\delta_{m+n,0}\,C,\qquad [C,\mathfrak{Vir}] = 0 .
\end{equation}
It is the central extension of the Witt algebra of vector fields $-z^{n+1}\partial_z$ by the
Gelfand--Fuks cocycle.
\end{definition}
\begin{leancode}
abbrev Virasoro : Type := LieAlgebra.ofTwoCocycle gfCocycle

noncomputable def L (n : ℤ) : Virasoro := ofProd gfCocycle (Witt.gen n, 0)

noncomputable def C : Virasoro := ofProd gfCocycle (0, 1)

theorem virasoro_bracket (m n : ℤ) :
    ⁅L m, L n⁆ = ((m : ℂ) - n) • L (m + n)
      + (if m + n = 0 then (((m : ℂ) ^ 3 - m) / 12) • C else 0) := by
\end{leancode}
(\lean{VOA.Virasoro}, \lean{VOA.L}, \lean{VOA.C}, \lean{VOA.virasoro_bracket},
\leanpath{VOA/\allowbreak Foundations/\allowbreak Virasoro.lean}. \lean{LieAlgebra.ofTwoCocycle} is Mathlib's central
extension of a Lie algebra by a $2$-cocycle; \lean{gfCocycle} is the Gelfand--Fuks cocycle on the
library's Witt algebra; \lean{ofProd} identifies the extension with the product
$\mathrm{Witt}\times\C$ as a vector space.)

\begin{physics}
$L_n$ are the modes of the stress tensor $T(z)=\sum L_nz^{-n-2}$: $L_{-1}$ generates translations, $L_0$
dilatations (the Hamiltonian in radial quantisation), $L_1$ special conformal transformations. The central
term is the conformal anomaly: $c$ measures the response of the theory to curvature (the Casimir energy on
a cylinder is $-\pi c/(6L)$) and counts degrees of freedom (a free boson has $c=1$, a free fermion $c=1/2$).
\end{physics}

\begin{definition}[Universal vacuum module]\label{l07:def:Vc}
For $c\in\C$, let $J(c)\subseteq U(\mathfrak{Vir})$ be the left ideal generated by the $L_n$ with
$n\ge-1$ and by $C-c$. The \emph{universal Virasoro vacuum module} is $V(c,0) := U(\mathfrak{Vir})/J(c)$,
with vacuum $\lsevenvac = [1]$ and $\mathfrak{Vir}$ acting by left multiplication. Write $L_n$ also for the
operator by which $L_n$ acts, and $\omega := L_{-2}\lsevenvac$.
\end{definition}
\begin{leancode}
abbrev UVir : Type := UniversalEnvelopingAlgebra ℂ Virasoro

def vacGens (c : ℂ) : Set UVir :=
  {a | (∃ n : ℤ, -1 ≤ n ∧ a = ι ℂ (L n)) ∨ a = ι ℂ C - algebraMap ℂ UVir c}

noncomputable def vacIdeal (c : ℂ) : Submodule UVir UVir :=
  Submodule.span _ (vacGens c)

abbrev VirUniv (c : ℂ) : Type := UVir ⧸ vacIdeal c

noncomputable def virVacRho (c : ℂ) : Virasoro →ₗ⁅ℂ⁆ Module.End ℂ (VirUniv c) :=
  (Algebra.lsmul ℂ ℂ (VirUniv c)).toLieHom.comp (ι ℂ)

noncomputable def vacVec (c : ℂ) : VirUniv c := Submodule.Quotient.mk 1
\end{leancode}
(\lean{VOA.VirModule.*}, \leanpath{VOA/\allowbreak Foundations/\allowbreak VirasoroModule.lean}.)
\begin{leancode}
noncomputable def Lop (n : ℤ) : End ℂ (VirUniv c) := virVacRho c (L n)

noncomputable def omegaVir : VirUniv c := Lop c (-2) (vacVec c)
\end{leancode}
(\lean{VOA.VirModule.Lop}, \lean{VOA.VirModule.omegaVir}, \leanpath{VOA/\allowbreak Examples/\allowbreak VirasoroField.lean}.)

\begin{dictionary}
\lean{UniversalEnvelopingAlgebra ℂ Virasoro} is Mathlib's $U(\mathfrak{Vir})$, with \lean{ι ℂ} the
canonical map $\mathfrak{Vir}\to U(\mathfrak{Vir})$. \lean{Submodule UVir UVir} spanned by a set is
the \emph{left} ideal it generates (a submodule of $U$ over itself acting by left multiplication),
and \lean{⧸} is the quotient module; so \lean{VirUniv c} is $V(c,0)$ exactly as in
Definition~\ref{l07:def:Vc}. \lean{virVacRho c} is the Lie homomorphism $\mathfrak{Vir}\to\End(V(c,0))$
given by left multiplication, and \lean{Lop c n} is the operator $L_n$ on $V(c,0)$. The
docstring of \lean{VirUniv} records a corrected naming error: this is the \emph{universal} module,
often written $V(c,0)$ or $\mathrm{Vir}^c$, \emph{not} the simple $L(c,0)$; the two differ exactly at
the degenerate central charges, which is the whole point of this chapter. The library also proves
$C$ acts by $c$ (\lean{vac_acts_central}) and $\lsevenvac\neq0$ for every $c$
(\lean{vacVec_ne_zero}, without using \lean{VOA.pbw}).
\end{dictionary}

By the Poincar\'e--Birkhoff--Witt theorem $V(c,0)$ has basis the ``creation monomials''
$L_{-n_1}\cdots L_{-n_k}\lsevenvac$ with $n_1\ge\cdots\ge n_k\ge2$; the monomial has $L_0$-eigenvalue
(weight) $n_1+\cdots+n_k$. The library encodes a monomial by the list $[n_1-2,\dots,n_k-2]$ of
natural numbers, defines the weight and the monomials, and grades $V(c,0)$ by the span of the
monomials of each weight:
\begin{leancode}
def energyVir : List ℕ → ℕ
  | [] => 0
  | (i :: rest) => (i + 2) + energyVir rest

noncomputable def negWordVir : List ℕ → VirUniv c
  | [] => vacVec c
  | n :: ws => Lop c (-(n : ℤ) - 2) (negWordVir ws)

noncomputable def gradingVir (n : ℤ) : Submodule ℂ (VirUniv c) :=
  Submodule.span ℂ {x | ∃ ws : List ℕ, (energyVir ws : ℤ) = n ∧ negWordVir c ws = x}
\end{leancode}
(\lean{VOA.VirModule.energyVir}, \leanpath{VOA/\allowbreak Examples/\allowbreak VirasoroGrading.lean};
\lean{VOA.VirModule.negWordVir}, \leanpath{VOA/\allowbreak Examples/\allowbreak VirasoroSpanning.lean};
\lean{VOA.VirModule.gradingVir}, \leanpath{VOA/\allowbreak Examples/\allowbreak VirasoroGrading.lean}.)
The \emph{sorted} lists of weight $n$ (partitions of $n$ into parts $\ge2$) index a basis
\lean{gradedBasis c n} of the weight-$n$ space, with $c$-independent index type
\lean{SortedOfEnergy n}; this is proved, \lean{pbw}-free, in
\leanpath{VOA/\allowbreak Examples/\allowbreak VirasoroPartitionBasis.lean}. So
\begin{equation*}
 \dim V(c,0)_n = \#\{\text{partitions of $n$ into parts}\ge2\}:\qquad 1,0,1,1,2,2,4,4,7,8,12,\dots
 \quad(n = 0,1,2,\dots),
\end{equation*}
independently of $c$. In particular $V(c,0)_1 = 0$, $V(c,0)_2 = \C\omega$, $V(c,0)_3 = \C L_{-3}\lsevenvac$,
$V(c,0)_4 = \mathrm{span}\{L_{-2}^2\lsevenvac,L_{-4}\lsevenvac\}$ (\lean{VirDefLow.gradingVir_two},
\lean{gradingVir_three}, \lean{LeeYang.gradingVir_four}).

\begin{physics}
$V(c,0)$ is the space of states of a CFT of central charge $c$ that are created from the vacuum by the stress
tensor alone: $L_{-n_1}\cdots L_{-n_k}\lsevenvac$ corresponds to the field $\,{:}\partial^{n_1-2}T\cdots\partial^{n_k-2}T{:}\,$
(up to normalisation). The relations $L_n\lsevenvac=0$ for $n\ge-1$ say the vacuum is invariant under the global
conformal group $\mathrm{SL}_2$ generated by $L_{-1},L_0,L_1$. ``Universal'' means no further relations are
imposed: every possible descendant of the identity is kept, even those of zero norm.
\end{physics}

\begin{definition}[Universal Virasoro VOA]\label{l07:def:virVOA}
$V(c,0)$ with the grading above, vacuum $\lsevenvac$, conformal vector $\omega = L_{-2}\lsevenvac$, central
charge $c$, and the unique vertex operators with $Y(\omega,z) = \sum_nL_nz^{-n-2}$, is a vertex operator
algebra of CFT type (Chapter~\ref{ch:L06}).
\end{definition}
\begin{leancode}
noncomputable def virasoroVOA : VertexOperatorAlgebra (VirUniv c) where
  grading := gradingVir c
  isInternal := gradingVir_isInternal c
  finrank := gradingVir_finiteDimensional c
  Y := virYA c
  vacuum := vacVec c
  vacuum_mem := vacVec_mem_gradingVir_zero c
  grading_mode := fun u hu v hv n => virYA_grading_mode c u hu v hv n
  vacuum_Y := virYA_vacuum c
  creation := virYA_creation c
  locality := virYA_local c
  regularity := virYA_regularity c
  conformalVector := omegaVir c
  conformal_mem := omegaVir_mem_gradingVir_two c
  centralCharge := c
  virasoro := virYA_virasoro c
  L0_grading := fun {_} v hv => virYA_L0_grading c v hv
  translation := virYA_omega_translation c
  vacuum_grade := gradingVir_zero c
  positivity := gradingVir_neg c
\end{leancode}
(\lean{VOA.VirModule.virasoroVOA}, \leanpath{VOA/\allowbreak Examples/\allowbreak VirasoroInstance.lean}.) The
vertex operator \lean{virYA c} is built by the reconstruction theorem of Chapter~\ref{ch:L03}
(Theorem~\ref{l03:thm:reconstruct}) from the field of $\omega$; we do not need its definition, only
that its conformal modes are the operators \lean{Lop c n} (\lean{virasoroVOA_L_eq_Lop}) and the
facts listed above.

\begin{physics}
This is the chiral algebra generated by the stress tensor alone, with the OPE $T(z)T(w)\sim
\frac{c/2}{(z-w)^4}+\frac{2T(w)}{(z-w)^2}+\frac{\partial T(w)}{z-w}$. Every two-dimensional CFT of central charge
$c$ contains (a quotient of) it: it is the ``identity sector'' or vacuum Verma module of the theory.
\end{physics}

\subsection{The simple quotient \texorpdfstring{$L(c,0)$}{L(c,0)}}

\begin{definition}[Simple Virasoro VOA]\label{l07:def:Lc}
$L(c,0)$ is the quotient of $V(c,0)$ by its maximal proper (graded) ideal; it is a simple vertex
operator algebra of central charge $c$.
\end{definition}
\begin{leancode}
def maxQuotientVOA (W : VertexOperatorAlgebra 𝒱) :
    VertexOperatorAlgebra (𝒱 ⧸ maxGradedIdeal W) :=
  quotientVOA W (maxGradedIdeal W) (maxGradedIdeal_isIdeal W) (maxGradedIdeal_isGradedIdeal W)

noncomputable def maxIdealVir (c : ℂ) : Submodule ℂ (VirUniv c) :=
  maxGradedIdeal (virasoroVOA c)

noncomputable def virasoroSimple (c : ℂ) :
    VertexOperatorAlgebra (VirUniv c ⧸ maxIdealVir c) :=
  maxQuotientVOA (virasoroVOA c)

theorem virasoroSimple_isSimple : (virasoroSimple c).IsSimple :=
  isSimple_maxQuotientVOA' (virasoroVOA c) (vacVec_ne_zero c)

theorem virasoroSimple_vacuum_ne_zero : (virasoroSimple c).vacuum ≠ 0 := by
\end{leancode}
(\lean{VOA.QuotVOA.maxQuotientVOA}, \leanpath{VOA/\allowbreak Core/\allowbreak QuotientVOA.lean};
\lean{VOA.VirSimple.maxIdealVir}, \lean{…virasoroSimple}, \lean{…virasoroSimple_isSimple},
\lean{…virasoroSimple_vacuum_ne_zero}, \leanpath{VOA/\allowbreak Examples/\allowbreak VirasoroSimple.lean}.)

\begin{dictionary}
\lean{maxGradedIdeal W} is the sum of all graded ideals not containing $\lsevenvac$, and
\lean{quotientVOA} the quotient VOA, both from Chapter~\ref{ch:L04} (Definition~\ref{l04:def:quot}).
\lean{IsSimple} is Definition~\ref{l04:def:simple}: nonzero, and only the ideals $0$ and
everything. Since every ideal of a VOA is graded (\lean{QuotVOA.isGradedIdeal_of_isIdeal}), the
``graded'' in \lean{maxGradedIdeal} costs nothing, and \lean{virasoroSimple c} is the standard
$L(c,0)$. Its grading is the image grading (\lean{virasoroSimple_grading}), its central charge is
$c$ (\lean{virasoroSimple_centralCharge}, by \lean{rfl}), its vacuum is nonzero, and its weight-one
space is zero (\lean{virasoroSimple_grading_one}). For $c\ne0$ its weight-two space is $\C\omega$
with $\omega\ne0$; for $c = 0$ the vector $\omega$ lies in the maximal ideal (it has norm $\frac c2 = 0$)
and $L(0,0) = \C\lsevenvac$ is the one-dimensional trivial VOA (\lean{zeroCharge_conformal_eq_zero},
Section~\ref{l07:sec:zero}). The statement \lean{virasoroSimple_grading_two} --- weight two is
$\C\cdot\omega$ --- holds for all $c$, with $\C\cdot0 = 0$ when $c = 0$.
\end{dictionary}

\begin{physics}
Passing to $L(c,0)$ is setting null states to zero. In a CFT, a state of zero norm that is orthogonal to everything
decouples from all correlation functions, and consistency (e.g.\ unitarity, or crossing symmetry in BPZ's
analysis) requires it to vanish; the resulting null-vector decoupling equations are the differential
equations that solve the minimal models. $L(c,0)$ is the vacuum sector of the CFT actually realised.
\end{physics}

\begin{definition}[Minimal-model central charges]\label{l07:def:cpq}
For natural numbers $p,q$, $\lsevencpq := 1 - 6(p-q)^2/(pq)$.
\end{definition}
\begin{leancode}
def minimalCentralCharge (p q : ℕ) : ℚ := 1 - 6 * ((p : ℚ) - q) ^ 2 / (p * q)
\end{leancode}
(\lean{VOA.VirModule.minimalCentralCharge}, \leanpath{VOA/\allowbreak Examples/\allowbreak VirasoroVOA.lean}.)
\begin{dictionary}
The value is a rational number; in \lean{virasoroSimple (minimalCentralCharge p q)} it is coerced
to $\C$ (Lean inserts the cast $\Q\to\C$ silently). Division by $pq = 0$ would give $0$ in Lean's
convention, but the theorems below always assume $p,q\ge2$. The formula is symmetric in $p,q$
(\lean{minimalCharge_swap}). Examples: $c_{2,3} = 0$, $c_{2,5} = -22/5$, $c_{3,4} = 1/2$,
$c_{2,7} = -68/7$, $c_{4,5} = 7/10$, $c_{5,6} = 4/5$. Every coprime pair $p,q\ge2$ gives a value
$<1$, and these values accumulate at $1$ (take $p-q$ fixed, $p\to\infty$).
\end{dictionary}

\begin{physics}
These are the central charges of the BPZ minimal models, the conformal theories with finitely many primary
fields under the Virasoro algebra. $c_{3,4}=1/2$ is the critical Ising model, $c_{4,5}=7/10$ the tricritical
Ising model, $c_{5,6}=4/5$ the three-state Potts model (as a Virasoro theory), $c_{2,5}=-22/5$ the Lee--Yang edge
singularity; the unitary series $c_{m,m+1}$ (Friedan--Qiu--Shenker) are the multicritical Ising points. The set is
discrete away from $c=1$, which is why the minimal models come in a list rather than in families.
\end{physics}

\subsection{The invariant form and the Kac determinant}\label{l07:sec:kac}

By Li's theorem (Theorem~\ref{l04:thm:canonical} of Chapter~\ref{ch:L04}) $V(c,0)$ carries a unique
invariant symmetric bilinear form with $\langle\lsevenvac,\lsevenvac\rangle = 1$; it satisfies
$\langle L_na,b\rangle = \langle a,L_{-n}b\rangle$. In Lean it is \lean{VirGram4.canB c}, and its
analogue on $L(c,0)$ is \lean{VirQuotForm.Bsimp c}:
\begin{leancode}
noncomputable def canB : InvariantForm (virasoroVOA c) :=
  canonicalInvariantForm (vacVec_ne_zero c) (virasoroVOA_qp c)

noncomputable def Bsimp : InvariantForm (virasoroSimple c) :=
  canonicalInvariantForm (virasoroSimple_vacuum_ne_zero c) (virasoroSimple_qp c)
\end{leancode}
(\leanpath{VOA/\allowbreak Examples/\allowbreak VirasoroGram4.lean}, \leanpath{VOA/\allowbreak Examples/\allowbreak VirasoroQuotientForm.lean}.)
\lean{InvariantForm} and \lean{canonicalInvariantForm} are Definition~\ref{l04:def:invform} and
Theorem~\ref{l04:thm:canonical}; the two arguments discharge the hypotheses $\lsevenvac\ne0$ and
$L_1\lsevenV_1 = 0$ (vacuous, as $\lsevenV_1 = 0$). The form is normalised (\lean{canB_normalized},
\lean{Bsimp_normalized}: $\langle\lsevenvac,\lsevenvac\rangle = 1$) and fully invariant.

Because $\langle L_na,b\rangle = \langle a,L_{-n}b\rangle$ and $L_n\lsevenvac = 0$ for $n\ge-1$, the pairing of two
creation monomials can be computed by moving the modes of the left monomial to the right as positive
modes and commuting them through with \eqref{l07:eq:vir} until they hit the vacuum. The result is a
polynomial in $c$. For example
\begin{equation*}
 \langle\omega,\omega\rangle = \langle\lsevenvac,L_2L_{-2}\lsevenvac\rangle = \tfrac c2,\qquad
 \langle L_{-3}\lsevenvac,L_{-3}\lsevenvac\rangle = 2c,
\end{equation*}
and at weight $4$, in the basis $(L_{-2}^2\lsevenvac, L_{-4}\lsevenvac)$, the Gram matrix is
\begin{equation}\label{l07:eq:gram4}
 G_4(c) = \begin{pmatrix} \tfrac{c(c+8)}{2} & 3c\\ 3c & 5c\end{pmatrix},\qquad
 \det G_4(c) = \tfrac{c^2(5c+22)}{2}.
\end{equation}
The \emph{level-$n$ vacuum Gram determinant} $D_n(c) := \det G_n(c)$ is a polynomial in $c$. The
radical of the form is the maximal ideal (Theorem~\ref{l04:thm:lirad}), so
$\dim L(c,0)_n = \mathrm{rank}\,G_n(c)$, and $L(c,0)\ne V(c,0)$ exactly when some $D_n(c)$ vanishes. The
classical result (Kac, Feigin--Fuchs; for the vacuum module, Gorelik--Kac) is that $D_n(c)$ vanishes
exactly at minimal-model charges $c = c_{p,q}$ with $(p-1)(q-1)\le n$, and that at the \emph{first} such
level $n = (p-1)(q-1)$ the zero at $c = \lsevencpq$ is \emph{simple}. For example $D_4$ vanishes simply at
$-22/5 = c_{2,5}$ and doubly at $0 = c_{2,3}$ (where the first singular level is $2$ and $D_2 = c/2$ has a
simple zero). This simple-zero statement is the one representation-theoretic input of the rigidity proof;
Section~\ref{l07:sec:all} states exactly which Lean theorems supply it.

\begin{example}[The Lee--Yang null vector]\label{l07:ex:leeyang}
At $c = -22/5$ the vector $s = 5L_{-2}^2\lsevenvac - 3L_{-4}\lsevenvac$ spans the kernel of $G_4(c)$:
indeed $5\cdot\frac{c(c+8)}2 - 3\cdot3c = \frac{5c^2+22c}{2} = 0$ and $5\cdot3c-3\cdot5c = 0$. So $s$ is
in the radical, hence in the maximal ideal, and $s = 0$ in $L(-22/5,0)$; $\dim L(-22/5,0)_4 = 1$. Its
norm as a polynomial in $c$ is
\begin{equation*}
 P(c) = \langle s,s\rangle_c = 25\cdot\tfrac{c(c+8)}2 - 2\cdot15\cdot3c + 9\cdot5c
 = \tfrac{25}{2}c^2 + 55c = \tfrac{5c}{2}(5c+22),
\end{equation*}
with $P(-22/5) = 0$ and $P'(-22/5) = 25\cdot(-\tfrac{22}{5}) + 55 = -55\ne0$. This number $-55$ is the
whole Lee--Yang rigidity theorem (Section~\ref{l07:sec:leeyang}).
\end{example}

%======================================================================
\section{Anatomy of a graded deformation of \texorpdfstring{$L(c,0)$}{L(c,0)}}\label{l07:sec:anatomy}
%======================================================================

From now on, unless said otherwise, $W = L(c,0)$, $\lsevenV$ is its state space, and
$D$ is a graded first-order deformation of $W$ with products $u^1_{(n)}v$. This section carries out
step 1 of the heuristic of Section~\ref{l07:sec:motivation}: it shows that $D$ produces, on the
first-order space $\lsevenV\oplus t\lsevenV$, an honest representation of the Virasoro algebra whose central
charge is $c+t\delta$ for a single number $\delta$, and that $\delta$ is the only gauge-invariant
information at the level of $\omega$.

\subsection{Two scalars and the gauge line}

Since $\lsevenV_0 = \C\lsevenvac$, $\lsevenV_1 = 0$ and $\lsevenV_2 = \C\omega$, gradedness of $D$ forces
\begin{equation*}
 \omega^1_{(3)}\omega\in\lsevenV_{2+2-3-1} = \lsevenV_0,\quad
 \omega^1_{(2)}\omega\in\lsevenV_1 = 0,\quad
 \omega^1_{(1)}\omega\in\lsevenV_2,\quad
 \omega^1_{(n)}\omega\in\lsevenV_{3-n} = 0\ (n\ge4),
\end{equation*}
so there are scalars $\alpha,\beta$ with
\begin{equation}\label{l07:eq:alphabeta}
 \omega^1_{(3)}\omega = \alpha\lsevenvac,\qquad \omega^1_{(1)}\omega = \beta\omega .
\end{equation}
Compare the undeformed values $\omega_{(3)}\omega = L_2L_{-2}\lsevenvac = \frac c2\lsevenvac$ and
$\omega_{(1)}\omega = L_0\omega = 2\omega$: to first order the deformation changes ``$\frac c2$'' by
$\alpha$ and ``$2$'' by $\beta$.
\begin{leancode}
theorem exists_conformal_parameters (D : FirstOrderDeformation W) (hD : D.IsGraded)
    (h2 : W.grading 2 = ℂ ∙ W.conformalVector) :
    ∃ α β : ℂ,
      D.smul₁ W.conformalVector 3 W.conformalVector = α • W.vacuum ∧
      D.smul₁ W.conformalVector 1 W.conformalVector = β • W.conformalVector := by
\end{leancode}
(\lean{VOA.VirDeformModes.exists_conformal_parameters}, \leanpath{VOA/\allowbreak Examples/\allowbreak VirasoroDeformationModes.lean};
$W$ is an arbitrary VOA here, and \lean{h2} is the hypothesis $\lsevenV_2 = \C\omega$.)

\begin{theorem}[The gauge line]\label{l07:thm:gaugeline}
Let $W$ be any VOA with $\lsevenV_2 = \C\omega$ and let $f$ be a coboundary. Then there is $\mu\in\C$
with $f(\omega)_{(1)}\omega = 2\mu\,\omega$ and $f(\omega)_{(3)}\omega = \mu c\,\lsevenvac$. Consequently, for a
trivial graded deformation the scalars of \eqref{l07:eq:alphabeta} satisfy $\alpha = \mu c$,
$\beta = 2\mu$, hence
\begin{equation}\label{l07:eq:delta-def}
 \delta := 2\alpha - \beta c = 0 .
\end{equation}
\end{theorem}
\begin{leancode}
theorem cmode_conformal_of_isCoboundary (W : VertexOperatorAlgebra 𝒱)
    (hV2 : W.grading 2 = ℂ ∙ W.conformalVector) {f : Cochain 𝒱} (hf : IsCoboundary W f) :
    ∃ μ : ℂ,
      cmode f W.conformalVector 1 W.conformalVector = (2 * μ) • W.conformalVector ∧
      cmode f W.conformalVector 3 W.conformalVector = (μ * W.centralCharge) • W.vacuum := by
\end{leancode}
(\lean{VOA.VirCDef.cmode_conformal_of_isCoboundary}, \leanpath{VOA/\allowbreak Examples/\allowbreak VirasoroCDeformation.lean};
\lean{cmode f u n v} is \lean{mode (f u) n v}, i.e.\ $f(u)_{(n)}v$.)

\begin{proof}[Proof on paper]
$f = \delta\varphi$ with $\varphi$ grade-preserving, $\varphi\lsevenvac = 0$. Then $\varphi\omega = \mu\omega$ for
some $\mu$ (as $\lsevenV_2 = \C\omega$) and $\varphi$ kills $\lsevenV_0 = \C\lsevenvac$. Hence
$(\delta\varphi)(\omega,1,\omega) = \mu\cdot2\omega - \varphi(2\omega) + \mu\cdot2\omega = 2\mu\omega$ and
$(\delta\varphi)(\omega,3,\omega) = \mu\frac c2\lsevenvac - \varphi(\frac c2\lsevenvac) + \mu\frac c2\lsevenvac = \mu c\lsevenvac$.
\end{proof}

\begin{dictionary}
This is the precise form of the physicist's ``rescaling $T$ cannot change $c$ alone'': the image
of the gauge group in the two slots $(\omega_{(1)}\omega,\omega_{(3)}\omega)$ is the line
$\C\cdot(2\omega, c\lsevenvac)$. The library also proves the converse (every point of the line is attained,
by $\frac\mu2\delta L_0$: \lean{VirCDef.exists_isCoboundary_on_gauge_line}), and for $c\ne0$ an
``iff'' on $L(c,0)$ (\lean{VirRigidStruct.attained_by_coboundary_iff}). The combination
$\delta = 2\alpha-\beta c$ is the linear functional vanishing on the line; it is the \emph{central
variation} of the deformation, and the rest of the proof is about showing $\delta = 0$ and that
$\delta = 0$ suffices. Caution: \eqref{l07:eq:alphabeta} are only two of infinitely many components
of $Y_1$ on $\omega\otimes\omega$ (the components $\omega^1_{(n)}\omega$ for $n\le0$ land in $\lsevenV_{3-n}$
and are unconstrained so far). Reducing all of them to $(\alpha,\beta)$ is exactly what the
reconstruction theorem of Section~\ref{l07:sec:reconstruction} does; the project notes record the
earlier claim ``the $(\omega,\omega)$ deformation data is two-dimensional'' as an undercount.
\end{dictionary}

\begin{status}
\lean{exists_conformal_parameters} and \lean{cmode_conformal_of_isCoboundary} are proved for an
arbitrary VOA; axioms \lean{[propext, Classical.choice, Quot.sound]}. Unconditional.
\end{status}

\begin{physics}
This is the statement ``rescaling $T$ cannot change $c$ alone''. A field redefinition $T\mapsto(1+t\mu)T$ multiplies
the fourth-order pole $c/2$ by $(1+t\mu)^2$ but also the second-order pole $2T$ by $(1+t\mu)$; only the ratio
$c/2:2$, i.e.\ $\delta=2\alpha-\beta c$, is redefinition-invariant. So in the chiral algebra generated by $T$ the
central charge is a genuine coordinate on the space of theories, not an artefact of normalisation --- in
contrast with the Heisenberg level, which rescaling the current removes.
\end{physics}

\subsection{The first-order space and its operators}

A state of the deformed algebra is $v + tv_1$, which the library writes as the pair
$(v,v_1)\in\lsevenV\times\lsevenV$. Multiplication by $t$ is the square-zero map
$\lsevenN(v,v_1) = (0,v)$, and a dual number $a+tb$ acts by $a + b\lsevenN$. The mode of a ``constant''
state $a = (a,0)$ acts by the block matrix
\begin{equation}\label{l07:eq:lifted}
 a^t_{(n)}(x,x_1) = \big(a_{(n)}x,\ a^1_{(n)}x + a_{(n)}x_1\big),
\end{equation}
which is \eqref{l07:eq:Yt} written in coordinates.
\begin{leancode}
def liftedMode (a : V) (n : ℤ) : End ℂ (V × V) where
  toFun x := (W.smul a n x.1, D.smul₁ a n x.1 + W.smul a n x.2)

def tangentShift : End ℂ (V × V) where
  toFun x := (0, x.1)

def dualScalarAction : DualNumber ℂ →ₐ[ℂ] End ℂ (V × V) where
  toFun z := z.fst • (1 : End ℂ (V × V)) + z.snd • tangentShift
\end{leancode}
(\lean{VOA.Deform.liftedMode}, \leanpath{VOA/\allowbreak Core/\allowbreak DeformationLiftedModes.lean};
\lean{VOA.Deform.tangentShift}, \lean{VOA.Deform.dualScalarAction},
\leanpath{VOA/\allowbreak Core/\allowbreak DeformationCentralModes.lean}; linearity proofs omitted.
\lean{DualNumber ℂ} is Mathlib's $\C[\varepsilon]/(\varepsilon^2)$, implemented as the trivial square-zero
extension, with \lean{z.fst}, \lean{z.snd} the two coordinates of $z = a+b\varepsilon$.)

The key observation is that the Taylor expansion of polynomials happens automatically:
for $p\in\C[X]$,
\begin{equation}\label{l07:eq:taylor}
 p(k\cdot1 + \delta\lsevenN) = p(k)\cdot1 + p'(k)\,\delta\,\lsevenN,
\end{equation}
because $\lsevenN^2 = 0$ (\lean{Deform.dualScalarAction_polynomial}). A central operator of the form
$k+\delta\lsevenN$ therefore records a ``central charge'' $k$ together with its first-order variation
$\delta$, and any polynomial identity in the central charge acquires a tangent coefficient
$p'(k)\delta$.

\subsection{The normalised Virasoro modes}

Write $L^1_n := \omega^1_{(n+1)}$ for the correction to $L_n = \omega_{(n+1)}$. The order-$t$ commutator
formula (the $\ell = 0$ slice of \eqref{l07:eq:jac1}) gives
\begin{equation*}
 [L^1_m,L_n] + [L_m,L^1_n] = \sum_{i\ge0}\binom{m+1}{i}\Big((\omega^1_{(i)}\omega)_{(m+n+2-i)}
   + (\omega_{(i)}\omega)^1_{(m+n+2-i)}\Big).
\end{equation*}
Only $i\in\{0,1,3\}$ contribute ($i = 2$ and $i\ge4$ land in $\lsevenV_1 = 0$ or negative weight), and the
$i = 1,3$ terms are known by \eqref{l07:eq:alphabeta}. The $i = 0$ term involves the unknown
$\omega^1_{(0)}\omega\in\lsevenV_3$ and the correction of $\omega_{(0)}\omega = L_{-1}\omega$; but
antisymmetry of the left side under $m\leftrightarrow n$ determines it (at $m = -1$; the potential
central term vanishes there because $\binom{2}{3} = 0$). The result, proved as
\lean{mixedComm_virasoro}, is
\begin{equation*}
 [L^1_m,L_n]+[L_m,L^1_n] = (m-n)L^1_{m+n} + \tfrac{(m-n)\beta}{2}L_{m+n}
 + \delta_{m+n,0}\tfrac{m^3-m}{6}\,\alpha .
\end{equation*}
The term $\frac{(m-n)\beta}2L_{m+n}$ is the effect of a rescaling of $\omega$; it is removed by the
normalisation $\widetilde L_n := L_n + t(L^1_n - \frac\beta2L_n)$. Then $\widetilde L_n$ satisfy the Virasoro
relations with central term $\frac{m^3-m}{12}(c + t(2\alpha-\beta c))$:
\begin{equation}\label{l07:eq:normcomm}
 [\widetilde L_m,\widetilde L_n] = (m-n)\widetilde L_{m+n} + \delta_{m+n,0}\frac{m^3-m}{12}\big(c\cdot1 + \delta\,\lsevenN\big),
 \qquad \delta = 2\alpha-\beta c .
\end{equation}
\begin{leancode}
noncomputable def normalizedL (D : FirstOrderDeformation W) (β : ℂ) (n : ℤ) : End ℂ (V × V) :=
  liftedMode D W.conformalVector (n+1) -
    (β/2) • (tangentShift * liftedMode D W.conformalVector (n+1))

theorem normalizedL_comm (D : FirstOrderDeformation W) (hD : D.IsGraded)
    (h1 : W.grading 1 = ⊥) (α β : ℂ)
    (ha : D.smul₁ W.conformalVector 3 W.conformalVector = α • W.vacuum)
    (hb : D.smul₁ W.conformalVector 1 W.conformalVector = β • W.conformalVector)
    (m n : ℤ) :
    ⁅normalizedL D β m, normalizedL D β n⁆ =
      ((m : ℂ)-n) • normalizedL D β (m+n) +
      if m+n=0 then (((m : ℂ)^3-m)/12) •
        (W.centralCharge • (1 : End ℂ (V × V)) +
          (2*α-β*W.centralCharge) • tangentShift) else 0 := by

theorem normalizedL_vacuum (D : FirstOrderDeformation W) (β : ℂ)
    (n : ℤ) (hn : -1 ≤ n) : normalizedL D β n (W.vacuum, 0) = 0 := by

theorem normalizedL_tangent (D : FirstOrderDeformation W) (β : ℂ) (n : ℤ) (x : V) :
    normalizedL D β n (0, x) = (0, W.L n x) := by

theorem normalizedL_neg_two_vacuum (D : FirstOrderDeformation W) (β : ℂ) :
    normalizedL D β (-2) (W.vacuum, 0) =
      (W.conformalVector, -(β/2) • W.conformalVector) := by
\end{leancode}
(\lean{VOA.VirDeformModes.*}, \leanpath{VOA/\allowbreak Examples/\allowbreak VirasoroDeformationModes.lean}; here $W$ is
any VOA with $\lsevenV_1 = 0$, hypothesis \lean{h1}.)

\begin{dictionary}
\lean{liftedMode D W.conformalVector (n+1)} is $\omega^t_{(n+1)} = L_n + tL^1_n$ in the form
\eqref{l07:eq:lifted}; \lean{tangentShift * liftedMode …} is $t\cdot\omega^t_{(n+1)} = tL_n$ (since $t^2 = 0$);
so \lean{normalizedL D β n} is $\widetilde L_n$. The bracket \lean{⁅·,·⁆} is the commutator in $\End(\lsevenV\times\lsevenV)$.
\lean{normalizedL_comm} is \eqref{l07:eq:normcomm}: the central term is
$\frac{m^3-m}{12}$ times the operator $c\cdot1 + (2\alpha-\beta c)\lsevenN$. \lean{normalizedL_vacuum}: the
constant vacuum $(\lsevenvac,0)$ is killed by $\widetilde L_n$ for $n\ge-1$ (from (D3)).
\lean{normalizedL_tangent}: on pure tangent states $(0,x)$ the normalised modes act by the
\emph{undeformed} $L_n$ --- this is the formal content of ``$t\cdot(\text{correction}) = 0$'' and will be
the crux of the obstruction. \lean{normalizedL_neg_two_vacuum}: the normalised conformal state is
$\widetilde\omega = (\omega,-\frac\beta2\omega)$.

\textbf{No hidden input.} The hypotheses are exactly: $D$ graded, $\lsevenV_1 = 0$, and the defining
equations of $\alpha,\beta$. In particular no deformed translation law, no deformed invariant form and
no Virasoro relation for the deformed modes is assumed; \eqref{l07:eq:normcomm} is derived from (D4).
\end{dictionary}

The library packages \eqref{l07:eq:normcomm} as a structure and instantiates it for every
$L(c,0)$:
\begin{leancode}
structure CentralVirasoroData (E : Type*) [AddCommGroup E] [Module ℂ E] where
  L : ℤ → End ℂ E
  central : End ℂ E
  comm : ∀ m n, ⁅L m, L n⁆ = ((m : ℂ)-n) • L (m+n) +
    if m+n=0 then (((m : ℂ)^3-m)/12) • central else 0
  central_comm : ∀ n, Commute central (L n)

theorem virasoroSimple_exists_normalizedData (c : ℂ)
    (D : FirstOrderDeformation (VirSimple.virasoroSimple c)) (hD : D.IsGraded) :
    ∃ (α β : ℂ) (A : CentralVirasoroData
        ((VirModule.VirUniv c ⧸ VirSimple.maxIdealVir c) ×
          (VirModule.VirUniv c ⧸ VirSimple.maxIdealVir c))),
      D.smul₁ (VirSimple.virasoroSimple c).conformalVector 3
        (VirSimple.virasoroSimple c).conformalVector = α • (VirSimple.virasoroSimple c).vacuum ∧
      D.smul₁ (VirSimple.virasoroSimple c).conformalVector 1
        (VirSimple.virasoroSimple c).conformalVector = β • (VirSimple.virasoroSimple c).conformalVector ∧
      A.L = normalizedL D β ∧
      A.central = c • 1 + (2*α-β*c) • tangentShift := by
\end{leancode}
(\lean{VOA.VirDeformModes.CentralVirasoroData}, \lean{…virasoroSimple_exists_normalizedData}, same file.)

\begin{dictionary}
\lean{CentralVirasoroData E} is a representation of the Virasoro relations on $E$ in which the central
element acts by an operator \lean{central} commuting with all $L_n$ (rather than by a scalar). The
theorem says: for every central charge $c$ (including $c = 0$, where $\omega = 0$ in $L(0,0)$ and
$\alpha = 0$ is forced) and every graded $D$ of $L(c,0)$ there are $\alpha,\beta$ as in \eqref{l07:eq:alphabeta}
and such a representation on $\lsevenV\times\lsevenV$, given by the normalised modes, with central operator
$c+\delta\lsevenN$. The existential quantifier over $\alpha,\beta$ is harmless: they are uniquely determined
by $D$ (because $\lsevenvac\ne0$ and, for $c\neq0$, $\omega\neq0$), and every later theorem takes them as
arguments together with their defining equations.
\end{dictionary}

\begin{status}
\lean{normalizedL_comm}, \lean{virasoroSimple_exists_normalizedData}: proved, axioms
\lean{[propext, Classical.choice, Quot.sound]}. Unconditional.
\end{status}

%======================================================================
\section{Reconstruction: zero central variation implies triviality}\label{l07:sec:reconstruction}
%======================================================================

This section carries out step 1 of the heuristic in full: a graded deformation of $L(c,0)$ whose
central variation $\delta$ vanishes is trivial. The argument has a general part, valid for any
\emph{simple} VOA, and a Virasoro part that produces the required subspace.

\subsection{Closed lifts in a simple VOA}

\begin{definition}[Deformed product, closed lift]\label{l07:def:closedlift}
On $\lsevenV\times\lsevenV$ (states $x = x_0 + tx_1$) define the deformed $n$-th product
\begin{equation*}
 x\star_n y := \big(x_{0(n)}y_0,\ x_{0(n)}y_1 + x_{1(n)}y_0 + x_0{}^1_{(n)}y_0\big),
\end{equation*}
the product of $Y^t$ written in coordinates. A subspace $S\subseteq\lsevenV\times\lsevenV$ is \emph{mode-closed}
if $x\star_ny\in S$ for $x,y\in S$; its \emph{reduction} is the first projection $S\to\lsevenV$; its
\emph{tangent kernel} is $\{v : (0,v)\in S\}$; it is \emph{grade-closed} if it contains the homogeneous
components $(x_{0,n},x_{1,n})$ of each of its elements.
\end{definition}
\begin{leancode}
def product (x : V × V) (n : ℤ) (y : V × V) : V × V :=
  (W.smul x.1 n y.1,
    W.smul x.1 n y.2 + W.smul x.2 n y.1 + D.smul₁ x.1 n y.1)

def ModeClosed (S : Submodule ℂ (V × V)) : Prop :=
  ∀ x ∈ S, ∀ y ∈ S, ∀ n : ℤ, product D x n y ∈ S

def reduction (S : Submodule ℂ (V × V)) : S →ₗ[ℂ] V :=
  (LinearMap.fst ℂ V V).comp S.subtype

def tangentKernel (S : Submodule ℂ (V × V)) : Submodule ℂ V :=
  S.comap (LinearMap.inr ℂ V V)

def GradeClosed (W : VertexOperatorAlgebra V) (S : Submodule ℂ (V × V)) : Prop :=
  ∀ x ∈ S, ∀ n : ℤ, (gradePart W n x.1, gradePart W n x.2) ∈ S
\end{leancode}
(\lean{VOA.Deform.Reconstruction.*}, \leanpath{VOA/\allowbreak Core/\allowbreak DeformationReconstruction.lean};
\lean{gradePart W n} is the projection $\lsevenV\to\lsevenV_n$ along the other weight spaces.)

\begin{physics}
$\lsevenV\times\lsevenV$ is the space of states of the perturbed theory, with $(x,x_1)$ the state $x+tx_1$, and $\star_n$ is
its perturbed OPE. A closed lift is a candidate for ``the same operator algebra written in a new basis'':
a family of perturbed operators, one over each unperturbed operator, closed under the perturbed OPE.
\end{physics}

\begin{theorem}[Reconstruction from a closed lift]\label{l07:thm:closedlift}
Let $W$ be simple and $D$ a first-order deformation. Suppose $S\subseteq\lsevenV\times\lsevenV$ is mode-closed,
grade-closed, contains $(\lsevenvac,0)$, does not contain $(0,\lsevenvac)$, and its reduction is surjective.
Then $D$ is trivial.
\end{theorem}
\begin{leancode}
theorem isTrivial_of_closed_lift (S : Submodule ℂ (V × V))
    (hsimple : W.IsSimple) (hclosed : ModeClosed D S)
    (hsurj : Function.Surjective (reduction S))
    (hsep : (0, W.vacuum) ∉ S) (hvac : (W.vacuum, 0) ∈ S)
    (hgrade : GradeClosed W S) : D.IsTrivial := by
\end{leancode}
(\lean{VOA.Deform.Reconstruction.isTrivial_of_closed_lift}, same file.)

\begin{proof}[Proof on paper]
\emph{The tangent kernel $K$ is an ideal.} If $v\in K$ and $u\in\lsevenV$, choose $x\in S$ with $x_0 = u$
(surjectivity); then $x\star_n(0,v) = (0,u_{(n)}v)\in S$, so $u_{(n)}v\in K$. \emph{It is zero.} By simplicity
$K = 0$ or $K = \lsevenV$; the latter would give $(0,\lsevenvac)\in S$. \emph{So $S$ is a graph.} The reduction is
injective (its kernel is $\{(0,v)\in S\} = 0$) and surjective, so $S = \{(v,\psi v)\}$ for a unique linear
$\psi:\lsevenV\to\lsevenV$. \emph{Properties of $\psi$.} $(\lsevenvac,0)\in S$ gives $\psi\lsevenvac = 0$; grade-closure gives
$\psi(\lsevenV_n)\subseteq\lsevenV_n$; mode-closure applied to $(u,\psi u)\star_n(v,\psi v)$ gives
$\psi(u_{(n)}v) = u_{(n)}\psi v + (\psi u)_{(n)}v + u^1_{(n)}v$, i.e.\ $u^1_{(n)}v = \delta(-\psi)(u,n,v)$.
So $D$ is the coboundary of $\varphi = -\psi$.
\end{proof}

\begin{dictionary}
\textbf{Hypotheses.} \lean{hsimple} is simplicity of $W$ (Definition~\ref{l04:def:simple}) --- this
is where the \emph{simple} quotient is essential; the universal $V(c,0)$ is not simple at the minimal
charges, and at generic $c$ a closed lift with the separation property need not exist (the
$\partial_c$-deformation has none). \lean{hclosed}, \lean{hsurj}, \lean{hsep}, \lean{hvac},
\lean{hgrade} are the five properties of the definition. \lean{Function.Surjective (reduction S)} is
surjectivity of the first projection restricted to $S$. The conclusion \lean{D.IsTrivial} is
Definition~\ref{l07:def:trivial}, with one global $\varphi$. Note that $\lsevenV\times\lsevenV$ itself is \emph{not}
assumed to be of CFT type --- its weight-zero part is two-dimensional, which is exactly why separation
(\lean{hsep}) is a hypothesis.

\textbf{Interpretation.} $S$ is the would-be image of an isomorphism $(\lsevenV_\lsevenR,Y)\to(\lsevenV_\lsevenR,Y^t)$; the
theorem turns it into the coordinate change. This is the ``reconstruction'' half of rigidity, and it is
completely general.
\end{dictionary}

\begin{status}
Proved for every simple VOA and every deformation; axioms \lean{[propext, Classical.choice, Quot.sound]}.
The same file proves the variant \lean{isTrivial_of_generated_separation} for the smallest closed lift
containing a generating set, where only separation remains to be checked.
\end{status}

\begin{physics}
If, over every operator of the unperturbed theory, one can find a perturbed operator such that the
perturbed operators close among themselves under the perturbed OPE and contain the identity but not ``$t$
times the identity'', then the perturbation is a field redefinition. Simplicity of $W$ is used the way a
physicist uses irreducibility of the vacuum sector: there are no nontrivial subsectors in which a would-be
obstruction could hide. Heuristic content: a perturbation is trivial exactly when the perturbed theory is
isomorphic to the unperturbed one.
\end{physics}

\subsection{The closed lift of a fixed-charge deformation}

Now let $W = L(c,0)$ and suppose $\delta = 0$, so the normalised modes $\widetilde L_n$ satisfy the
Virasoro relations with \emph{scalar} central charge $c$ and kill $(\lsevenvac,0)$ for $n\ge-1$. By the
universal property of $V(c,0) = U(\mathfrak{Vir})/J(c)$ there is a unique $\mathfrak{Vir}$-module map
\begin{equation*}
 F : V(c,0)\to\lsevenV\times\lsevenV,\qquad F(\lsevenvac) = (\lsevenvac,0),\qquad F(L_nv) = \widetilde L_nF(v)
\end{equation*}
(\lean{VirDeformModes.fixedEval}). Let $S = F(V(c,0))$. Then:
\begin{enumerate}
\item the first component of $F$ is the projection $\pi : V(c,0)\to L(c,0)$, because the first component
  of $\widetilde L_n$ is $L_n$; hence the reduction of $S$ is surjective;
\item $F$ maps $V(c,0)_n$ into $\lsevenV_n\times\lsevenV_n$ (the normalised modes shift both coordinates by the
  expected weight, \lean{normalizedL_mem_grading}), hence $S$ is grade-closed;
\item separation: the weight-zero part of $S$ is $F(\C\lsevenvac) = \C(\lsevenvac,0)$, which does not contain
  $(0,\lsevenvac)$;
\item mode-closure: $S$ is generated from $(\lsevenvac,0)$ by the modes of the corrected conformal state
  $\widetilde\omega = (\omega,-\frac\beta2\omega)$, whose deformed products are exactly the $\widetilde L_n$
  (\lean{corrected_conformal_product}); the order-$t$ iterate formula (the $m = 0$ slice of
  \eqref{l07:eq:jac1}, proved on $\lsevenV\times\lsevenV$ as \lean{Deform.Reconstruction.product_iterate}) shows that the states whose
  deformed modes preserve $S$ are closed under deformed products, so all of $S$ preserves $S$.
\end{enumerate}
Theorem~\ref{l07:thm:closedlift} applies, using simplicity of $L(c,0)$.

\begin{theorem}[Fixed-charge reconstruction]\label{l07:thm:zerodelta}
Let $D$ be a graded first-order deformation of $L(c,0)$, $c\in\C$ arbitrary, with scalars $\alpha,\beta$
as in \eqref{l07:eq:alphabeta}. If $2\alpha-\beta c = 0$ then $D$ is trivial.
\end{theorem}
\begin{leancode}
variable (c : ℂ) (D : FirstOrderDeformation (virasoroSimple c))

theorem isTrivial_of_zero_central_variation (hD : D.IsGraded) (α β : ℂ)
    (ha : D.smul₁ (virasoroSimple c).conformalVector 3
      (virasoroSimple c).conformalVector = α • (virasoroSimple c).vacuum)
    (hb : D.smul₁ (virasoroSimple c).conformalVector 1
      (virasoroSimple c).conformalVector = β • (virasoroSimple c).conformalVector)
    (hz : 2*α-β*c = 0) : D.IsTrivial := by
\end{leancode}
(\lean{VOA.VirDeformModes.isTrivial_of_zero_central_variation},
\leanpath{VOA/\allowbreak Examples/\allowbreak VirasoroFixedReconstruction.lean}. This file is assigned to
Chapter~\ref{ch:L06} but its theorem is indispensable here; its companion \lean{isTrivial_of_fixed_modes}
takes the \lean{CentralVirasoroData} with central operator $c\cdot1$ directly.)

\begin{dictionary}
All hypotheses are visible: gradedness, the defining equations of $\alpha$ and $\beta$, and $\delta = 0$.
There is no hypothesis on $c$ (it holds at every central charge, degenerate or not) and none on the
representation theory of $L(c,0)$ beyond what the library proves (simplicity, $\lsevenV_1 = 0$, $\lsevenV_2 = \C\omega$,
the universal property of $V(c,0)$). So the theorem reduces graded rigidity of $L(c,0)$ to the single
scalar equation $\delta = 0$:
\begin{equation*}
 \text{\lean{IsRigidGraded (virasoroSimple c)}}\iff\text{every graded $D$ has }2\alpha-\beta c = 0,
\end{equation*}
the ``$\Leftarrow$'' by this theorem and ``$\Rightarrow$'' by the gauge line
(Theorem~\ref{l07:thm:gaugeline}). This is the precise sense in which ``the only possible modulus of
$L(c,0)$ is the central charge''.
\end{dictionary}

\begin{status}
Proved, axioms \lean{[propext, Classical.choice, Quot.sound]}; no \lean{VOA.pbw}. Unconditional.
\end{status}
\begin{physics}
``The Virasoro vacuum sector is determined by $c$.'' If a perturbation of $L(c,0)$ does not change the central
charge, then the perturbed stress tensor still satisfies the Virasoro algebra with the same $c$, all states are
still its descendants, and the perturbed theory is the same theory in new coordinates. This is the
first-order form of the familiar fact that the $TT$ OPE has only one free parameter.
\end{physics}


%======================================================================
\section{The obstruction: null relations force zero central variation}\label{l07:sec:obstruction}
%======================================================================

We now carry out steps 2 and 3 of the heuristic: at a central charge where $L(c,0)$ has a relation
whose norm polynomial has a simple zero, every graded deformation has $\delta = 0$. Two formulations
appear in the library: an early, purely linear-algebraic ``bridge'' through deformed bilinear forms,
which turned out to be the wrong tool for the simple quotient, and the word-pairing obstruction that
the final proof uses. We present both, the first as an instructive near-miss.

\subsection{Deformed bilinear forms and Jantzen layers}

Let $E$ be a finite-dimensional vector space with a symmetric bilinear form $B_0$ and another bilinear
form $B_1$, thought of as the first-order deformation $B_t = B_0+tB_1$. A vector $s$ in the radical of
$B_0$ (a ``relation at $t = 0$'') \emph{lifts to first order} if there is $s_1$ with
$(B_0+tB_1)(s+ts_1,\cdot) = 0$ modulo $t^2$, i.e.\ $B_1(s,v) + B_0(s_1,v) = 0$ for all $v$.

\begin{proposition}[Lifting criterion]\label{l07:prop:lift}
For symmetric $B_0$ on finite-dimensional $E$ and any $s\in E$: $s$ lifts to first order if and only if
$B_1(s,v) = 0$ for all $v$ in the radical of $B_0$. Hence, writing $F_1 = \mathrm{rad}\,B_0$ and
$F_2 = \{s\in F_1 : B_1(s,F_1) = 0\}$ (the first two Jantzen layers), $F_2$ is exactly the set of
relations that lift.
\end{proposition}
\begin{leancode}
def radBil (B : E →ₗ[K] E →ₗ[K] K) : Submodule K E := LinearMap.ker B

theorem lift_iff_pairing_zero [FiniteDimensional K E] {B₀ B₁ : E →ₗ[K] E →ₗ[K] K}
    (hB₀ : ∀ x y : E, B₀ x y = B₀ y x) (s : E) :
    (∃ s₁ : E, ∀ v : E, B₁ s v + B₀ s₁ v = 0) ↔ ∀ v ∈ radBil B₀, B₁ s v = 0 := by

def F1 (B₀ : E →ₗ[K] E →ₗ[K] K) : Submodule K E := radBil B₀

def F2 (B₀ B₁ : E →ₗ[K] E →ₗ[K] K) : Submodule K E where
  carrier := {s | s ∈ F1 B₀ ∧ ∀ v ∈ F1 B₀, B₁ s v = 0}
\end{leancode}
(\lean{VOA.DefForm.*}, \leanpath{VOA/\allowbreak Core/\allowbreak DeformedForm.lean}, over an arbitrary field
\lean{K}; submodule proofs of \lean{F2} omitted.)

\begin{proof}[Proof on paper]
``$\Rightarrow$'': if $v\in\mathrm{rad}B_0$ then $B_0(s_1,v) = B_0(v,s_1) = 0$ by symmetry. ``$\Leftarrow$'': the
functional $-B_1(s,\cdot)$ kills $\mathrm{rad}B_0$; for symmetric $B_0$ on a finite-dimensional space the
functionals of the form $B_0(x,\cdot)$ are exactly those killing $\mathrm{rad}B_0$ (dimension count), so
$-B_1(s,\cdot) = B_0(s_1,\cdot)$ for some $s_1$.
\end{proof}

The \emph{rigidity bridge} applies this to $E = \lsevenV_k$ with $B_0 = $ the level-$k$ Shapovalov form
$\langle a,b\rangle_k := \varepsilon(a_{(2k-1)}b)$ (where $a_{(2k-1)}b\in\lsevenV_0 = \C\lsevenvac$ and $\varepsilon$ reads off
the coefficient of $\lsevenvac$) and $B_1 = \langle a,b\rangle^1_k := \varepsilon(a^1_{(2k-1)}b)$, the same formula
with $Y_1$. It proves that the Shapovalov form is symmetric (from skew-symmetry), agrees up to the sign
$(-1)^k$ with the canonical invariant form on quasi-primary vectors, and contains every Virasoro
singular vector of a descendant-spanned weight space in its radical. Its main theorem is:

\begin{leancode}
theorem F2_lt_F1_of_self_pairing_ne_zero (hvac : W.vacuum ≠ 0)
    (D : FirstOrderDeformation W) (hD : D.IsGraded) {k : ℤ} {s : W.grading k}
    (hs : s ∈ DefForm.radBil (shapForm W hvac k))
    (hself : shapForm₁ D hD hvac k s s ≠ 0) :
    DefForm.F2 (shapForm W hvac k) (shapForm₁ D hD hvac k) < DefForm.F1 (shapForm W hvac k) :=
\end{leancode}
(\lean{VOA.RigidBridge.F2_lt_F1_of_self_pairing_ne_zero}, \leanpath{VOA/\allowbreak Core/\allowbreak RigidityBridge.lean};
\lean{shapForm W hvac k} and \lean{shapForm₁ D hD hvac k} are the two forms just described, defined
via \lean{vacCoeff W hvac : W.grading 0 →ₗ[ℂ] ℂ}, the inverse of $\lambda\mapsto\lambda\lsevenvac$.)

\begin{dictionary}
In words: if $s$ is a relation of the level-$k$ Shapovalov form and the \emph{deformed} self-pairing
$\langle s,s\rangle^1_k$ is nonzero, then $s$ does not survive to first order, so $F_2\subsetneq F_1$ (\lean{<} on
submodules is strict inclusion). The file's own docstring states loudly that this is an obstruction
statement, not rigidity: ``Nothing here proves \lean{IsRigid}''.
\end{dictionary}

\begin{warning}\label{l07:warn:bridge}
\textbf{The bridge hypothesis is unsatisfiable exactly where rigidity holds.} The library proves
(\lean{NonzeroDeform.shapForm₁_self_eq_zero_of_isTrivial}) that for every \emph{trivial} graded deformation
the deformed self-pairing of a radical vector vanishes:
\noindent File \leanpath{VOA/Examples/NonzeroDeformation.lean}.
\begin{leancode}
theorem shapForm₁_self_eq_zero_of_isTrivial (W : VertexOperatorAlgebra 𝒱)
    (hvac : W.vacuum ≠ 0) (D : Deform.FirstOrderDeformation W) (hD : D.IsGraded)
    (hT : D.IsTrivial) {k : ℤ} {s : W.grading k}
    (hs : s ∈ DefForm.radBil (shapForm W hvac k)) :
    shapForm₁ D hD hvac k s s = 0 := by
\end{leancode}
So the hypothesis \lean{hself} can only be met by a non-trivial graded deformation, and on a graded-rigid
VOA it is never met. Moreover on the simple quotient $L(c,0)$ the Shapovalov form is nondegenerate, so
its radical is zero and the bridge has nothing to act on. The bridge is a correct theorem about the
universal $V(c,0)$ (``a singular vector of $V(c_{p,q},0)$ does not survive a deformation that moves its
norm''), but it is not a route to rigidity of $L(c,0)$, and the final proof does not use it. It is kept
because the linear algebra of Proposition~\ref{l07:prop:lift} is correct and because it documents a
recorded dead end: the project notes also refute an earlier ``Jantzen bridge lemma''
$\mathrm{rad}_R = F_2\oplus\varepsilon F_1$, which negated its own hypothesis, and that false lemma appears
nowhere in the library.
\end{warning}

\begin{physics}
In a family of theories, the Gram matrix of states becomes $B_0+tB_1$. A null state at $t=0$ survives to first order
(remains null after a small correction) exactly when its first-order change of norm, paired against the other
null states, vanishes. This is the first step of the Jantzen filtration, the representation-theorist's way of
following null vectors along a family of central charges or highest weights.
\end{physics}

\subsection{Words, their adjoints, and the norm polynomial}

The correct mechanism is not a deformed form on $L(c,0)$ but a computation in the representation
$\widetilde L_n$ of Section~\ref{l07:sec:anatomy}, paired with the \emph{undeformed} invariant form.

\begin{definition}[Words and norm polynomials]\label{l07:def:words}
A \emph{word} is a finite list $w = [m_1,\dots,m_k]$ of integers; in a representation $L$ of the
Virasoro relations it acts by $L_w := L_{m_1}\cdots L_{m_k}$. Its \emph{adjoint} is
$w^\dagger := [-m_k,\dots,-m_1]$ (reverse and negate), so that $\langle L_wx,y\rangle = \langle x,L_{w^\dagger}y\rangle$ for an
invariant form. For a word $w$, the \emph{vacuum polynomial} $P_w\in\C[X]$ is defined by normal
ordering: commute positive and zero modes to the right using \eqref{l07:eq:vir}, with the central
element kept as a formal variable $X$, until they annihilate the vacuum; $P_w$ is the coefficient of the
empty creation word. For a finite combination $s = \sum_ir_iL_{w_i}\lsevenvac$ the \emph{norm polynomial} is
\begin{equation*}
 P_{w,r}(X) := \sum_{i,j}r_ir_j\,P_{w_i^\dagger w_j}(X),\qquad\text{so that}\qquad P_{w,r}(c) = \langle s,s\rangle_{V(c,0)} .
\end{equation*}
\end{definition}
\begin{leancode}
noncomputable def virWord (L : ℤ → End ℂ V) : List ℤ → End ℂ V
  | [] => 1
  | n::w => L n * virWord L w

theorem virWord_adjoint {W : VertexOperatorAlgebra V} (B : InvariantForm W)
    (w : List ℤ) (x y : V) :
    B.form (virWord W.L w x) y =
      B.form x (virWord W.L (w.reverse.map Neg.neg) y) := by

noncomputable def centralPolynomial : NormalForm →ₗ[ℂ] Polynomial ℂ :=
  Finsupp.linearCombination ℂ (fun t => if t.2=[] then X^t.1 else 0)

noncomputable def wordPolynomial (w : List ℤ) : Polynomial ℂ :=
  centralPolynomial (wordNormal w)

noncomputable def relationPolynomial (w : ι → List ℤ) (r : ι → ℂ) : Polynomial ℂ :=
  ∑ i, ∑ j, (r i * r j) • wordPolynomial ((w i).reverse.map Neg.neg ++ w j)
\end{leancode}
(\lean{VOA.VirDeformModes.virWord}, \lean{…virWord_adjoint}, \leanpath{VOA/\allowbreak Examples/\allowbreak VirasoroWordPairing.lean};
\lean{VOA.VirDeformModes.NormalOrdering.centralPolynomial}, \lean{…wordPolynomial},
\leanpath{VOA/\allowbreak Examples/\allowbreak VirasoroUniversalPairing.lean}; \lean{…relationPolynomial},
\leanpath{VOA/\allowbreak Examples/\allowbreak VirasoroUniformObstruction.lean}.)

\begin{dictionary}
\lean{virWord L w} is $L_w$ (\lean{1} is the identity, \lean{*} composition, so $L_{[m]++w} = L_m\circ L_w$).
\lean{w.reverse.map Neg.neg} is $w^\dagger$ and \lean{++} is concatenation, so
\lean{(w i).reverse.map Neg.neg ++ w j} is $w_i^\dagger w_j$. \lean{wordNormal w} is the normal-ordered form
of $L_w\lsevenvac$, a finite formal combination of (central power, creation word) labels built by a recursive
commutation algorithm (\leanpath{VOA/\allowbreak Examples/\allowbreak VirasoroNormalOrdering.lean}); \lean{centralPolynomial}
keeps the terms with empty creation word, turning central power $k$ into $X^k$. \lean{relationPolynomial w r}
is $P_{w,r}$; the index type \lean{ι} is any finite type.

\textbf{Why this is the right polynomial.} Two theorems pin down the meaning, so the reader need not
trust the normal-ordering algorithm: (i) at every $c$, $P_{u^\dagger v}(c) = \langle L_u\lsevenvac,L_v\lsevenvac\rangle$ in
$V(c,0)$ for the canonical form (\lean{universal_gram_entry}), hence
$P_{w,r}(c) = \langle s,s\rangle$ (\lean{relationPolynomial_eval}); (ii) in \emph{any}
\lean{CentralVirasoroData} with a vacuum-like vector, the normal-ordering evaluation reproduces
$L_wv$ (\lean{evaluate_wordNormal}). Both are proved; there is no free choice in $P_{w,r}$.
\end{dictionary}
\begin{leancode}
theorem universal_gram_entry (c : ℂ) (u v : List ℤ) :
    (wordPolynomial (u.reverse.map Neg.neg ++ v)).eval c =
      (VirGram4.canB c).form
        (virWord (VirModule.virasoroVOA c).L u (VirModule.vacVec c))
        (virWord (VirModule.virasoroVOA c).L v (VirModule.vacVec c)) := by
\end{leancode}
(\lean{VOA.VirDeformModes.NormalOrdering.universal_gram_entry}, \leanpath{VOA/\allowbreak Examples/\allowbreak VirasoroUniversalPairing.lean}.)

\begin{physics}
A word is a product of Virasoro generators; its vacuum polynomial is the vacuum expectation value
$\langle0|L_{m_1}\cdots L_{m_k}|0\rangle$ computed with only the Virasoro algebra, as a polynomial in $c$. The norm
polynomial of a combination $s$ is the two-point function $\langle s|s\rangle$ of the corresponding descendant of the
identity, as a function of the central charge. Its zeros are the central charges at which $s$ becomes a null
state --- the vacuum row of the Kac table.
\end{physics}

\subsection{The word-pairing obstruction}

\begin{theorem}[Relation obstruction]\label{l07:thm:relation}
Let $c\in\C$. Suppose $s := \sum_ir_iL_{w_i}\lsevenvac = 0$ in $L(c,0)$ for a finite family of words $w_i$ and
coefficients $r_i$, and suppose the norm polynomial satisfies $P'_{w,r}(c)\ne0$. Then every graded
first-order deformation of $L(c,0)$ has central variation $\delta = 0$, and therefore
$L(c,0)$ is graded rigid.
\end{theorem}
\begin{leancode}
theorem virWord_pairing_obstruction {W : VertexOperatorAlgebra V}
    (D : FirstOrderDeformation W) (β : ℂ) (B : InvariantForm W)
    {ι : Type*} [Fintype ι] (w : ι → List ℤ) (r : ι → ℂ)
    (hr : ∑ i, r i • virWord W.L (w i) W.vacuum = 0) :
    B.form W.vacuum
      ((∑ i, r i • virWord (normalizedL D β) ((w i).reverse.map Neg.neg))
        (∑ j, r j • virWord (normalizedL D β) (w j) (W.vacuum,0))).2 = 0 := by

theorem simple_normalized_word_tangent (c : ℂ)
    (D : FirstOrderDeformation (VirSimple.virasoroSimple c)) (hD : D.IsGraded)
    (α β : ℂ)
    (ha : D.smul₁ (VirSimple.virasoroSimple c).conformalVector 3
      (VirSimple.virasoroSimple c).conformalVector = α • (VirSimple.virasoroSimple c).vacuum)
    (hb : D.smul₁ (VirSimple.virasoroSimple c).conformalVector 1
      (VirSimple.virasoroSimple c).conformalVector = β • (VirSimple.virasoroSimple c).conformalVector)
    (w : List ℤ) :
    (VirQuotForm.Bsimp c).form (VirSimple.virasoroSimple c).vacuum
      (virWord (normalizedL D β) w ((VirSimple.virasoroSimple c).vacuum,0)).2 =
        (wordPolynomial w).derivative.eval c * (2*α-β*c) := by

theorem graded_rigidity_of_relation (c : ℂ) (w : ι → List ℤ) (r : ι → ℂ)
    (hr : ∑ i, r i • virWord (VirSimple.virasoroSimple c).L (w i)
      (VirSimple.virasoroSimple c).vacuum = 0)
    (hd : (relationPolynomial w r).derivative.eval c ≠ 0) :
    IsRigidGraded (VirSimple.virasoroSimple c) := by
\end{leancode}
(\lean{VOA.VirDeformModes.virWord_pairing_obstruction}, \leanpath{VOA/\allowbreak Examples/\allowbreak VirasoroWordPairing.lean};
\lean{VOA.VirDeformModes.NormalOrdering.simple_normalized_word_tangent},
\leanpath{VOA/\allowbreak Examples/\allowbreak VirasoroUniversalPairing.lean};
\lean{VOA.VirDeformModes.NormalOrdering.graded_rigidity_of_relation},
\leanpath{VOA/\allowbreak Examples/\allowbreak VirasoroUniformObstruction.lean}, where \lean{{ι : Type*} [Fintype ι]} is a
section variable.)

\begin{proof}[Proof on paper]
Let $D$ be graded with scalars $\alpha,\beta$ and normalised modes $\widetilde L_n$ (central operator
$c+\delta\lsevenN$). Lift the relation: $\widetilde s := \sum_jr_j\widetilde L_{w_j}(\lsevenvac,0)$. Its first coordinate is
$\sum_jr_jL_{w_j}\lsevenvac = s = 0$, so $\widetilde s = (0,s')$ for some $s'\in\lsevenV$. Apply the lifted adjoint
combination $\widetilde A := \sum_ir_i\widetilde L_{w_i^\dagger}$ and pair the tangent coordinate with the vacuum using the
\emph{undeformed} invariant form $B$ of $L(c,0)$:
\begin{itemize}
\item \emph{It vanishes.} On pure tangent states the normalised modes act by the undeformed $L_n$
  (\lean{normalizedL_tangent}), so $\widetilde A(0,s') = (0,As')$ with $A = \sum_ir_iL_{w_i^\dagger}$, and by
  invariance $B(\lsevenvac,As') = B(\sum_ir_iL_{w_i}\lsevenvac,s') = B(s,s') = B(0,s') = 0$. This is
  \lean{virWord_pairing_obstruction}, valid for any VOA, any $\beta$ and any invariant form.
\item \emph{It equals $P'_{w,r}(c)\,\delta$.} Expanding,
  $\widetilde A\widetilde s = \sum_{i,j}r_ir_j\widetilde L_{w_i^\dagger w_j}(\lsevenvac,0)$. In any representation of the Virasoro
  relations with central operator $\lsevenK$ and a vector killed by $L_n$, $n\ge-1$, normal ordering gives
  $L_w(\lsevenvac,0) = P_w(\lsevenK)(\lsevenvac,0) + (\text{creation words of positive weight applied to }(\lsevenvac,0))$; the
  vacuum pairing kills the positive-weight terms (weight orthogonality, Proposition~\ref{l04:thm:formbasics}),
  and by \eqref{l07:eq:taylor} with $\lsevenK = c+\delta\lsevenN$ the tangent coordinate of $P_w(\lsevenK)(\lsevenvac,0)$ is
  $P_w'(c)\delta\lsevenvac$. Summing, the tangent pairing is $P'_{w,r}(c)\delta$ (with $B(\lsevenvac,\lsevenvac) = 1$). This is
  \lean{simple_normalized_word_tangent} and its sum \lean{relation_tangent}.
\end{itemize}
Hence $P'_{w,r}(c)\delta = 0$, so $\delta = 0$ (\lean{central_variation_zero_of_relation}), and
Theorem~\ref{l07:thm:zerodelta} makes $D$ trivial.
\end{proof}

\begin{dictionary}
\textbf{Hypotheses of \lean{graded_rigidity_of_relation}.} \lean{hr}: the relation $s = 0$ holds in the
\emph{simple} algebra \lean{virasoroSimple c}, with its own modes \lean{(virasoroSimple c).L} --- an ordinary
statement about $L(c,0)$, not about deformations. \lean{hd}: the derivative \lean{derivative.eval c} of the
norm polynomial is nonzero at $c$. Neither mentions a deformation; the theorem has no deformation-theoretic
hypothesis at all. In particular it does \emph{not} assume determinant preservation, a deformed invariant
form, or a deformed singular vector --- the three ways in which earlier versions of the argument were
circular. \textbf{Conclusion:} \lean{IsRigidGraded} of Definition~\ref{l07:def:rigid}.

\textbf{Why the derivative, and why at $c$.} By \lean{relationPolynomial_eval}, $P_{w,r}(c) = \langle s,s\rangle$ in
$V(c,0)$; since $s$ maps to $0$ in $L(c,0)$, it lies in the radical and $P_{w,r}(c) = 0$ automatically. So the
content of \lean{hd} is that this zero is \emph{simple}. A double zero would give no information
(Warning~\ref{l07:warn:doubleroot}).

\textbf{How it could be vacuous.} If no relation with $P'\neq0$ existed at any $c$, the theorem would be empty.
It is not: at $c = -22/5$ the Lee--Yang relation of Example~\ref{l07:ex:leeyang} has $P' = -55$; at Ising,
$P' = 84084$ (Section~\ref{l07:sec:worked}); and at every $\lsevencpq$ the vacuum Gram determinant supplies one
(Section~\ref{l07:sec:all}). At generic $c$, $L(c,0) = V(c,0)$ has no relations at all, so the theorem says
nothing there --- consistent with $V(c,0)$ not being rigid.
\end{dictionary}

\begin{status}
\lean{virWord_pairing_obstruction}, \lean{simple_normalized_word_tangent}, \lean{graded_rigidity_of_relation}:
proved, axioms \lean{[propext, Classical.choice, Quot.sound]}. The theorem is unconditional; its two
hypotheses are ordinary algebraic facts to be checked at the charge of interest. The same file also gives
\lean{graded_rigidity_of_word_gram}, which replaces \lean{hr} by ``$r$ is a row-kernel vector of the
polynomial Gram matrix of a word basis evaluated at $c$'', so that the hypothesis can be checked entirely in
the universal module.
\end{status}

\begin{physics}
This is the null-vector argument for rigidity. In the minimal model the null state $s$ is set to zero. Perturb
the theory; to first order the only effect (after field redefinitions) is $c\to c+t\delta$. The relation $s=0$ must
persist, but the norm of $s$ in the perturbed theory is $P(c+t\delta)=tP'(c)\delta$, and a vanishing state has
zero norm. So $\delta=0$ unless the norm vanishes to second order. Physically: the minimal-model central
charge is pinned by null-vector decoupling, and to first order it cannot move.
\end{physics}

\subsection{The determinant form of the obstruction}

For a uniform theorem it is more convenient to use the whole vacuum Gram matrix. Index the partitions of
$n$ into parts $\ge2$ by \lean{SortedOfEnergy n} and let $G_n(X)$ be the matrix of norm polynomials
$P_{w_i^\dagger w_j}(X)$ of the corresponding creation words; $\det G_n(c) = D_n(c)$ is the vacuum Gram
determinant of Section~\ref{l07:sec:kac}.
\begin{leancode}
def partitionWord (n : ℕ) (w : SortedOfEnergy n) : List ℤ :=
  w.val.map (fun (i : ℕ) => -(i:ℤ)-2)

noncomputable def partitionGramPolynomial (n : ℕ) :
    Matrix (SortedOfEnergy n) (SortedOfEnergy n) (Polynomial ℂ) :=
  fun i j => wordPolynomial ((partitionWord n i).reverse.map Neg.neg ++ partitionWord n j)

theorem partitionGramPolynomial_eval (c : ℂ) (n : ℕ) :
    (partitionGramPolynomial n).map (Polynomial.eval c) =
      GramMat.gramMatrix (VirGram4.canB c) (n:ℤ) (gradedBasis c n) := by
\end{leancode}
(\lean{VOA.VirDeformModes.NormalOrdering.*}, \leanpath{VOA/\allowbreak Examples/\allowbreak VirasoroPartitionObstruction.lean};
\lean{GramMat.gramMatrix B n e} is the matrix $(B(e_i,e_j))$ of a basis $e$ of $\lsevenV_n$.)

\begin{theorem}[Simple determinant root obstruction]\label{l07:thm:det}
If $D_n(c) = 0$ and $D_n'(c)\ne0$ for some $n$, then $L(c,0)$ is graded rigid.
\end{theorem}
\begin{leancode}
theorem actual_partition_determinant_zero (c : ℂ)
    (D : FirstOrderDeformation (virasoroSimple c)) (hD : D.IsGraded)
    (α β : ℂ)
    (ha : D.smul₁ (virasoroSimple c).conformalVector 3
      (virasoroSimple c).conformalVector = α • (virasoroSimple c).vacuum)
    (hb : D.smul₁ (virasoroSimple c).conformalVector 1
      (virasoroSimple c).conformalVector = β • (virasoroSimple c).conformalVector)
    (n : ℕ) (hroot : (partitionGramPolynomial n).det.eval c = 0) :
    (partitionGramPolynomial n).det.aeval (inl c + inr (2*α-β*c) : DualNumber ℂ) = 0 := by

theorem graded_rigidity_of_simple_determinant (c : ℂ) (n : ℕ)
    (hroot : (partitionGramPolynomial n).det.eval c = 0)
    (hsimple : (partitionGramPolynomial n).det.derivative.eval c ≠ 0) :
    IsRigidGraded (virasoroSimple c) := by
\end{leancode}
(\lean{VOA.VirDeformModes.NormalOrdering.*}, \leanpath{VOA/\allowbreak Examples/\allowbreak VirasoroDeterminantObstruction.lean};
\lean{inl c + inr δ} is the dual number $c+\delta\varepsilon$ and \lean{aeval} evaluates a polynomial at it.)

\begin{proof}[Proof on paper]
Let $D$ be graded. \emph{Step 1: the determinant does not move.} Since $D_n(c) = 0$ there is a nonzero
$r$ with $\sum_ir_i\langle e_i,e_j\rangle = 0$ for all $j$, so $\sum_ir_ie_i$ lies in the radical and
$\sum_ir_i\pi(e_i) = 0$ in $L(c,0)$. Lift: $\sum_ir_i\widetilde L_{w_i}(\lsevenvac,0) = (0,x)$ with $x\in L(c,0)_n$. The
ordinary words $\pi(e_j)$ span $L(c,0)_n$, so $x = \sum_js_j\pi(e_j)$, and since $t\cdot(y,y_1) = (0,y)$ the
dual-number vector $\rho = r - ts$ satisfies $\sum_i\rho_i\widetilde L_{w_i}(\lsevenvac,0) = 0$. Pairing with the lifted
adjoint rows as in the previous proof shows $\rho^{T}G_n(c+t\delta) = 0$ over the dual numbers. As $r\neq0$, some
$\rho_i$ is a unit, so the matrix $G_n(c+t\delta)$ has a kernel vector with a unit entry and its determinant is
zero. \emph{Step 2: Taylor.} $0 = D_n(c+t\delta) = D_n(c) + tD_n'(c)\delta = tD_n'(c)\delta$, so $\delta = 0$
(\lean{VirRigidTangent.tangent_eq_zero_of_simple_root}). \emph{Step 3:} Theorem~\ref{l07:thm:zerodelta}.
\end{proof}

\begin{dictionary}
\lean{hroot} and \lean{hsimple} are exactly ``$D_n$ has a simple zero at $c$''; nothing about deformations is
assumed. Determinant preservation ($D_n(c+t\delta) = 0$) is a \emph{conclusion}
(\lean{actual_partition_determinant_zero}), proved from the actual deformation, not a hypothesis. This is the
``determinant'' twin of Theorem~\ref{l07:thm:relation}, and it is the form used for the general theorem.
\end{dictionary}

\begin{status}
Both theorems proved, axioms \lean{[propext, Classical.choice, Quot.sound]}. Unconditional.
\end{status}

\begin{warning}\label{l07:warn:doubleroot}
\textbf{An isolated degenerate charge is not enough; the root must be simple.} One might think that
rigidity follows because $\lsevencpq$ is an isolated zero of the level-$n$ determinant. It does not: over the
dual numbers, $X^2$ has the nonzero root $\varepsilon$, although $0$ is its only complex root. The library records
this as a theorem,
\begin{leancode}
theorem isolated_root_has_nonzero_tangent :
    (∀ c : ℂ, c ^ 2 = 0 ↔ c = 0) ∧
      (DualNumber.eps : DualNumber ℂ) ≠ 0 ∧
      (X ^ 2 : Polynomial ℂ).aeval (DualNumber.eps : DualNumber ℂ) = 0 := by
\end{leancode}
(\lean{VOA.VirRigidTangent.isolated_root_has_nonzero_tangent}, \leanpath{VOA/\allowbreak Examples/\allowbreak VirasoroRigidityTangent.lean}),
and the theorem that is used is the simple-root test
\lean{tangent_eq_zero_of_simple_root}: if $p'(c)\ne0$ and $p(c+\delta\varepsilon) = 0$ then $\delta = 0$. Concretely,
$D_4(c) = c^2(5c+22)/2$ has a \emph{double} zero at $c = 0$, so level $4$ says nothing at $c = 0$ (the useful
level there is $2$). A second, related trap is also recorded: ``the degenerate locus is discrete'' is false as
a statement over all levels, since the charges $\lsevencpq$ accumulate at $c = 1$; the correct and sufficient
statement is per level (for fixed $n$, $D_n$ is a nonzero polynomial with finitely many roots), and the
argument above uses only one level at a time.
\end{warning}
\begin{physics}
The determinant $D_n(c)$ is the Kac determinant of the vacuum module at level $n$; its zeros list the central
charges where a null descendant of the identity appears. A simple zero means that, as $c$ moves off the
minimal value, exactly one state stops being null at first order. The theorem says the minimal model cannot
be perturbed: any perturbation would have to move $c$, and moving $c$ would turn the decoupled null state back
into a physical state of nonzero norm.
\end{physics}


%======================================================================
\section{Worked cases}\label{l07:sec:worked}
%======================================================================

Before the general theorem, the library proved graded rigidity case by case, each time by exhibiting a
relation with a nonzero norm derivative and applying the mechanism of Theorem~\ref{l07:thm:relation} (in
the earliest cases, a hand-specialised version of it). These cases are worth reading: each is a complete,
checkable instance of the argument, with explicit numbers, and together they show the mechanism is not
vacuous. All of them are now also consequences of the general theorem of Section~\ref{l07:sec:all}.

\subsection{Central charge zero: \texorpdfstring{$M(2,3)$}{M(2,3)}}\label{l07:sec:zero}

At $c = 0$ the conformal vector has norm $\langle\omega,\omega\rangle = c/2 = 0$ and spans $V(0,0)_2$, so it lies in
the radical, i.e.\ in the maximal ideal: $\omega = 0$ in $L(0,0)$. Then $\alpha\lsevenvac = \omega^1_{(3)}\omega = 0$, so
$\alpha = 0$ and $\delta = 2\alpha - \beta\cdot0 = 0$; Theorem~\ref{l07:thm:zerodelta} applies.
\begin{leancode}
theorem zeroCharge_conformal_eq_zero : (virasoroSimple 0).conformalVector = 0 := by

theorem zeroCharge_graded_rigidity : Deform.IsRigidGraded (virasoroSimple 0) := by
\end{leancode}
(\lean{VOA.VirDeformModes.*}, \leanpath{VOA/\allowbreak Examples/\allowbreak VirasoroZeroChargeRigidity.lean}.)
\begin{dictionary}
In fact $L(0,0) = \C\lsevenvac$ is the one-dimensional VOA (every positive-weight creation monomial contains a
mode applied to $\omega = 0$, or more simply, $\omega = 0$ generates). Its graded rigidity is a triviality, and
the library's proof is correspondingly short. It is included honestly: $c_{2,3} = 0$, so the minimal model
$M(2,3)$ is the trivial theory, and the general theorem covers it with no exception (the vacuum Gram
determinant $D_2(c) = c/2$ has a simple zero at $0$).
\end{dictionary}


\begin{physics}
$M(2,3)$ is the trivial CFT with one state; its rigidity is the statement that a one-point theory cannot be
deformed. It is included because the minimal-model list starts there, and because $c=0$ is a point where the
level-two null vector is $T$ itself: at $c=0$ the stress tensor of the vacuum sector has zero norm and decouples.
\end{physics}

\subsection{Lee--Yang: \texorpdfstring{$M(2,5)$, $c = -22/5$}{M(2,5), c = -22/5}}\label{l07:sec:leeyang}

The relation is $s = 5L_{-2}^2\lsevenvac - 3L_{-4}\lsevenvac = 0$ in $L(-\frac{22}5,0)$
(Example~\ref{l07:ex:leeyang}). In any representation with central operator $\lsevenK$ and a vacuum-like vector
$v$, normal ordering gives the adjoint-word identity
\begin{equation*}
 (5L_2L_2 - 3L_4)(5L_{-2}L_{-2} - 3L_{-4})v = \big(\tfrac{25}{2}\lsevenK^2 + 55\lsevenK\big)v,
\end{equation*}
\lean{vir_singular4_polynomial}. With $\lsevenK = -\frac{22}5 + \delta\lsevenN$ the right side applied to
$(\lsevenvac,0)$ is $(0,-55\delta\lsevenvac)$ (\lean{leeYang_singular_central_polynomial}; the constant term
vanishes because $-22/5$ is a root and the tangent term is $P'(-22/5) = -55$). The vanishing pairing
forces $-55\delta = 0$.
\begin{leancode}
theorem leeYang_singular_central_polynomial (δ : ℂ) (v : V) :
    let K : End ℂ (V × V) := LeeYang.leeYangC • 1 + δ • tangentShift
    (25/2 : ℂ) • K (K (v, 0)) + (55 : ℂ) • K (v, 0) =
      (0, (-55 * δ) • v) := by

theorem leeYang_graded_rigidity : Deform.IsRigidGraded (virasoroSimple LeeYang.leeYangC) := by
\end{leancode}
(\lean{VOA.VirDeformModes.leeYang_singular_central_polynomial},
\leanpath{VOA/\allowbreak Examples/\allowbreak VirasoroDeformationGram4.lean};
\lean{VOA.VirDeformModes.leeYang_graded_rigidity}, \leanpath{VOA/\allowbreak Examples/\allowbreak VirasoroFixedReconstruction.lean};
\lean{LeeYang.leeYangC} is \lean{-22 / 5 : ℂ}.)


\begin{physics}
The Lee--Yang model describes the edge singularity of the Ising model in an imaginary magnetic field; it is the
simplest non-unitary CFT, with one nontrivial primary of weight $-1/5$. Its vacuum null vector at level four,
$(L_{-2}^2-\frac35L_{-4})\lsevenvac$, is the relation $\,{:}TT{:}\propto\partial^2T$ among chiral fields; the number $-55$
measures how fast this relation fails when $c$ moves off $-22/5$, and its nonvanishing pins $c$.
\end{physics}

\subsection{Ising and \texorpdfstring{$M(2,7)$}{M(2,7)} at weight six}

$c_{3,4} = \frac12$ and $c_{2,7} = -\frac{68}7$ both have first singular level $(p-1)(q-1) = 6$. The weight-six
space has the basis $L_{-6}\lsevenvac$, $L_{-4}L_{-2}\lsevenvac$, $L_{-3}^2\lsevenvac$, $L_{-2}^3\lsevenvac$, and the null relations
have integral coefficients:
\begin{equation*}
 \text{Ising: }(-108,-264,93,64),\qquad M(2,7):\ (-38,-77,-7,49).
\end{equation*}
They are certified as row-kernel vectors of the canonical Gram matrix \lean{canGram6} at the respective
charge (\lean{ising_gram6_row}, \lean{m27_gram6_row}) and hence vanish in $L(c,0)$. The norm polynomials,
computed by normal ordering with the central operator kept,
\begin{equation*}
 P_{\mathrm{Is}}(\lsevenK) = -83316\lsevenK + 165096\lsevenK^2 + 3072\lsevenK^3,\qquad
 P_{2,7}(\lsevenK) = 67830\lsevenK + \tfrac{48951}{2}\lsevenK^2 + \tfrac{7203}{4}\lsevenK^3,
\end{equation*}
vanish at their charges and have derivatives $84084$ and $102102$ there. (Check for Ising:
$-83316\cdot\frac12+165096\cdot\frac14+3072\cdot\frac18 = -41658+41274+384 = 0$ and
$-83316+165096+3\cdot3072\cdot\frac14 = 84084$.)
\begin{leancode}
def weight6Words : Fin 4 → List ℤ := ![[-6],[-4,-2],[-3,-3],[-2,-2,-2]]

def isingCoefficients : Fin 4 → ℂ := ![-108,-264,93,64]

def m27Coefficients : Fin 4 → ℂ := ![-38,-77,-7,49]

theorem ising_weight6_central {V : Type*} [AddCommGroup V] [Module ℂ V]
    (δ : ℂ) (v : V) :
    let K : End ℂ (V × V) := isingC • 1 + δ • tangentShift
    (-83316 : ℂ) • K (v,0) + (165096 : ℂ) • (K^2) (v,0) +
      (3072 : ℂ) • (K^3) (v,0) = (0, (84084*δ) • v) := by

theorem ising_graded_rigidity : Deform.IsRigidGraded (virasoroSimple isingC) := by

theorem m27_graded_rigidity : Deform.IsRigidGraded (virasoroSimple m27C) := by
\end{leancode}
(\lean{VOA.VirDeformModes.*}, \leanpath{VOA/\allowbreak Examples/\allowbreak VirasoroWeight6Norm.lean} and
\leanpath{VOA/\allowbreak Examples/\allowbreak VirasoroWeight6Rigidity.lean}; \lean{isingC} is \lean{1 / 2},
\lean{m27C} is \lean{-68 / 7}, both in \lean{VirGram56}.)

\begin{dictionary}
The normal forms were \emph{proposed} by a Python script
and each reduction step was then \emph{proved} in Lean from the Virasoro commutators, with the central
element kept as an operator so that the tangent coefficient is retained. Python is not trusted: it only
suggests the identities that Lean checks. The source states this, and it matters for the trust model: no
numerical evidence enters these theorems.
\end{dictionary}


\begin{physics}
The Ising CFT ($c=1/2$, primaries $1,\sigma,\varepsilon$ of weights $0,\frac1{16},\frac12$) is the continuum limit of the
critical two-dimensional Ising model, equivalently a free Majorana fermion. Heuristically, its vacuum relation at weight six
reflects that, for a single real fermion $\psi$, products of $T=-\frac12\psi\partial\psi$ satisfy an extra identity (from
$\psi^2=0$); $M(2,7)$ is a non-unitary model with the same first singular level. The nonzero derivatives
$84084$ and $102102$ say these relations cannot survive any change of $c$.
\end{physics}

\subsection{Tricritical Ising: \texorpdfstring{$M(4,5)$, $c = 7/10$}{M(4,5), c = 7/10}}

Here the first singular level is $3\cdot4 = 12$, where $V(c,0)_{12}$ has dimension $21$. The library certifies an
explicit weight-twelve combination of the $21$ partition words with integer coefficients, proves it is
annihilated by $L_1$ and $L_2$ at $c = 7/10$ (hence lies in the maximal ideal and vanishes in $L(\frac7{10},0)$),
proves it is nonzero at every charge (its ``vacuum survival'' value is $8000$), and computes its norm
polynomial, a polynomial of degree six whose derivative at $7/10$ is $128451135137139360$:
\begin{leancode}
noncomputable def tricriticalC : ℂ := 7/10

theorem tricritical_weight12_central {V : Type*} [AddCommGroup V] [Module ℂ V]
    (δ : ℂ) (v : V) :
    let K : End ℂ (V × V) := tricriticalC • 1+δ • tangentShift
    (321372476328566880:ℂ) • K (v,0) +
      (-1083581144832396000:ℂ) • (K^2) (v,0) +
      (866474649917232000:ℂ) • (K^3) (v,0) +
      (36252776782320000:ℂ) • (K^4) (v,0) +
      (527065804800000:ℂ) • (K^5) (v,0) +
      (2949120000000:ℂ) • (K^6) (v,0) = (0,(128451135137139360*δ) • v) := by

theorem tricritical_graded_rigidity : Deform.IsRigidGraded (virasoroSimple tricriticalC) := by

theorem minimal_four_five_graded_rigidity :
    Deform.IsRigidGraded (virasoroSimple (minimalCentralCharge 4 5)) := by
\end{leancode}
(\lean{VOA.VirDeformModes.*}, \leanpath{VOA/\allowbreak Examples/\allowbreak VirasoroTricriticalRigidity.lean}; the certificate,
with 555 private ordered-word reductions, is in \leanpath{VOA/\allowbreak Examples/\allowbreak VirasoroTricriticalNorm.lean}.)


\begin{physics}
The tricritical Ising model ($c=7/10$) describes the tricritical point of the Blume--Capel model and is the
simplest superconformal minimal model; its first vacuum null vector sits at weight twelve, which is why the
explicit certificate is large. Rigidity says the tricritical point is isolated among chiral algebras generated
by $T$.
\end{physics}

\subsection{The families \texorpdfstring{$M(2,q)$}{M(2,q)} and \texorpdfstring{$M(p,3)$}{M(p,3)}}

Two infinite families were done uniformly before the general theorem, each by constructing the relevant null
vector by an explicit recursion and proving its norm derivative nonzero.

\emph{$M(2,q)$, $q$ odd.} In Kac's parametrisation $c(t) = 13-6t-6/t$, the choice $t = (N+1)/2$ makes the
$(1,N)$ Kac weight $h_{1,N}(t) = \frac{(N-t)^2-(t-1)^2}{4t}$ equal to $0$, and $c((N+1)/2) = c_{2,N+1}$
(\lean{bsaCharge_eq_minimal}). The vacuum module therefore contains the $(1,N)$ singular vector of
Benoit and Saint-Aubin, built by a finite recursion in the $L_{-1},L_{-2}$ (the library constructs it from the
Ward identities for $L_1$, $L_2$, with only the denominators $r(N-r)$, without assuming a singular-vector
formula). For even $N$ the norm derivative is identified with a ``diagonal response'' computed by a
three-term recurrence and a contiguous relation, and shown nonzero.
\begin{leancode}
noncomputable def bsaCharge (N : ℕ) : ℂ :=
  13 - 3*((N:ℂ)+1) - 12/((N:ℂ)+1)

theorem bsa_even_graded_rigidity (k : ℕ) :
    Deform.IsRigidGraded (VirSimple.virasoroSimple (bsaCharge (2*(k+2)))) :=

theorem minimal_two_graded_rigidity (q : ℕ) (hq : 2≤q) (hcop : Nat.Coprime 2 q) :
    Deform.IsRigidGraded (VirSimple.virasoroSimple (minimalCentralCharge 2 q)) := by

theorem minimal_right_two_graded_rigidity (p : ℕ) (hp : 2≤p) (hcop : Nat.Coprime p 2) :
    Deform.IsRigidGraded (VirSimple.virasoroSimple (minimalCentralCharge p 2)) := by
\end{leancode}
(\lean{VOA.VirModule.bsaCharge}, \leanpath{VOA/\allowbreak Examples/\allowbreak VirasoroBSARecursion.lean}; the three theorems,
\leanpath{VOA/\allowbreak Examples/\allowbreak VirasoroBoundaryRigidity.lean}. \lean{bsa_even_graded_rigidity k} covers $q = 2k+5$;
$q = 3$ is the zero-charge case; coprimality with $2$ supplies oddness.)

\emph{$M(p,3)$.} At $t = 3/p$ the $(p-1,2)$ Kac weight is $0$ and the first vacuum singular level is
$2(p-1)$. The library constructs the singular vector by a ``fusion'' recursion (a triangular system dividing
only by charge-independent nonzero diagonals, nonresonant by coprimality with $3$), transports the
construction to a moving charge, proves the moving-norm identity
$\langle S(t),S(t)\rangle = -d_R(t)\,C(t)$ with $d_R$ vanishing simply at $t_0 = 3/p$ and $C(t_0)\ne0$ (the latter by a finite
hypergeometric-resonance argument), and concludes that the norm derivative is nonzero.
\begin{leancode}
theorem minimal_three_graded_rigidity (p : ℕ) (hp : 2≤p) (hcop : Nat.Coprime p 3) :
    Deform.IsRigidGraded (VirSimple.virasoroSimple (minimalCentralCharge p 3)) := by

theorem minimal_left_three_graded_rigidity (q : ℕ) (hq : 2≤q) (hcop : Nat.Coprime 3 q) :
    Deform.IsRigidGraded (VirSimple.virasoroSimple (minimalCentralCharge 3 q)) := by
\end{leancode}
(\lean{VOA.InverseDS.VirasoroVerma.*}, \leanpath{VOA/\allowbreak Examples/\allowbreak VirasoroMinimalThreeRigidity.lean}; the
construction occupies the \lean{VirasoroFusion*}, \lean{VirasoroMinimalThree*}, \lean{VirasoroThreeMoving*}
and \lean{VirasoroMovingNormDerivative} files of this chapter.)

\begin{status}
Every theorem of this section is proved; \texttt{\#print axioms} gives
\lean{[propext, Classical.choice, Quot.sound]} for \lean{zeroCharge_graded_rigidity},
\lean{leeYang_graded_rigidity}, \lean{ising_graded_rigidity}, \lean{m27_graded_rigidity},
\lean{tricritical_graded_rigidity}, \lean{minimal_four_five_graded_rigidity}, \lean{bsa_even_graded_rigidity},
\lean{minimal_two_graded_rigidity}, \lean{minimal_right_two_graded_rigidity},
\lean{minimal_three_graded_rigidity} and \lean{minimal_left_three_graded_rigidity}
(Section~\ref{l07:sec:axioms}). No \lean{VOA.pbw}. All unconditional statements of graded first-order
rigidity. Full names: the tricritical theorems are \lean{VOA.VirDeformModes.tricritical_graded_rigidity}
and \lean{VOA.VirDeformModes.minimal_four_five_graded_rigidity}; the three-label theorems are
\lean{VOA.InverseDS.VirasoroVerma.minimal_three_graded_rigidity} and
\lean{VOA.InverseDS.VirasoroVerma.minimal_left_three_graded_rigidity}; the two-label theorems are
\lean{VOA.VirModule.minimal_two_graded_rigidity} and \lean{VOA.VirModule.minimal_right_two_graded_rigidity}
(not all in \lean{VOA.VirModule}, despite the headline's namespace).
\end{status}

\begin{warning}\label{l07:warn:tautology}
\textbf{A ``structural'' rigidity theorem that was a tautology.} Before the cases above were proved, the file
\leanpath{VOA/\allowbreak Examples/\allowbreak VirasoroRigidityStructural.lean} stated rigidity of Ising, $M(2,7)$ and Lee--Yang from a
graded-dimension jump, under a hypothesis
\begin{leancode}
def RealisedByNondegenerate (U : VertexOperatorAlgebra 𝒰) (n : ℤ) (d : ℕ) : Prop :=
  (∃ D : FirstOrderDeformation U, D.IsGraded ∧ ¬ D.IsTrivial) → finrank ℂ (U.grading n) = d

theorem isRigidGraded_of_finrank_lt (U : VertexOperatorAlgebra 𝒰) (n : ℤ) (d : ℕ)
    (hE : RealisedByNondegenerate U n d) (hlt : finrank ℂ (U.grading n) < d) :
    IsRigidGraded U := by
\end{leancode}
intended to mean ``a modulus would move $U$ along the family $c\mapsto L(c,0)$, preserving graded dimensions''.
An adversarial audit showed that, given the (proved) jump \lean{hlt}, the hypothesis \lean{hE} is
\emph{logically equivalent to the conclusion}: \lean{hE} unfolds to ``a nontrivial graded deformation exists
$\Rightarrow$ (a false equation)'', i.e.\ to \lean{IsRigidGraded U}. The library now records the certificate
\begin{leancode}
theorem realisedByNondegenerate_iff_isRigidGraded
    {V : Type*} [AddCommGroup V] [Module ℂ V] (W : VertexOperatorAlgebra V)
    (n : ℤ) (d : ℕ) (hlt : Module.finrank ℂ (W.grading n) < d) :
    VirRigidStruct.RealisedByNondegenerate W n d ↔ Deform.IsRigidGraded W := by
\end{leancode}
(\lean{VOA.VirRigidTangent.realisedByNondegenerate_iff_isRigidGraded},
\leanpath{VOA/\allowbreak Examples/\allowbreak VirasoroRigidityTangent.lean}). So the instances
\lean{isRigidGraded_virasoroSimple_ising} etc.\ in that file read ``if $L(c,0)$ is graded rigid then it is graded
rigid''. They remain in the library (with honest docstrings) but are \emph{not} the rigidity theorems; the
unconditional theorems of this section are. The lesson, which the project applies throughout: a hypothesis
must be checked for \emph{uniform derivability from the conclusion}, not merely for truth. (The unconditional
content of that file --- the gauge line and the dimension jump
$\dim L(\frac12,0)_6<\dim L(1,0)_6$ --- is correct.)
\end{warning}

\begin{physics}
$M(2,q)$ are the non-unitary ``Lee--Yang series'' (they include Lee--Yang at $q=5$); their vacuum null vector
is of Benoit--Saint-Aubin type, the explicit $(1,s)$ null vectors found in the study of degenerate fields, whose
decoupling gives order-$s$ differential equations for correlators. $M(p,3)$ include the Ising model
$M(4,3)$. In both families the physics is the same as in the single cases: the null vector pins the central
charge to first order.
\end{physics}


%======================================================================
\section{All minimal models}\label{l07:sec:all}
%======================================================================

\begin{theorem}[Graded first-order rigidity of the Virasoro minimal models]\label{l07:thm:main}
Let $p,q\ge2$ be coprime integers and $\lsevencpq = 1-6(p-q)^2/(pq)$. Every graded first-order deformation of
the simple Virasoro vertex operator algebra $L(\lsevencpq,0)$ is trivial: there is a grade-preserving linear map
$\varphi$ with $\varphi(\lsevenvac) = 0$ such that $u^1_{(n)}v = (\varphi u)_{(n)}v - \varphi(u_{(n)}v) + u_{(n)}\varphi v$ for all
$u,v,n$.
\end{theorem}
\begin{leancode}
theorem minimal_graded_rigidity (p q : ℕ) (hp : 2 ≤ p) (hq : 2 ≤ q)
    (hcop : p.Coprime q) :
    IsRigidGraded (virasoroSimple (minimalCentralCharge p q)) := by
  apply minimal_rigidity_of_infinite_kac_roots p q hp hq hcop
  intro r s hr hs _
  exact infinite_integer_screening_roots r s hr hs
\end{leancode}
(\lean{VOA.VirModule.minimal_graded_rigidity}, \leanpath{VOA/\allowbreak Examples/\allowbreak VirasoroMinimalRigidityComplete.lean};
the three-line proof is shown because it is informative: it is the composition of two named theorems.)

\begin{dictionary}
\textbf{Decoding the statement.}
\begin{itemize}
\item \lean{(p q : ℕ)}: natural numbers; \lean{hp}, \lean{hq}: $p,q\ge2$; \lean{hcop : p.Coprime q}:
  $\gcd(p,q) = 1$. These are exactly the standard labels of the minimal models; the order of $p,q$ is
  irrelevant (the formula is symmetric and both orders are allowed), and $p = q$ is excluded by coprimality.
  There is no unitarity assumption: non-unitary models such as Lee--Yang are included.
\item \lean{minimalCentralCharge p q} is Definition~\ref{l07:def:cpq}, a rational number, silently cast to
  $\C$.
\item \lean{virasoroSimple c} is $L(c,0)$, the quotient of the universal vacuum VOA $V(c,0)$ by its maximal
  proper ideal (Definition~\ref{l07:def:Lc}); it is simple, with nonzero vacuum, for every $c$.
\item \lean{IsRigidGraded} is Definition~\ref{l07:def:rigid}, examined in Section~\ref{l07:sec:fidelity-rigid}.
\end{itemize}

\textbf{Every hypothesis is the intended one, and there are no others.} The statement has no hypothesis about
singular vectors, determinants, characters, vacuum survival, norm derivatives, or deformations; earlier
versions of the theorem (\lean{minimal_graded_rigidity_of_kac_singular}, \lean{graded_rigidity_of_verma_kac},
\lean{minimal_rigidity_of_infinite_kac_roots}) carried such inputs as explicit hypotheses, and the final
theorem discharges them, as its displayed proof shows.

\textbf{Same as the intended theorem?} Yes, in the sense of Section~\ref{l07:sec:fidelity-rigid}: it is
\emph{graded} first-order rigidity of the vacuum VOA of every minimal model, over the dual numbers, modulo
graded vacuum-preserving coordinate changes, with the conformal vector allowed to move. It does not assert
rigidity among ungraded deformations (where the library's \lean{IsRigid} is false for a definitional reason,
Warning~\ref{l07:warn:isrigid}), higher-order or analytic rigidity, rigidity of the minimal-model modules or
of the full non-chiral CFT, or the rationality of $L(\lsevencpq,0)$ (which is Wang's theorem, a separate matter).

\textbf{Could it be vacuous?} The hypotheses are satisfiable ($p = 2$, $q = 3$, \dots); the object
$L(\lsevencpq,0)$ is a genuine nonzero simple VOA (for $(p,q)\ne(2,3),(3,2)$ it is infinite-dimensional, with
$\omega\neq0$); the deformation notion contains every genuine graded first-order deformation and admits nonzero
graded cocycles (e.g.\ $\delta L_0$); the conclusion is not a tautology of the definitions: it fails for
$V(c,0)$ at every $c$, machine-checked (\lean{VOA.}\allowbreak\lean{VirPoly.}\allowbreak\lean{not_}\allowbreak\lean{isRigidGraded_}\allowbreak\lean{virasoroVOA}, Remark~\ref{l07:rem:nonvacuity}).

\textbf{The label $(2,3)$.} At $(p,q) = (2,3)$ (or $(3,2)$) the central charge is $c_{2,3} = 0$, and
$L(0,0) = \C\lsevenvac$ is the one-dimensional VOA (Section~\ref{l07:sec:zero}); its graded rigidity is
trivially true, since $Y_1\lsevenvac = 0$ is forced by (D2). This instance of the theorem carries no content.
It is harmless: the other labels do.
\end{dictionary}

\begin{physics}
Every Virasoro minimal model --- Ising, tricritical Ising, Potts, Lee--Yang and all non-unitary $M(p,q)$ --- has a chiral
algebra with no exactly marginal deformations (to first order, among weight-preserving perturbations): the
minimal models are isolated points, not members of continuous families. This matches the physical picture
in which $c$ is quantised by null-vector decoupling and the Kac table, and the classification of minimal
models as a discrete list. Honest caveats: this is a statement about the chiral vacuum algebra only (not about
the full non-chiral CFT with its modular-invariant combination of sectors, nor about its modules), and about
first order only. That the full CFT is isolated too is widely believed, but it is a different statement
(about weight-$(1,1)$ fields of the non-chiral theory) and is not what is proved.
\end{physics}

\subsection{The proof on paper}

Fix coprime $p,q\ge2$, put $c_0 = \lsevencpq$ and $N = (p-1)(q-1)$, and let $D$ be a graded first-order deformation
of $L(c_0,0)$.
\begin{enumerate}
\item \emph{Normal form} (Section~\ref{l07:sec:anatomy}). $D$ determines scalars $\alpha,\beta$ and normalised modes
  $\widetilde L_n$ on $\lsevenV\oplus t\lsevenV$ satisfying the Virasoro relations with central operator $c_0+\delta\lsevenN$,
  $\delta = 2\alpha-\beta c_0$.
\item \emph{Reconstruction} (Theorem~\ref{l07:thm:zerodelta}). If $\delta = 0$ then $D$ is trivial.
\item \emph{Obstruction} (Theorem~\ref{l07:thm:det}). If the vacuum Gram determinant $D_N(c) = \det G_N(c)$ has a
  simple zero at $c_0$, then $\delta = 0$.
\item \emph{Representation theory} (Chapter~\ref{ch:L06}). $D_N(c_0) = 0$ and $D_N'(c_0)\ne0$.
\end{enumerate}
Steps 1--3 are this chapter's; step 4 is a statement purely about the universal Virasoro vacuum module, and is
the deep input. In the literature it follows from the Kac determinant formula for the vacuum module (Gorelik and
Kac): $D_N(c)$ is, up to a nonzero constant, $\prod(c-c_{a,b})^{\dim L(h;c_{a,b})_N}$ over primitive labels, and
at the first level $N = (p-1)(q-1)$ the factor $c-c_{p,q}$ occurs exactly once (the lowest weight space of the
irreducible quotient is one-dimensional) while every other factor is nonzero at $c_0$ (distinct coprime label
pairs with ratios $>1$ give distinct charges). The Lean proof does \emph{not} cite this; it proves step 4 from
scratch along the following route, all of which lives in files of Chapter~\ref{ch:L06}:
\begin{enumerate}[label=(\alph*)]
\item \emph{Kac parametrisation.} $c(t) = 13-6t-6/t$, $h_{r,s}(t) = \frac{(rt-s)^2-(t-1)^2}{4t}$; for $p>q$ one takes
  $t = p/q$, so that $c(p/q) = c_{p,q}$, and the label $(q-1,p-1)$ has $h = 0$ at level $(p-1)(q-1)$.
\item \emph{Gram roots from free fields.} For positive integers $r,s,m$, at $t = m^2$, an explicit finite
  ``screening'' polynomial in a rank-one charged Fock space is a nonzero state of energy $rs$ annihilated by
  $L_1$ and $L_2$; a transfer argument (equal level dimensions of Fock and Verma, injectivity below the first
  degenerate level) turns it into a zero of the Verma Gram determinant at $(c(t),h_{r,s}(t))$ and level $rs$.
  This gives, for each label, the infinite family of roots $t\in\{1,4,9,\dots\}$:
  \lean{infinite_integer_screening_roots}.
\item \emph{All parameters.} Along a Kac curve the determinant is a Laurent polynomial in $t$; infinitely many
  roots force it to vanish identically (\lean{kacGram_root_of_infinite}). Strong induction on the level then
  produces actual nonzero singular vectors at nonreal $t$, and degree exhaustion
  $\sum_{\lambda\vdash n}\ell(\lambda) = \sum_{rs\le n}p(n-rs)$ identifies the Verma determinant with the Kac product;
  continuity extends it to all $t\ne0$.
\item \emph{From Verma to vacuum.} The coefficient of the top power of $h$ in the level-$(n+1)$ Verma determinant
  equals $2^{p(n)}\det G^{\mathrm{vac}}_{n+1}(c)\det G^{\mathrm{Verma}}_n(c,1)$; comparing with the Kac product, whose
  relevant factor vanishes simply at $t = p/q$ (derivative $q(q-p)/(2p)\neq0$) while the weight-one factor does not
  vanish, the chain rule gives $D_N(c_0) = 0\neq D_N'(c_0)$ (\lean{vacuum_simple_of_verma_kac}).
\end{enumerate}
The two Lean theorems composed in the displayed proof are the end points of this route:
\begin{leancode}
theorem minimal_rigidity_of_infinite_kac_roots (p q : ℕ) (hp : 2 ≤ p)
    (hq : 2 ≤ q) (hcop : p.Coprime q)
    (hroots : ∀ r s : ℕ, 0 < r → 0 < s → r*s ≤ (p-1)*(q-1) →
      ∃ S : Set ℂ, S.Infinite ∧ (∀ t ∈ S, t ≠ 0) ∧
        ∀ t ∈ S, (vermaGramMatrix (vermaBSACharge t)
          (fusionKacWeight r s t) (r*s)).det = 0) :
    IsRigidGraded (virasoroSimple (minimalCentralCharge p q)) := by

theorem infinite_integer_screening_roots (r s : ℕ) (hr : 0 < r) (hs : 0 < s) :
    ∃ S : Set ℂ, S.Infinite ∧ (∀ t ∈ S, t ≠ 0) ∧
      ∀ t ∈ S, (vermaGramMatrix (vermaBSACharge t)
        (fusionKacWeight r s t) (r*s)).det = 0 := by
\end{leancode}
(\lean{VOA.InverseDS.VirasoroVerma.minimal_rigidity_of_infinite_kac_roots},
\leanpath{VOA/\allowbreak Examples/\allowbreak VirasoroKacRootBootstrap.lean};
\lean{VOA.VirasoroScreening.infinite_integer_screening_roots},
\leanpath{VOA/\allowbreak Examples/\allowbreak VirasoroScreeningKacRoots.lean}; both files belong to Chapter~\ref{ch:L06}.)

\begin{dictionary}
\lean{vermaBSACharge t} is $c(t) = 13-6t-6/t$ and \lean{fusionKacWeight r s t} is $h_{r,s}(t)$;
\lean{vermaGramMatrix c h n} is the level-$n$ Gram (Shapovalov) matrix of the Verma module $M(c,h)$ in the
partition basis. The hypothesis \lean{hroots} of the first theorem is a statement about Verma modules only ---
for each positive label $(r,s)$ with $rs\le(p-1)(q-1)$, infinitely many distinct nonzero parameters $t$ at which
$\det G^{\mathrm{Verma}}_{rs}(c(t),h_{r,s}(t)) = 0$ --- and the second theorem supplies it for \emph{all} positive
labels, with $S = \{1,4,9,\dots\}$. The source of the first theorem notes that isolated collision parameters could
not satisfy \lean{hroots}; it genuinely requires infinitely many roots, and gets them. The internal route
(a)--(d) replaces the external citation of Gorelik--Kac, which the project's paper note
(\texttt{Notes/Virasoro\_Uniform\_Rigidity\_Argument.tex}) keeps as an alternative exposition, not a proof
dependency.
\end{dictionary}

\begin{status}
\lean{VOA.VirModule.minimal_graded_rigidity}: \textbf{graded first-order rigidity, proved, unconditional}
(no hypotheses beyond the label conditions; the notion is non-vacuous, since $V(c,0)$ is proved not graded rigid,
Remark~\ref{l07:rem:nonvacuity}). \texttt{\#print axioms} (run on this
commit, see Section~\ref{l07:sec:axioms}):
\begin{plaincode}
'VOA.VirModule.minimal_graded_rigidity' depends on axioms: [propext, Classical.choice, Quot.sound]
\end{plaincode}
\lean{VOA.pbw} does \emph{not} appear: the result does not use \lean{VOA.pbw}. Its hypotheses are only
the label conditions. (The project's build record reports an aggregate \lean{lake build} of 9,874 jobs and
a full sequential axiom audit with the same outcome at the time the theorem landed; for this book we re-ran
\texttt{\#print axioms} on the built tree and did not re-run the aggregate build.) The repository's
\texttt{README.md} now records the theorem as proved, with the same graded first-order qualification.
\end{status}

%======================================================================
\section{Other files of this chapter}\label{l07:sec:others}
%======================================================================

For completeness, here is what the remaining assigned files contain; none of them is needed for the statements
above beyond what has been described.
\begin{itemize}
\item \leanpath{Core/\allowbreak DeformationCurrentModes}, \leanpath{Core/\allowbreak DeformationPoleForm},
  \leanpath{Core/\allowbreak DeformationNormalization}, and the current part of
  \leanpath{Core/\allowbreak DeformationLiftedModes}, \leanpath{Core/\allowbreak DeformationCentralModes}: the weight-one
  (affine) analogue of Section~\ref{l07:sec:anatomy} --- lifted current modes, the invariant pole pairing, the
  level variation $k+\delta\lsevenN$, and gauge normalisation via $H^2(\mathfrak g,\mathfrak g) = 0$. They serve the affine
  rigidity theorems of Chapter~\ref{ch:L09}. \leanpath{Core/\allowbreak DeformationIteration}: the iterate formula on
  $\lsevenV\times\lsevenV$ used in Section~\ref{l07:sec:reconstruction}.
\item \leanpath{Core/\allowbreak RankOneTransfer} (completing the square for a positive-definite Hermitian matrix minus a rank-one
  term), \leanpath{Core/\allowbreak GlobalConformalTowers} (norms of global descendants
  $\langle L_{-1}^ru,L_{-1}^rv\rangle = r!(2m)_r\langle u,v\rangle$ for quasi-primaries of weight $m$), \leanpath{Core/\allowbreak HorizontalSaturation}
  ($\mathfrak{sl}_2$ ladder saturation): linear-algebra tools used by the graded-unitarity work of Chapter~\ref{ch:L19}.
\item \leanpath{Examples/\allowbreak VirasoroDeformationLow}, \leanpath{VirasoroRigidityGramLift}: the tridegree tables of
  Section~\ref{l07:sec:tridegree} and the level-$4$/$6$ Gram matrices over $\C[X]$ evaluated at dual numbers.
  \leanpath{VirasoroCDeformation}: the gauge line and the conditional non-rigidity of $V(c,0)$.
  \leanpath{VirasoroPolyStructure}, \leanpath{VirasoroPolyFields}, \leanpath{VirasoroPolyTensor},
  \leanpath{VirasoroPolyEnergy}, \leanpath{VirasoroPolyDeriv}, \leanpath{VirasoroCDefBuild},
  \leanpath{VirasoroCDefJacobi}, \leanpath{VirasoroCDefFinal}: the structure constants of $V(c,0)$ as polynomials in $c$ on
  the sorted-word basis, their $c$-derivative $\partial_cY$ (\lean{cDeformation}), and the unconditional non-rigidity
  \lean{not_isRigidGraded_virasoroVOA} (Section~\ref{l07:sec:nonvacuity}). (These files postdate the assignment list of
  the first edition.)
  \leanpath{VirasoroDeformationReconstruction}: generation of $L(c,0)$ by $\omega$ and the generated-carrier form of
  reconstruction.
\item \leanpath{VirasoroBSA*}, \leanpath{VirasoroVermaBSA*}, \leanpath{VirasoroDescendantBSATransport}: the
  Benoit--Saint-Aubin $(1,N)$ singular vectors, in the vacuum module and in general Verma modules
  (\lean{vermaBSA_nonzero_highest_all}: a nonzero highest-weight vector at level $N$ of $M(c(t),h_{1,N}(t))$ for every
  $t\ne0$, $N>0$). \leanpath{VirasoroFusion*}, \leanpath{VirasoroVacuumFusion*}, \leanpath{VirasoroMinimalThree*},
  \leanpath{VirasoroThreeMoving*}, \leanpath{VirasoroMovingNormDerivative}, \leanpath{VirasoroMinimalFourVacuum},
  \leanpath{VirasoroMinimalFusionArithmetic}, \leanpath{VirasoroMinimalChargeArithmetic}: the fusion construction of
  singular vectors and its arithmetic (nonresonance at prime-power labels; uniqueness of the primitive label of a
  charge). \leanpath{VirasoroTranslationTangent}, \leanpath{VirasoroTranslationFirstForm}: the translated first-order
  form, giving the Verma-to-vacuum coefficient identity of step (d).
\item \leanpath{VirasoroSalazarInterface}: definitions of predicates expressing Salazar's admissible-monomial basis
  theorem for the boundary simple Virasoro module, with implications between them; nothing is assumed or asserted,
  and no axiom is added.
\end{itemize}
Several of these files contain conditional theorems whose hypotheses were later discharged (for instance
\lean{graded_rigidity_of_bsa_scalar}, \lean{minimal_graded_rigidity_of_kac_singular},
\lean{graded_rigidity_of_verma_kac}); their docstrings say ``this conditional theorem does not assert that the
scalar input holds''. They are correct as stated and are superseded by the unconditional theorems quoted in this
chapter.

%======================================================================
\section{Verification record}\label{l07:sec:axioms}
%======================================================================

When this chapter was first written we ran \texttt{\#print axioms} on the headline declarations, on the already-built tree at
an earlier commit (not re-run for this edition), with a scratch file importing the relevant modules and the command
\begin{plaincode}
export PATH="$HOME/.elan/bin:$PATH"
lake env lean -DmaxSynthPendingDepth=3 Axioms.lean
\end{plaincode}
(the flag avoids the audit trap of Section~\ref{ch15:sec:trap}: a bare \texttt{lake env lean} drops
\texttt{maxSynthPendingDepth}). Every one of the following reported exactly
\begin{plaincode}
depends on axioms: [propext, Classical.choice, Quot.sound]
\end{plaincode}
\begin{quote}\small
\lean{VirModule.minimal_graded_rigidity};
\lean{VirDeformModes.NormalOrdering.graded_rigidity_of_relation};
\lean{VirDeformModes.NormalOrdering.graded_rigidity_of_simple_determinant};
\lean{VirDeformModes.isTrivial_of_zero_central_variation};
\lean{VirDeformModes.virasoroSimple_exists_normalizedData}; \lean{VirDeformModes.normalizedL_comm};
\lean{VirDeformModes.virWord_pairing_obstruction}; \lean{Deform.Reconstruction.isTrivial_of_closed_lift};
\lean{Deform.coboundaries_le_cocycles}; \lean{Deform.isRigid_iff_H_trivial};
\lean{Deform.isRigidGraded_of_gradedCocycles_le};
\lean{Deform.gradedCocyclesAt_le_gradedCoboundariesAt_of_isRigid}; \lean{DefForm.lift_iff_pairing_zero};
\lean{RigidBridge.F2_lt_F1_of_self_pairing_ne_zero}; \lean{RigidBridge.shapForm_symm};
\lean{NonzeroDeform.not_isRigid_virasoroVOA}; \lean{NonzeroDeform.isGraded_of_isTrivial};
\lean{NonzeroDeform.shapForm₁_self_eq_zero_of_isTrivial}; \lean{NonzeroDeform.exists_nonzero_deformation};
\lean{VirCDef.not_isRigidGraded_of_isCDeformation}; \lean{VirCDef.cmode_conformal_of_isCoboundary};
\lean{VirDeformModes.zeroCharge_graded_rigidity}; \lean{VirDeformModes.leeYang_graded_rigidity};
\lean{VirDeformModes.ising_graded_rigidity}; \lean{VirDeformModes.m27_graded_rigidity};
\lean{VirDeformModes.tricritical_graded_rigidity}; \lean{VirDeformModes.minimal_four_five_graded_rigidity};
\lean{VirModule.minimal_two_graded_rigidity}; \lean{VirModule.minimal_right_two_graded_rigidity};
\lean{VirModule.bsa_even_graded_rigidity}; \lean{InverseDS.VirasoroVerma.minimal_three_graded_rigidity};
\lean{InverseDS.VirasoroVerma.minimal_left_three_graded_rigidity};
\lean{InverseDS.VirasoroVerma.minimal_rigidity_of_infinite_kac_roots};
\lean{InverseDS.VirasoroVerma.minimal_graded_rigidity_of_kac_singular};
\lean{VirasoroScreening.infinite_integer_screening_roots};
\lean{VirRigidTangent.tangent_eq_zero_of_simple_root}; \lean{VirRigidTangent.isolated_root_has_nonzero_tangent};
\lean{VirRigidStruct.isRigidGraded_of_finrank_lt};
\lean{VirRigidTangent.realisedByNondegenerate_iff_isRigidGraded}.
\end{quote}
(All names in the \lean{VOA} namespace.) The PBW theorem \lean{VOA.pbw} appears in none of them. We did not
re-run the aggregate \texttt{lake build} for this chapter; the tree had been built, and the \texttt{olean} files of
the imported modules were present. A clean \texttt{\#print axioms} does not report explicit hypotheses; here the
headline theorems have none beyond the label conditions, so ``unconditional'' is justified for them, while the
reduction theorems (\lean{graded_rigidity_of_relation}, \lean{graded_rigidity_of_simple_determinant},
\lean{minimal_rigidity_of_infinite_kac_roots}, \ldots) are unconditional \emph{implications} whose hypotheses are
visible in their statements.

\textbf{Re-run 2026-10-27.} At a later commit (not re-run for this edition), with the same command on a
scratch file, \texttt{\#print axioms} again gave exactly \lean{[propext, Classical.choice, Quot.sound]} for
\lean{VirModule.minimal_graded_rigidity}, \lean{InverseDS.VirasoroVerma.minimal_three_graded_rigidity},
\lean{InverseDS.VirasoroVerma.minimal_left_three_graded_rigidity}, \lean{VirDeformModes.tricritical_graded_rigidity},
\lean{VirDeformModes.minimal_four_five_graded_rigidity}, \lean{VirCDef.cmode_conformal_of_isCoboundary},
\lean{VirCDef.not_isTrivial_of_omega_witness}, \lean{VirCDef.not_isRigidGraded_of_omega_witness},
\lean{VirCDef.not_isTrivial_of_isCDeformation}, \lean{VirCDef.not_isRigidGraded_of_isCDeformation} and
\lean{NonzeroDeform.not_isRigid_virasoroVOA} (all in \lean{VOA}). No aggregate \texttt{lake build} was run.

Release audit 2026-10-27: \texttt{Notes/Textbook/referee/RELEASE\_AUDIT\_RIGIDITY.md}. Its verdict on the
headline and its special cases is ``announceable with caveats''; the two caveats (graded first-order only;
non-vacuity argued on paper only) are stated in the opening of the chapter, in
Section~\ref{l07:sec:fidelity-rigid} and in Remark~\ref{l07:rem:nonvacuity}. The second caveat has since been
discharged by \lean{VOA.}\allowbreak\lean{VirPoly.}\allowbreak\lean{not_}\allowbreak\lean{isRigidGraded_}\allowbreak\lean{virasoroVOA} (release follow-ups audit,
\texttt{Notes/Textbook/referee/RELEASE\_AUDIT\_FOLLOWUPS.md}, 2026-10-28: ``with caveats: holds for every $c$; the
deformation moves $c$''; it reports standard axioms for \lean{cDeformation}, \lean{isCDeformation_cDeformation},
\lean{cDeformation_isGraded} and \lean{not_isRigidGraded_virasoroVOA}). We did not re-run \texttt{\#print axioms} on these
for this chapter. The remaining caveat is that rigidity is graded first-order only (plus the remark on the label
$(2,3)$).

%======================================================================
\section{Referee's checklist}\label{l07:sec:checklist}
%======================================================================

\begingroup\small
\begin{longtable}{>{\raggedright\arraybackslash}p{0.29\textwidth}>{\raggedright\arraybackslash}p{0.33\textwidth}>{\raggedright\arraybackslash}p{0.2\textwidth}>{\raggedright\arraybackslash}p{0.1\textwidth}}
\toprule
Lean declaration & Intended meaning & Hypotheses / status & Check\\
\midrule\endhead
\lean{Deform.}\allowbreak\lean{FirstOrderDeformation} & $Y+tY_1$ is a VA over $\C[t]/t^2$, same space and vacuum & definition; (D1)--(D4), \lean{finsum} non-vacuous & Def.~\ref{l07:def:fod}\\
\lean{…FirstOrderDeformation.}\allowbreak\lean{IsGraded} & $Y_1$ respects weights (VA1 at order $t$) & definition (predicate) & Def.~\ref{l07:def:graded}\\
\lean{Deform.}\allowbreak\lean{IsCoboundary}, \lean{…IsTrivial} & $Y_1 = \delta\varphi$, $\varphi$ graded, $\varphi\lsevenvac = 0$ & definition & Def.~\ref{l07:def:trivial}\\
\lean{Deform.}\allowbreak\lean{coboundaries_}\allowbreak\lean{le_}\allowbreak\lean{cocycles} & field redefinitions are deformations & none; unconditional & Thm.~\ref{l07:thm:cobcoc}\\
\lean{Deform.IsRigid} & every deformation trivial & definition; \textbf{TRAP}: false for every $V(c,0)$ (definitional) & Warn.~\ref{l07:warn:isrigid}\\
\lean{Deform.}\allowbreak\lean{IsRigidGraded} & every graded first-order deformation trivial (graded $H^2 = 0$, vacuum fixed; $\ge$ Linshaw--Qi) & definition; the notion used & Def.~\ref{l07:def:rigid}, §\ref{l07:sec:fidelity-rigid}\\
\lean{Deform.}\allowbreak\lean{isRigid_}\allowbreak\lean{iff_}\allowbreak\lean{H_}\allowbreak\lean{trivial} & rigidity $\iff$ $H = 0$ & none; unconditional & §\ref{l07:sec:deformations}\\
\lean{Deform.}\allowbreak\lean{GradedToGlobalRigidity} & tridegree-wise $\Rightarrow$ global & \textbf{named target, not proved}, unused & Warn.~\ref{l07:warn:tridegree}\\
\lean{VirCDef.}\allowbreak\lean{not_}\allowbreak\lean{isRigidGraded_}\allowbreak\lean{of_}\allowbreak\lean{isCDeformation} & $\partial_c$ would refute rigidity of $V(c,0)$ & conditional; the predicate is inhabited by \lean{VirPoly.cDeformation} & Warn.~\ref{l07:warn:nonvacuity}\\
\lean{VirPoly.}\allowbreak\lean{not_}\allowbreak\lean{isRigidGraded_}\allowbreak\lean{virasoroVOA} & $V(c,0)$ is not graded first-order rigid, every $c$ & none; unconditional (the witness $\partial_cY$ moves $c$) & §\ref{l07:sec:nonvacuity}\\
\lean{VOA.}\allowbreak\lean{VirCDef.}\allowbreak\lean{not_}\allowbreak\lean{isTrivial_}\allowbreak\lean{of_}\allowbreak\lean{omega_}\allowbreak\lean{witness} & a deformation of $V(c,0)$ with $\omega^1_{(1)}\omega = 0\neq\omega^1_{(3)}\omega$ is non-trivial & the two $\omega$-values; unconditional & Rem.~\ref{l07:rem:nonvacuity}\\
\lean{VirCDef.}\allowbreak\lean{cmode_}\allowbreak\lean{conformal_}\allowbreak\lean{of_}\allowbreak\lean{isCoboundary} & gauge line: coboundaries have $\delta = 0$ & $\lsevenV_2 = \C\omega$; unconditional & Thm.~\ref{l07:thm:gaugeline}\\
\lean{VirDeformModes.}\allowbreak\lean{normalizedL_}\allowbreak\lean{comm} & normalised modes: Virasoro with centre $c+\delta\lsevenN$ & graded, $\lsevenV_1 = 0$; unconditional & §\ref{l07:sec:anatomy}\\
\lean{Deform.}\allowbreak\lean{Reconstruction.}\allowbreak\lean{isTrivial_}\allowbreak\lean{of_}\allowbreak\lean{closed_}\allowbreak\lean{lift} & closed separated lift $\Rightarrow$ trivial & $W$ simple + 5 lift properties & Thm.~\ref{l07:thm:closedlift}\\
\lean{VirDeformModes.}\allowbreak\lean{isTrivial_}\allowbreak\lean{of_}\allowbreak\lean{zero_}\allowbreak\lean{central_}\allowbreak\lean{variation} & $\delta = 0\Rightarrow$ trivial, any $c$ & graded; unconditional & Thm.~\ref{l07:thm:zerodelta}\\
\lean{DefForm.}\allowbreak\lean{lift_}\allowbreak\lean{iff_}\allowbreak\lean{pairing_}\allowbreak\lean{zero} & lifting criterion for null vectors & $B_0$ symmetric, fin.-dim.; unconditional & Prop.~\ref{l07:prop:lift}\\
\lean{RigidBridge.}\allowbreak\lean{F2_}\allowbreak\lean{lt_}\allowbreak\lean{F1_}\allowbreak\lean{of_}\allowbreak\lean{self_}\allowbreak\lean{pairing_}\allowbreak\lean{ne_}\allowbreak\lean{zero} & a null vector with $\langle s,s\rangle^1\ne0$ does not lift & hypothesis unsatisfiable for trivial $D$; not used & Warn.~\ref{l07:warn:bridge}\\
\lean{…NormalOrdering.}\allowbreak\lean{graded_}\allowbreak\lean{rigidity_}\allowbreak\lean{of_}\allowbreak\lean{relation} & relation with simple-zero norm $\Rightarrow$ rigid & relation in $L(c,0)$, $P'(c)\ne0$; unconditional & Thm.~\ref{l07:thm:relation}\\
\lean{…NormalOrdering.}\allowbreak\lean{graded_}\allowbreak\lean{rigidity_}\allowbreak\lean{of_}\allowbreak\lean{simple_}\allowbreak\lean{determinant} & simple zero of $D_n$ $\Rightarrow$ rigid & $D_n(c) = 0\neq D'_n(c)$; unconditional & Thm.~\ref{l07:thm:det}\\
\lean{VirRigidTangent.}\allowbreak\lean{isolated_}\allowbreak\lean{root_}\allowbreak\lean{has_}\allowbreak\lean{nonzero_}\allowbreak\lean{tangent} & isolated root $\not\Rightarrow$ rigid & none & Warn.~\ref{l07:warn:doubleroot}\\
\lean{VirRigidStruct.}\allowbreak\lean{isRigidGraded_}\allowbreak\lean{of_}\allowbreak\lean{finrank_}\allowbreak\lean{lt} & ``dimension jump'' rigidity & \textbf{TAUTOLOGY}: hypothesis $\iff$ conclusion & Warn.~\ref{l07:warn:tautology}\\
\lean{VOA.}\allowbreak\lean{VirDeformModes.}\allowbreak\lean{zeroCharge_}\allowbreak\lean{graded_}\allowbreak\lean{rigidity} & $M(2,3)$, $c = 0$, $L(0,0) = \C\lsevenvac$: trivially rigid & none; unconditional & §\ref{l07:sec:zero}\\
\lean{VirDeformModes.}\allowbreak\lean{leeYang_}\allowbreak\lean{graded_}\allowbreak\lean{rigidity} & Lee--Yang $c = -22/5$ & none; unconditional & §\ref{l07:sec:leeyang}\\
\lean{VirDeformModes.}\allowbreak\lean{ising_}\allowbreak\lean{graded_}\allowbreak\lean{rigidity}, \lean{m27_graded_rigidity} & Ising, $M(2,7)$ & none; unconditional & §\ref{l07:sec:worked}\\
\lean{VOA.}\allowbreak\lean{VirDeformModes.}\allowbreak\lean{tricritical_}\allowbreak\lean{graded_}\allowbreak\lean{rigidity} (+ \lean{minimal_four_five}) & $M(4,5)$ & none; unconditional & §\ref{l07:sec:worked}\\
\lean{VOA.}\allowbreak\lean{VirModule.}\allowbreak\lean{minimal_}\allowbreak\lean{two_}\allowbreak\lean{graded_}\allowbreak\lean{rigidity} (+ \lean{right}) & $M(2,q)$, $M(p,2)$ & labels; unconditional & §\ref{l07:sec:worked}\\
\lean{VOA.}\allowbreak\lean{InverseDS.}\allowbreak\lean{VirasoroVerma.}\allowbreak\lean{minimal_}\allowbreak\lean{three_}\allowbreak\lean{graded_}\allowbreak\lean{rigidity} (+ \lean{left}) & $M(p,3)$, $M(3,q)$ & labels; unconditional & §\ref{l07:sec:worked}\\
\lean{…minimal_}\allowbreak\lean{rigidity_}\allowbreak\lean{of_}\allowbreak\lean{infinite_}\allowbreak\lean{kac_}\allowbreak\lean{roots} & infinitely many Verma Gram roots $\Rightarrow$ rigid & \lean{hroots} (supplied by screening) & §\ref{l07:sec:all}\\
\lean{VirModule.}\allowbreak\lean{minimal_}\allowbreak\lean{graded_}\allowbreak\lean{rigidity} (headline) & every graded first-order deformation of every minimal model $L(c_{p,q},0)$ is trivial & $p,q\ge2$ coprime; \textbf{unconditional}, no \lean{pbw}; non-vacuity machine-checked ($V(c,0)$ not graded rigid) & Thm.~\ref{l07:thm:main}\\
\bottomrule
\end{longtable}
\endgroup

%======================================================================
\section{Notes and references}\label{l07:sec:notes}
%======================================================================

\textbf{Deformation theory.} The cohomological theory of deformations of associative algebras is due to
M.~Gerstenhaber, ``On the deformation of rings and algebras'', \emph{Annals of Mathematics} 79 (1964), 59--103;
the Lie-algebra analogue to A.~Nijenhuis and R.~W.~Richardson, ``Cohomology and deformations in graded Lie
algebras'', \emph{Bull.\ Amer.\ Math.\ Soc.} 72 (1966), 1--29. For vertex algebras, a cohomology theory in which
the second cohomology classifies first-order deformations of grading-restricted vertex algebras was developed by
Y.-Z.~Huang, ``A cohomology theory of grading-restricted vertex algebras'', \emph{Comm.\ Math.\ Phys.} 327 (2014),
279--307, and ``First and second cohomologies of grading-restricted vertex algebras'', \emph{Comm.\ Math.\ Phys.}
327 (2014), 261--278. The library constructs no cohomology theory and no comparison with Huang's is formalised.

\textbf{The target statements.} The affine counterpart, rigidity of $L_{-4/3}(\mathfrak{sl}_2)$, is the subject of the
paper of A.~Linshaw and F.~Qi, arXiv:2601.12017v2 (2026), whose Theorem~3.1 setting (regular deformations, pole
coefficient) is formalised in \leanpath{VOA/\allowbreak Core/\allowbreak Deformation.lean}; see Chapter~\ref{ch:L09}. The
Virasoro theorem of this chapter was formulated within the project; the project's working note
\texttt{Notes/Virasoro\_Uniform\_Rigidity\_Argument.tex} gives a short paper proof using the vacuum Kac
determinant of M.~Gorelik and V.~Kac, ``On simplicity of vacuum modules'', \emph{Advances in Mathematics} 211
(2007), 621--677, arXiv:math-ph/0606002, Theorem~4.2.1; the Lean proof does not depend on that citation.

\textbf{Virasoro representation theory.} The Kac determinant: V.~G.~Kac, ``Contravariant form for
infinite-dimensional Lie algebras and superalgebras'', in \emph{Group Theoretical Methods in Physics}, Lecture
Notes in Physics 94, Springer (1979), 441--445; B.~L.~Feigin and D.~B.~Fuchs, ``Verma modules over the Virasoro
algebra'', in \emph{Topology} (Leningrad 1982), Lecture Notes in Mathematics 1060, Springer (1984), 230--245. The
Benoit--Saint-Aubin null vectors: L.~Benoit and Y.~Saint-Aubin, ``Degenerate conformal field theories and explicit
expressions for some null vectors'', \emph{Phys.\ Lett.\ B} 215 (1988), 517--522. Fusion constructions of singular
vectors: e.g.\ A.~Kent, ``Singular vectors of the Virasoro algebra'', \emph{Phys.\ Lett.\ B} 273 (1991), 56--62. Free-field
(screening) realisations: V.~S.~Dotsenko and V.~A.~Fateev, ``Conformal algebra and multipoint correlation functions
in 2D statistical models'', \emph{Nucl.\ Phys.\ B} 240 (1984), 312--348.

\textbf{Vertex operator algebras of Virasoro type.} I.~B.~Frenkel and Y.~Zhu, ``Vertex operator algebras associated
to representations of affine and Virasoro algebras'', \emph{Duke Math.\ J.} 66 (1992), 123--168, construct $V(c,0)$
and $L(c,0)$ as VOAs; W.~Wang, ``Rationality of Virasoro vertex operator algebras'', \emph{Int.\ Math.\ Res.\ Notices}
1993, no.~7, 197--211, proves that $L(c_{p,q},0)$ is rational with irreducible modules given by the Kac table. Neither
rationality nor the module classification is used or proved in this chapter.

\textbf{Physics.} The minimal models: A.~A.~Belavin, A.~M.~Polyakov and A.~B.~Zamolodchikov, ``Infinite conformal
symmetry in two-dimensional quantum field theory'', \emph{Nucl.\ Phys.\ B} 241 (1984), 333--380; the unitary series:
D.~Friedan, Z.~Qiu and S.~Shenker, ``Conformal invariance, unitarity, and critical exponents in two dimensions'',
\emph{Phys.\ Rev.\ Lett.} 52 (1984), 1575--1578. For conformal perturbation theory and marginal deformations see
P.~Di~Francesco, P.~Mathieu and D.~S\'en\'echal, \emph{Conformal Field Theory}, Springer (1997), and for the
Lee--Yang model J.~L.~Cardy, ``Conformal invariance and the Yang--Lee edge singularity in two dimensions'',
\emph{Phys.\ Rev.\ Lett.} 54 (1985), 1354--1356. Standard VOA background: T.~Gannon, \emph{Moonshine beyond the Monster},
Cambridge University Press (2006), §§3.1, 5.1; J.~Lepowsky and H.~Li, \emph{Introduction to Vertex Operator Algebras and
Their Representations}, Birkh\"auser (2004).

\textbf{Project records.} The dead ends of Warnings~\ref{l07:warn:isrigid}--\ref{l07:warn:tautology} are also
recorded in the project's internal records; the book's own catalogue of dead ends is Section~\ref{ch15:sec:deadends}.

\endgroup
